\documentclass{article}
\PassOptionsToPackage{table}{xcolor}

\usepackage{iclr2027_conference,times}
\usepackage{enumitem}
\usepackage{graphicx}
\usepackage{amsmath,amsfonts,amssymb}
\usepackage{microtype}
\usepackage{tikz}
\usetikzlibrary{arrows.meta,backgrounds,positioning,shapes.geometric,calc,decorations.pathmorphing,fadings}
\usepackage{hyperref}
\usepackage{url}
\usepackage{booktabs}
\usepackage[compatibility=false]{caption}
\usepackage{subcaption}
\usepackage{flafter}
\usepackage{wrapfig}
\usepackage{xcolor}
% Four-band Exp 1 heatmap. Pastel fills keep black text legible; printed values
% remain the primary encoding.
\definecolor{heatred}{HTML}{F4CCCC}
\definecolor{heatorange}{HTML}{FCE5CD}
\definecolor{heatyellow}{HTML}{FFF2CC}
\definecolor{heatgreen}{HTML}{D9EAD3}

\hypersetup{
  colorlinks=true,
  linkcolor=[rgb]{0.10,0.25,0.55},
  citecolor=[rgb]{0.10,0.25,0.55},
  urlcolor=[rgb]{0.10,0.25,0.55}
}

\graphicspath{{figures/}{icons/}}
\newcommand{\micon}[1]{%
  \IfFileExists{icons/#1}{%
    \raisebox{-0.25\height}{\includegraphics[height=1em]{icons/#1}}%
  }{}%
}

\title{Can language models forecast inductively?}
% Keep the submission source itself anonymous: supplementary source archives are
% covered by ICLR's double-blind policy, even though the style hides this field.
\author{Anonymous authors \\
Paper under double-blind review}

% \iclrfinalcopy % Uncomment only for the camera-ready version.


% The ICLR style sets \flushbottom, so LaTeX stretches page glue to justify the
% bottom edge. Left alone, most of that slack lands in \textfloatsep and opens a
% large gap between a top float and the text under it. Removing the stretch
% component (keeping the shrink) pins the float-to-text distance and pushes the
% slack into the section skips instead, which absorb it far less visibly.
\setlength{\textfloatsep}{16pt plus 0pt minus 4pt}
\setlength{\intextsep}{12pt plus 0pt minus 3pt}
\setlength{\floatsep}{12pt plus 0pt minus 3pt}

\begin{document}
\maketitle

% Share the current replication estimates with the introduction.
\input{tables/exp3_rebuilt_data}
%============================================================================
% BEGIN frontmatter.tex
%============================================================================
\begin{abstract}
Judgmental forecasting requires inferring predictive relationships from sparse
observations and contextual knowledge. We ask whether language models recover
and use these relationships. Across prediction markets and synthetic worlds,
models selectively update forecasts, infer relationships from context, and
revise them as direct evidence accumulates. Larger tested models infer arbitrary
conventions; intervening on one model's internal representation changes its
forecast. Training improves forecasts across domains, but an eight-seed
replication weakens structural-transfer claims. A supervised pilot learns
response-form selection in distribution where the tested RL recipe fails;
its boundary is arbitrary labels, not a new domain. Accuracy gains on
compositions of familiar mechanisms exceed a fixed prior, but a gain intervention
does not establish transferred evidence use. Historical-market training lowers
mean error on unseen event families. These findings support context-sensitive
inductive forecasting while distinguishing in-domain learnability, transfer,
and episode-specific evidence use.\footnotemark
\end{abstract}
\footnotetext{%
  \micon{huggingface.png}\;%
  \href{https://huggingface.co/datasets/anonymous-review/forecasting-study-data}%
  {\textbf{Study dataset and artifacts}}\qquad
  \micon{gh_logo.png}\;%
  \href{https://github.com/anonymous-review/forecasting-study-code}%
  {\textbf{Code repository}}\quad
  \textit{Links anonymized for review.}}

\section{Introduction}
\label{sec:introduction}

\setlength{\columnsep}{6pt}
\begin{wrapfigure}{r}{0.44\textwidth}
\vspace{-1.2em}
\centering
\includegraphics[width=\linewidth,trim=0 72pt 0 0,clip]{teaser_revised.pdf}
\vspace{-0.45cm}
\caption{\textbf{Motivation.} Judgmental forecasts combine observations with
contextual knowledge about the process that generated them. Frontier models,
asked to forecast this question, average $P(H)\approx0.68$ (see
App.~\ref{app:carnival-coin-probe}).}
\vspace{0pt}
\label{fig:world-knowledge-evidence}
\end{wrapfigure}

Judgmental forecasting is a demanding form of reasoning central to decisions
about political, social, and economic futures. Unlike settings with abundant
historical data for statistical extrapolation, judgmental forecasting requires
integrating sparse, noisy evidence with contextual knowledge to form calibrated
beliefs about future outcomes
\citep{tetlock2016superforecasting,mellers2014psychological,lawrence2006judgmental}.
This creates an inductive challenge: forecasters must
infer from limited observations the relationships that govern how evidence
impacts outcomes, and use those relationships to anticipate how events may
unfold
\citep{tenenbaum2006theory,griffiths2010probabilistic,tenenbaum2011mind,lake2017building}.
Figure~\ref{fig:world-knowledge-evidence} illustrates this identification problem
in a low-stakes setting. A researcher observes a carnival coin-flipping machine
and, after seeing two heads, must estimate the probability that the next flip
will be heads. A judgmental forecast may incorporate contextual knowledge (e.g.,
whether the game is rigged) to infer the process generating the observed
outcomes.

Language models (LMs) now approach crowd accuracy in judgmental forecasting
\citep{halawi2024approaching,schoenegger2024wisdom}, but accurate predictions need
not imply recovery of the underlying task structure. Models recover latent
structure in some settings \citep{li2023emergent}, yet can predict accurately
using relationships that fail to generalize in others \citep{vafa2025foundation}.
This distinction matters when decision-makers ask how and whether a forecast
should change under alternative evidence or events. Answering such questions
requires predictive relationships that generalize across changes in context.

\noindent
\textbf{Present work.}
We study whether LMs exhibit inductive reasoning in forecasting tasks through a
series of progressively more stringent behavioral tests spanning both real-world
prediction markets and controlled synthetic environments. First, we evaluate
whether LMs respond selectively to new evidence, distinguishing
outcome-relevant information from topically related but orthogonal information
(\hyperref[sec:exp1]{Exp.~1}). Second, we test whether LMs infer predictive
relationships from described context before statistical evidence accumulates. We
construct synthetic forecasting problems in which the numerical observations are
identical across conditions and only the described setting varies, isolating the
contribution of background knowledge to the LM's forecast
(\hyperref[sec:exp2]{Exp.~2}). Finally, we test whether relationships learned
during fine-tuning transfer to new settings. We fine-tune LMs with reinforcement
learning from verifiable rewards (RLVR) on both synthetic tasks and historical
prediction market data, diagnose component failures with a supervised pilot,
and evaluate on held-out structures and domains
(\hyperref[sec:exp3]{Exp.~3}, \hyperref[sec:exp4]{Exp.~4}).

\noindent\textbf{Results.}
All nine systems move $1.9$--$7.5\times$ farther for outcome-relevant than topical
evidence, conditional on updating (\hyperref[sec:exp1]{Exp.~1}). In Coin City,
forecasts track context-selected relationships ($r=.52$--$.70$) and revise them
as direct evidence arrives; larger models infer episode-local conventions, with
a causal intervention linking the induced representation to Qwen3-14B's forecast
(\hyperref[sec:exp2]{Exp.~2}). Training transfers across domains, but its
advantage over an answer-matched shuffled control under joint domain and structure
shift is uncertain at eight seeds (Qwen3-8B:
$\eEightJointEst{}$, 95\% CI $\eEightJointCi{}$). A supervised pilot learns response-form selection in distribution and, trained
natively, fails under labels but not domain. A separate composition benchmark improves
accuracy over a fixed prior and an answer-matched shuffled control, without
clear evidence of transferred gain tracking (\hyperref[sec:exp3]{Exp.~3}). On 318 held-out markets, training lowers mean
error for every tested base and several checkpoints reach crowd performance,
although the Qwen3-32B improvement is uncertain
(\hyperref[sec:exp4]{Exp.~4}).
%============================================================================
% END frontmatter.tex
%============================================================================
%============================================================================
% BEGIN inductive_forecasting.tex
%============================================================================
\section{What does it mean to forecast inductively?}
\label{sec:inductive-forecasting}

Forecasting is a useful testbed for inductive reasoning because it makes
otherwise hidden choices observable. A forecaster must decide which evidence
matters, which population provides the relevant comparison, how strongly that
evidence should affect an outcome, and when direct observations should override
an existing belief. Each of these choices implies a different numerical
forecast. An accurate prediction alone does not tell us whether the forecaster
made those choices for the right reasons.

This distinction matters because a forecast is rarely the last question. We often
want to know what evidence should change the prediction, whether the same
evidence should matter differently in another population, and what new
observation should trigger revision. A model that reaches the right number
through a familiar base rate, a topical cue, or some other surface regularity may
still fail under exactly these changes. Average accuracy can therefore conceal an
important weakness: two models can give the same forecast while relying on very
different relationships between evidence and outcomes.

We use \emph{inductive forecasting} to describe the stronger behavior we care
about here. An inductive forecaster should not only identify whether evidence is
relevant, but also infer the relationship that determines how that evidence
should affect the outcome. When context changes which comparison is appropriate,
the forecast should change with it. When direct evidence becomes available, that
evidence should be able to revise the earlier contextual inference. And when the
same underlying relationship appears in an unfamiliar setting, the forecaster
should be able to use it without depending on the original surface form.

\begin{figure}[!t]
\centering
\includegraphics[width=\linewidth]{exp1_protocol_four_panels.pdf}
\caption{\textbf{When should a forecast change?} An agent forecasts an
unresolved question, then receives one fictional update with search off. Some
updates support the outcome, some oppose it, and some concern the same topic but
do not bear on the outcome. The labels remain hidden; we ask whether the model
moves when the news matters and whether it moves the right way.}
\label{fig:exp1-pipeline}
\end{figure}

These requirements extend beyond forecasting. A persuasion model should respond
not only to a message, but to who delivers it. A coordination model should
distinguish shared incentives from what participants know about one another. A
policy model should distinguish similar language from the institutional pathway
through which an intervention takes effect. In each case, the conditional
relationship is often the object of interest: who trusts whom, which population
responds to which signal, or which mechanism connects an action to an outcome.
Matching an average outcome while missing these dependencies can produce a
convincing static prediction and a misleading response when the setting changes
\citep{park2023generative,aher2023simulate,argyle2023out,bisbee2024synthetic}.

This suggests a stronger way to evaluate forecasting systems. Rather than asking
only whether a model reaches the correct number, we can ask whether its forecast
changes when the relevant evidence, comparison, or relationship changes. We
therefore proceed through increasingly demanding tests: first whether models
distinguish relevant evidence from topical information, then whether they infer
how strongly evidence should matter from context, and finally whether those
relationships can be revised and carried into new settings.
%============================================================================
% END inductive_forecasting.tex
%============================================================================
%============================================================================
% BEGIN evidence_ladder.tex
%============================================================================
\section{Do forecasts change for the right reason?}
\label{sec:evidence}

A key forecasting skill is to appropriately incorporate new evidence as it
arrives; for example, classic work on judgmental forecasting identifies repeated
belief revision as a hallmark of strong human forecasters
\citep{tetlock2005expert}. We argue that effective updating requires a forecaster
to determine: (1) which evidence is relevant in predicting an outcome and (2) the
extent to which the new evidence should affect the forecast. We assess these
capacities in two experiments. First, we pair real-world prediction-market
questions with controlled synthetic evidence, allowing us to test whether models
update selectively and in the direction the evidence warrants. Second, we use a
controlled synthetic environment in which the same evidence has different
predictive effects across contexts, allowing us to test whether models can infer
the relationship that determines how strongly the evidence should affect their
forecast.

\subsection{Exp.~1: Can models directionally update?}
\label{sec:exp1}

\noindent\textbf{Design.}
We ask nine systems to forecast 100 unresolved binary Polymarket questions
with web search, then revise each forecast after one fictional news update with
search disabled (Fig.~\ref{fig:exp1-pipeline}). Updates support YES, support NO,
or concern the same topic without bearing on the outcome; these labels are hidden
from the model. This yields 30,094 valid updates
(App.~\ref{app:exp1}).

\noindent\textbf{Finding: models track the intended packet direction.}
Every system leaves its forecast unchanged more often for topical controls.
Conditional on moving, models move $1.9$--$7.5\times$ farther for directional
than topical updates. For revisions of at least three percentage points,
directional consistency ranges from $.887$ to $.990$, versus $.688$ for a
word-only classifier evaluated on held-out markets. Human validation and
uncertainty estimates are in App.~\ref{app:exp1-threshold-robustness}.

\begin{figure}[!b]
\centering
\begin{subfigure}[c]{0.525\textwidth}
\centering
\includegraphics[width=.92\linewidth]{exp1_direction_selective_updating.pdf}
\caption{}
\label{fig:exp1-bars}
\end{subfigure}\hspace{0.01\textwidth}
\begin{subfigure}[c]{0.455\textwidth}
\centering
\footnotesize
\setlength{\tabcolsep}{3pt}
\begin{tabular*}{\linewidth}{@{}l@{\extracolsep{\fill}}rccc@{}}
\toprule
\textbf{Model} & \textbf{Size} & \textbf{Abst.} & \textbf{Sens.}$\uparrow$ & \textbf{EHC$_{\geq3}$}$\uparrow$ \\
\midrule
\multicolumn{5}{@{}l}{\emph{Hosted API (parameters undisclosed)}}\\
\micon{claude.png}~Opus~4.8 & --- & 71.4\% & \cellcolor{heatgreen}$7.0\times$ & \cellcolor{heatgreen}.990 \\
\micon{openai.png}~GPT-5.4 & --- & 8.4\% & \cellcolor{heatgreen}$7.5\times$ & \cellcolor{heatgreen}.987 \\
\micon{deepseek.png}~V4-Pro & --- & 44.4\% & \cellcolor{heatyellow}$4.9\times$ & \cellcolor{heatyellow}.977 \\
\midrule
\multicolumn{5}{@{}l}{\emph{Open-weight, descending parameters} $\downarrow$}\\
\micon{qwen.pdf}~2.5-72B & 72B & 24.2\% & \cellcolor{heatorange}$3.2\times$ & \cellcolor{heatyellow}.963 \\
\micon{llamma.pdf}~3.3-70B & 70B & 10.7\% & \cellcolor{heatyellow}$5.7\times$ & \cellcolor{heatyellow}.965 \\
\micon{llamma.pdf}~3.1-70B & 70B & 2.5\% & \cellcolor{heatyellow}$4.5\times$ & \cellcolor{heatyellow}.961 \\
\micon{qwen.pdf}~2.5-32B & 32B & 11.1\% & \cellcolor{heatorange}$2.8\times$ & \cellcolor{heatyellow}.964 \\
\micon{qwen.pdf}~2.5-14B & 14B & 26.1\% & \cellcolor{heatred}$2.2\times$ & \cellcolor{heatorange}.943 \\
\micon{llamma.pdf}~3.1-8B & 8B & 1.1\% & \cellcolor{heatred}$1.9\times$ & \cellcolor{heatred}.887 \\
\bottomrule
\end{tabular*}
\caption{}
\label{fig:exp1-table}
\label{tab:consistency}
\end{subfigure}
\caption{\textbf{Models move much more for evidence than for topical
chatter, and usually move in the implied direction.} \textbf{(a)} Average
revision after supporting, contrary, and topical updates. \textbf{(b)} The
table summarizes how often models leave topical updates alone, how much farther
they move for directional evidence conditional on moving (Sens.), and how often
a move of at least three percentage points has the right sign (EHC$_{\geq3}$). Intervals and robustness checks are in
App.~\ref{app:exp1-threshold-robustness}.}
\label{fig:exp1-updating}
\end{figure}

\noindent\textbf{Takeaway.}
Lexical cues identify original-packet relevance ($.989$ accuracy). On 20
unvalidated context-reversal families, GPT-5.6 reverses both directions in
85\% of pairs, with 0.5-point broken-link movement; the original constrained
local deployments reach 25--35\% and 9.75--13.25 points in the matched
single-repeat comparison. A matched Qwen3 follow-up reaches 60\% paired
reversal with thinking enabled (App.~\ref{app:exp1-context-reversal}). This
exploratory extension uses different inference settings and requires independent
material validation. Experiment~2
tests contextual relationship selection under a controlled generative process.

\subsection{Exp.~2: can models infer how strongly evidence should matter?}
\label{sec:exp2}

\begin{figure}[!t]
\centering
\includegraphics[
  width=.90\linewidth,
  trim=11pt 14.5pt 3.5pt 7pt,
  clip
]{exp2_coin_city_prompt_arms.pdf}
\caption{\textbf{Coin City turns contextual reasoning into a paired
comparison.} Every arm shares the query and numerical cases. The context arm
adds a description connecting City C to one displayed media environment; the
opposite-context arm changes only that sentence. Target cases then accumulate,
allowing local evidence to confirm or correct the contextual choice.}
\label{fig:coin-city-prompt}
\end{figure}

\noindent\textbf{Design.}
Experiment~1 tests which evidence warrants an update. Forecasting also requires
inferring how strongly that evidence should matter. In the opening coin-flip
example (Fig.~\ref{fig:world-knowledge-evidence}), the same observation warrants
different predictions depending on the process believed to have generated it.

We systematize this ambiguity in \emph{Coin City}: the same campaign news moves
polls strongly in one reference city and weakly in another. Context identifies the
better analogue for target City C, while zero to four City C cases provide direct
evidence that can confirm or override it. Each city's polling response follows
\begin{equation}
  \Delta p_{ij}=g_j x_{ij}+\epsilon_{ij},
  \qquad \epsilon_{ij}\sim\mathcal{N}(0,5^2),
  \label{eq:coin-city-dgp}
\end{equation}
where $x_{ij}$ is signed campaign news and $g_j$ is the city's stable sensitivity,
which is never stated: the model must infer it from examples.

We compare target cases alone, reference cases without target context, correct
context, and opposite context, holding numerical evidence fixed across the paired
arms (Fig.~\ref{fig:coin-city-prompt}). A fifth arm replaces familiar descriptions
with episode-varying labels whose meaning must be inferred from the references.
Forecast error tests whether context helps; the forecast-implied response to news
tests which reference the model selects (App.~\ref{app:coin-city-results}).

\noindent\textbf{Finding: the forecast follows the contextual clue.}
Figure~\ref{fig:coin-city-results} shows that, when City C provides few direct
observations, context guides which reference city's response is used to forecast.
The correct context lowers forecast error, while the wrong analogue increases
error. Moreover, the effect implies the model is quantifying the importance of
the contextual clue: after controlling for the underlying response regime, the
forecast-implied slope correlates $r=.51$--$.86$ with the fitted slope of the
cue-selected reference city, while its correlation with the other reference
remains near zero. As City C observations accumulate, context matters less, as
forecasts increasingly rely on direct local evidence.

\begin{figure}[!t]
\centering
\includegraphics[width=\linewidth]{exp2_coin_city_results.pdf}
\vspace{3pt}
\begin{minipage}{\linewidth}
\centering
\footnotesize
\setlength{\tabcolsep}{4pt}
\resizebox{\linewidth}{!}{%
\begin{tabular}{l r rrrr rr}
\toprule
& & \multicolumn{4}{c}{Forecast MAE [95\% CI]} & \multicolumn{2}{c}{$\rho(\hat g,g_C)$} \\
\cmidrule(lr){3-6}\cmidrule(lr){7-8}
Model & Paired $n$ & Target only & No C context & $+$ C context &
$+$ wrong C context & No context & $+$ context \\
\midrule
\micon{claude.png}~Opus~4.8 & 250 & 2.32 [2.20,\,2.46] & 2.65 [2.41,\,2.89] & \textbf{1.82} [1.64,\,2.02] & 4.20 [3.87,\,4.54] & $-.06$ & \textbf{.70} \\
\micon{claude.png}~Opus~5 & 250 & 2.36 [2.12,\,2.59] & 2.62 [2.38,\,2.84] & \textbf{1.76} [1.57,\,1.99] & 4.25 [3.92,\,4.57] & $-.01$ & \textbf{.62} \\
\micon{openai.png}~GPT-5.6 & 250 & 3.60 [3.28,\,3.94] & 2.52 [2.33,\,2.75] & \textbf{1.60} [1.42,\,1.77] & 4.23 [3.90,\,4.55] & $-.04$ & \textbf{.68} \\
\micon{deepseek.png}~V4-Pro & 250 & \textbf{2.45} [2.27,\,2.61] & 2.84 [2.55,\,3.15] & 2.47 [2.19,\,2.77] & 4.36 [4.01,\,4.74] & $-.03$ & \textbf{.52} \\
\micon{kimi.png}~K3 & 250 & 3.35 [3.04,\,3.67] & 2.55 [2.32,\,2.78] & \textbf{1.89} [1.66,\,2.13] & 4.31 [3.98,\,4.62] & $-.02$ & \textbf{.59} \\
\micon{gemini.png}~3.6 Flash & 250 & 2.53 [2.35,\,2.73] & 2.55 [2.31,\,2.80] & \textbf{2.17} [1.89,\,2.48] & 4.29 [3.94,\,4.66] & $.08$ & \textbf{.56} \\
\bottomrule
\end{tabular}
}
\end{minipage}
\caption{\textbf{Context improves forecasts before target cases and helps assign
the response pattern across six systems.} Giving the correct analogue lowers
error before City C has a history; naming the wrong analogue raises it. As direct
City C cases accumulate, forecasts rely less on either contextual cue. Full
results by model are in App.~\ref{app:coin-city-results}.}
\label{fig:coin-city-results}
\end{figure}

\subsubsection{Infer a convention without stable semantics}
\label{sec:novel-mapping}

However, a different explanation for the relational mapping above is that the
model may already associate familiar cues, like ``national news,'' with a
stronger response and ``local news'' with a weaker one, rather than actually
inferring the mapping from the reference cases. We remove this shortcut by
replacing these cues with the arbitrary labels KIV and ZOR. Their meanings are
randomized across episodes, so neither label has any stable association with
strong or weak responsiveness. Twelve of fifteen tested systems recover the episode-local mapping before any
target cases, including all tested local checkpoints above 8B
(App.~\ref{app:coin-city-symbol-context}), although recovery does not always
improve forecast error. When we further challenge this
inference by making the two response regimes harder to distinguish, recovery
weakens, but sufficiently large models are still able to recover the
relationship, with the threshold varying by model family
(App.~\ref{app:coin-city-robustness}).

\begin{figure}[!t]
\centering
\begin{minipage}[c]{0.46\linewidth}
\includegraphics[width=\linewidth]{exp2_symbol_decoding.pdf}
\end{minipage}\hfill
\begin{minipage}[c]{0.50\linewidth}
\caption{\textbf{Larger models infer an episode-local convention, then yield
 to direct evidence.} On 88 new episodes, four larger Qwen and Llama checkpoints
(14--72B) encode the changing KIV/ZOR mapping; three smaller checkpoints
(4--8B) do not reliably encode it. The decoded state follows context before
City C has a history and follows the target relationship as local cases accumulate.}
\label{fig:coin-city-symbol-decoding}
\end{minipage}
\end{figure}

This pattern matches what we would expect if models are inferring a latent
mapping: doing so becomes harder as the underlying regimes become less
distinguishable. \emph{But when a model does recover that mapping, what has it
actually inferred?} We ask whether the episode-local KIV/ZOR relationship can be
decoded from the model's hidden states. In larger open models, the decoded
representation follows the context-selected relationship before City C has its
own history, and then shifts toward City C's observed behavior as direct evidence
accumulates (Fig.~\ref{fig:coin-city-symbol-decoding}).

But decoding still leaves open whether this representation actually matters for
the forecast. Figure~\ref{fig:coin-city-causal-patch} tests this directly by
swapping the label-state representation between otherwise matched Qwen3-14B
prompts. Before City C provides any direct observations, the intervention changes
both the analogue the model names and its numerical forecast. After four City C
cases, the same swap has no reliable effect. Across episodes, the intervention
produces a regime-aligned symmetric forecast shift of \(+.333\) before target
cases and \(-.035\) after four; the latter falls within the sign-flip null range. For Qwen3-14B, this shows that
the induced mapping helps determine the forecast when target evidence is sparse,
but gives way once direct observations reveal City C's own response relationship
(App.~\ref{app:coin-city-causal-patch}).

\begin{figure}[!t]
\vspace{-0.4em}
\centering
\includegraphics[width=\linewidth]{exp2_causal_patch.pdf}
\vspace{-0.7em}
\caption{\textbf{Changing the induced convention changes the forecast when
evidence is sparse.} In the illustrated pair, swapping the internal
representation moves the forecast from 51.3 to 54.0 and changes which reference
city the model names; reversing the swap reverses the movement. As more City C
behavior is revealed, the same label-state swap has no reliable effect. Full controls in
Table~\ref{tab:coin-city-causal-patch}.}
\label{fig:coin-city-causal-patch}
\vspace{-1.0em}
\end{figure}

\noindent\textbf{Takeaway.}
Models use context to select a predictive relationship and revise that selection
from local evidence. Arbitrary labels remove a familiar-semantic shortcut;
decoding reveals an episode-local representation, and intervention links it to
Qwen3-14B's forecast when target evidence is absent. Experiment~3 asks which parts
of this behavior training carries into new domains and response structures.
%============================================================================
% END evidence_ladder.tex
%============================================================================
%============================================================================
% BEGIN transfer_and_grounding.tex
%============================================================================
\section{Can relationship use survive a change of world?}
\label{sec:transfer}

Finetuning may teach evidence use or merely typical answers. We test whether
training transfers relationships to new worlds and improves unseen-event forecasts.

\subsection{Exp.~3: what training transfers, and what it does not}
\label{sec:exp3}

\input{tables/exp3_rebuilt_data}
\newcommand{\transfergridlines}{15}
%============================================================================
% BEGIN fig_exp3_transfer_grid.tex
%============================================================================
% Compact inline orientation graphic for the Experiment 3 transfer design: a 2x2
% grid crossing generative mechanism (rows) with surface domain (columns).
%
% Row cards carry the mechanism detail: a two-node vs three-node graph with the
% mediator drawn hollow (it is never named to the model) and a persistence loop,
% plus the six-period impulse response on labelled axes. Bar heights are the
% normalised responses to a unit shock under worlds.transition at the midpoint of
% each sampling band (direct: gain .57; mediated: gain .57, rho .61, phi .30,
% lambda .81), i.e. A = (1, 0, 0, 0, 0, 0) and B = (.81, .74, .52, .34, .21, .13).
%
% Geometry note: cards have fixed dimensions and every label is a single line
% with no text width, so nothing can rewrap into a neighbour.
% Callers may set \transfergridlines to match the paragraph they anchor this
% to; wrapfig cannot wrap a section heading, so the count must not overrun it.
\providecommand{\transfergridlines}{13}
\begin{wrapfigure}[\transfergridlines]{r}{0.48\textwidth}
\vspace{-0.8em}
\centering
\definecolor{miniCityFill}{HTML}{E8F2FB}
\definecolor{miniCityEdge}{HTML}{4F7EA8}
\definecolor{miniHarborEdge}{HTML}{3F8E82}
\definecolor{miniTestFill}{HTML}{E7F4EA}
\definecolor{miniTestEdge}{HTML}{6FA97F}
\definecolor{miniTestChip}{HTML}{5E9470}
% Coin discs reuse Figure 1's yellow pair exactly -- the fill and stroke of its
% NAIVE box -- rather than the saturated gold used before. Figure 1's other warm
% pair (#FFE6CC on #D79B00, its FREQUENTIST box) is the orange one and is not
% used here. The glyph inside each disc still carries its domain colour.
\definecolor{miniCoinFill}{HTML}{FFF2CC}
\definecolor{miniCoinRim}{HTML}{D6B656}
\definecolor{miniMute}{HTML}{8A93A3}
\definecolor{miniInk}{HTML}{3F4550}
\definecolor{miniBar}{HTML}{7C93AD}
% Coin City: a skyline stamped into a coin.
\newcommand{\coincity}[2]{%
  \begin{scope}[shift={(#1,#2)},scale=0.70]
    \fill[miniCoinFill] (0,0) circle (.37);
    \draw[miniCoinRim,line width=.75pt] (0,0) circle (.37);
    \fill[miniCityEdge] (-.23,-.18) rectangle (-.10,.04);
    \fill[miniCityEdge] (-.07,-.18) rectangle (.07,.16);
    \fill[miniCityEdge] (.10,-.18) rectangle (.24,.09);
    \draw[miniCityEdge,line width=.65pt] (-.27,-.19)--(.27,-.19);
  \end{scope}}
% Coin Harbor: a small vessel and waves stamped into a coin.
\newcommand{\coinharbor}[2]{%
  \begin{scope}[shift={(#1,#2)},scale=0.70]
    \fill[miniCoinFill] (0,0) circle (.37);
    \draw[miniCoinRim,line width=.75pt] (0,0) circle (.37);
    \fill[miniHarborEdge] (-.24,-.05)--(.24,-.05)--(.15,-.18)--(-.15,-.18)--cycle;
    \draw[miniHarborEdge,line width=.7pt]
      (-.01,-.05)--(-.01,.15)--(.15,.02)--(-.01,.02);
    \draw[miniHarborEdge,line width=.55pt]
      (-.26,-.24)..controls(-.18,-.19)and(-.10,-.29)..(-.02,-.24)
      ..controls(.06,-.19)and(.14,-.29)..(.26,-.24);
  \end{scope}}
% Impulse response on labelled axes: response (Delta y) against period (t).
\newcommand{\impulse}[3]{% #1 = axis corner x, #2 = baseline y, #3 = bar heights
  \draw[-{Stealth[length=2.6pt,width=2.2pt]},draw=miniInk!65,line width=.5pt]
    (#1,#2) -- (#1,#2+0.40);
  \draw[-{Stealth[length=2.6pt,width=2.2pt]},draw=miniInk!65,line width=.5pt]
    (#1,#2) -- (#1+0.72,#2);
  \node[anchor=south,font=\sffamily\scriptsize,text=miniInk,inner sep=1pt]
    at (#1,#2+0.40) {$\Delta y$};
  \node[anchor=west,font=\sffamily\scriptsize,text=miniInk,inner sep=1.5pt]
    at (#1+0.72,#2) {$t$};
  \foreach \i/\h in {#3}{%
    \fill[miniBar] (#1+0.055+\i*0.105,#2) rectangle (#1+0.055+\i*0.105+0.070,#2+\h);}}
\resizebox{\linewidth}{!}{%
\begin{tikzpicture}[
  x=1cm, y=1cm,
  cell/.style={rounded corners=3pt, minimum width=1.95cm, minimum height=1.25cm,
    line width=.6pt, inner sep=0pt},
  traincell/.style={cell, draw=miniCityEdge, fill=miniCityFill},
  heldcell/.style={cell, draw=miniTestEdge, fill=miniTestFill,
    dash pattern=on 1.9pt off 1.5pt},
  rowcard/.style={rounded corners=3pt, minimum width=2.60cm,
    minimum height=1.25cm, inner sep=0pt, draw=miniMute!30, fill=black!3,
    line width=.5pt},
  dname/.style={anchor=south, font=\sffamily\bfseries\scriptsize},
  rname/.style={anchor=center, font=\sffamily\scriptsize, text=miniInk},
  note/.style={anchor=center, font=\sffamily\scriptsize, text=miniInk},
  gnode/.style={circle, draw=miniInk!75, fill=white, line width=.5pt,
    inner sep=0pt, minimum size=0.26cm, font=\sffamily\tiny, text=miniInk},
  hidden/.style={gnode, draw=miniMute, fill=miniMute!14,
    dash pattern=on 1.1pt off .9pt},
  garrow/.style={-{Stealth[length=2.4pt,width=2pt]}, draw=miniInk!70,
    line width=.45pt, shorten >=.4pt, shorten <=.4pt},
  chip/.style={anchor=center, rounded corners=1.6pt, draw=none, text=white,
    inner xsep=3pt, inner ysep=1.1pt, font=\sffamily\bfseries\tiny},
]
% ---- domain marks and names ----------------------------------------------
\coincity{0.975}{2.56}
\coinharbor{3.075}{2.56}
\node[dname, text=miniCityEdge]   at (0.975,1.80) {COIN CITY};
\node[dname, text=miniHarborEdge] at (3.075,1.80) {COIN HARBOR};

% ---- Structure A: direct, one period -------------------------------------
\node[rowcard] at (-1.50,1.00) {};
\node[rname] at (-1.50,1.45) {\textbf{Structure A} direct};
\node[gnode] (ax) at (-2.51,0.82) {$x$};
\node[gnode] (ay) at (-2.07,0.82) {$y$};
\draw[garrow] (ax) -- (ay);
\impulse{-1.22}{0.62}{0/0.360, 1/0.000, 2/0.000, 3/0.000, 4/0.000, 5/0.000}

% ---- Structure B: mediated, persistent -----------------------------------
\node[rowcard] at (-1.50,-0.42) {};
\node[rname] at (-1.50,0.03) {\textbf{Structure B} mediated};
\node[gnode]  (bx) at (-2.51,-0.60) {$x$};
\node[hidden] (bm) at (-2.07,-0.60) {$m$};
\node[gnode]  (by) at (-1.63,-0.60) {$y$};
\draw[garrow] (bx) -- (bm);
\draw[garrow] (bm) -- (by);
\draw[garrow] (bm) to[out=118,in=62,looseness=4.5] (bm);
\impulse{-1.22}{-0.80}{0/0.292, 1/0.265, 2/0.188, 3/0.123, 4/0.077, 5/0.048}

% ---- the four evaluation cells -------------------------------------------
\node[traincell] at (0.975,1.00) {};
\node[heldcell]  at (3.075,1.00) {};
\node[heldcell]  at (0.975,-0.42) {};
\node[heldcell]  at (3.075,-0.42) {};

\node[chip, fill=miniCityEdge] at (0.975,1.26) {TRAINED};
\node[chip, fill=miniTestChip] at (3.075,1.26) {TEST};
\node[chip, fill=miniTestChip] at (0.975,-0.16) {TEST};
\node[chip, fill=miniTestChip] at (3.075,-0.16) {TEST};

\node[note] at (0.975,0.80) {in-distribution};
\node[note] at (3.075,0.80) {domain shift};
\node[note] at (0.975,-0.62) {structure shift};
\node[note] at (3.075,-0.62) {joint shift};
\end{tikzpicture}%
}
\vspace{-0.3em}
\caption{\textbf{Transfer grid.} Training uses shading;
all four are evaluated. Columns vary domain; rows vary response mechanism.}
\label{fig:exp3-transfer-grid}
\vspace{0.1em}
\end{wrapfigure}
%============================================================================
% END fig_exp3_transfer_grid.tex
%============================================================================

\textbf{Why matching matters.}
The parent experiment compares byte-identical prompts, identical target-answer
multisets, and equal training budgets, changing only case--answer pairing.
Figure~\ref{fig:exp3-transfer-grid} crosses a domain change---campaign news and
candidate support become shipping bulletins and cargo flow---with a change
from direct to persistent responses (App.~\ref{app:exp3-complete-training-setup}).

\noindent\textbf{Finding: domain transfer is stronger than structural transfer.}
Training generalizes across the surface change. Across both arms and
\eEightG{} seeds, Qwen3-8B reduces domain-shift response error by
\eEightDomainBaseGain{} points \eEightDomainBaseGainCi{}, with every seed agreeing.
The matched arm reaches \eEightDomainMatched{} versus the base's
\eEightDomainBase{}.
Because unparsed draws take the full
range penalty, we report base contrasts twice: the joint-cell gain is
\eEightJointBaseGain{} \eEightJointBaseGainCi{} fail-closed and
\eEightJointBaseGainParsed{} \eEightJointBaseGainParsedCi{} on parsed draws.

%============================================================================
% BEGIN fig_exp3_results_grid.tex
%============================================================================
% Experiment 3 transfer results, laid out on the same 2x2 grid as the design
% figure (Figure~\ref{fig:exp3-transfer-grid}) so the reader reads results in the
% cells they were introduced in: mechanism down the rows, surface domain across
% the columns, the trained cell shaded blue and the three held-out cells green.
%
% Every number comes from tables/exp3_coin_qwen3_8b_stochastic_data.tex,
% regenerated by coin_city_structural/make_qwen3_8b_eight_seed_outputs.py. Bar
% lengths share one scale
% (\cciPeak) across all four cells so cells are comparable by eye. Do not put
% formatting logic here: emphasis on an interval that excludes zero is decided
% by the generator and arrives pre-set in \cci*delta.
\input{tables/exp3_coin_qwen3_8b_stochastic_data}
\begin{figure}[!ht]
\centering
\definecolor{gCityFill}{HTML}{E8F2FB}
\definecolor{gCityEdge}{HTML}{4F7EA8}
\definecolor{gHarborEdge}{HTML}{3F8E82}
\definecolor{gTestFill}{HTML}{E7F4EA}
\definecolor{gTestEdge}{HTML}{6FA97F}
\definecolor{gTestChip}{HTML}{5E9470}
\definecolor{gCoinFill}{HTML}{FFF2CC}
\definecolor{gCoinRim}{HTML}{D6B656}
\definecolor{gInk}{HTML}{3F4550}
\definecolor{gMute}{HTML}{8A93A3}
\definecolor{gBarBase}{HTML}{C3CBD4}
\definecolor{gBarLess}{HTML}{A9B4C0}
\definecolor{gBarPrior}{HTML}{7C93AD}
\definecolor{gBarMatched}{HTML}{4F7EA8}
\newcommand{\gcoincity}[2]{%
  \begin{scope}[shift={(#1,#2)},scale=0.70]
    \fill[gCoinFill] (0,0) circle (.37);
    \draw[gCoinRim,line width=.75pt] (0,0) circle (.37);
    \fill[gCityEdge] (-.23,-.18) rectangle (-.10,.04);
    \fill[gCityEdge] (-.07,-.18) rectangle (.07,.16);
    \fill[gCityEdge] (.10,-.18) rectangle (.24,.09);
    \draw[gCityEdge,line width=.65pt] (-.27,-.19)--(.27,-.19);
  \end{scope}}
\newcommand{\gcoinharbor}[2]{%
  \begin{scope}[shift={(#1,#2)},scale=0.70]
    \fill[gCoinFill] (0,0) circle (.37);
    \draw[gCoinRim,line width=.75pt] (0,0) circle (.37);
    \fill[gHarborEdge] (-.24,-.05)--(.24,-.05)--(.15,-.18)--(-.15,-.18)--cycle;
    \draw[gHarborEdge,line width=.7pt] (-.01,-.05)--(-.01,.15)--(.15,.02)--(-.01,.02);
    \draw[gHarborEdge,line width=.55pt]
      (-.26,-.24)..controls(-.18,-.19)and(-.10,-.29)..(-.02,-.24)
      ..controls(.06,-.19)and(.14,-.29)..(.26,-.24);
  \end{scope}}
% One results card. #1 cx, #2 cy, #3 cell name, #4 chip text, #5 chip fill,
% #6..#8 base/prior/episode-matched MAE; the contrast is placed separately.
\newcommand{\gcard}[8]{%
  \node[cellname] at (#1-2.66,#2+0.78) {#3};
  \node[chip, fill=#5] at (#1+2.66,#2+0.78) {#4};
  \foreach \arm/\val/\col/\dy in {%
      base/#6/gBarBase/0.34, prior/#7/gBarPrior/0.00,
      episode-matched/#8/gBarMatched/-0.34}{%
    \node[armlab] at (#1-2.66,#2+\dy) {\arm};
    \fill[\col, rounded corners=1pt]
      (#1-0.55,#2+\dy-0.075) rectangle (#1-0.55+\val*2.55/\cciPeak,#2+\dy+0.075);
    \node[armval] at (#1-0.43+\val*2.55/\cciPeak,#2+\dy) {\val};}}
\resizebox{.96\linewidth}{!}{%
\begin{tikzpicture}[
  x=1cm, y=1cm,
  card/.style={rounded corners=3pt, minimum width=5.80cm, minimum height=1.95cm,
    line width=.6pt, inner sep=0pt},
  traincard/.style={card, draw=gCityEdge, fill=gCityFill},
  testcard/.style={card, draw=gTestEdge, fill=gTestFill,
    dash pattern=on 1.9pt off 1.5pt},
  dname/.style={anchor=south, font=\sffamily\bfseries\small},
  rname/.style={anchor=east, align=right, font=\sffamily\small, text=gInk},
  cellname/.style={anchor=west, font=\sffamily\bfseries\footnotesize, text=gInk},
  chip/.style={anchor=east, rounded corners=1.6pt, draw=none, text=white,
    inner xsep=3pt, inner ysep=1.1pt, font=\sffamily\bfseries\scriptsize},
  armlab/.style={anchor=west, font=\sffamily\scriptsize, text=gInk},
  armval/.style={anchor=west, font=\sffamily\scriptsize, text=gInk},
  delta/.style={anchor=center, font=\sffamily\scriptsize, text=gInk},
  legend/.style={anchor=east, font=\sffamily\scriptsize},
]
% ---- domain columns ------------------------------------------------------
\gcoincity{2.90}{3.22}
\gcoinharbor{9.00}{3.22}
\node[dname, text=gCityEdge]   at (2.90,2.40) {COIN CITY};
\node[dname, text=gHarborEdge] at (9.00,2.40) {COIN HARBOR};

% ---- what the bars mean, and which way is good --------------------------
\node[legend, text=gInk]  at (-0.20,2.62) {bars: response MAE};
\node[legend, text=gMute] at (-0.20,2.28) {$\leftarrow$ lower is better};

% ---- mechanism rows ------------------------------------------------------
\node[rname] at (-0.20,1.25) {\textbf{Structure A}\\ direct};
\node[rname] at (-0.20,-0.85) {\textbf{Structure B}\\ mediated};

% ---- the four result cards ----------------------------------------------
\node[traincard] at (2.90,1.25) {};
\node[testcard]  at (9.00,1.25) {};
\node[testcard]  at (2.90,-0.85) {};
\node[testcard]  at (9.00,-0.85) {};

\gcard{2.90}{1.25}{in-distribution}{TRAINED}{gCityEdge}
  {\cciIDbase}{\cciIDprior}{\cciIDmatched}
\node[delta] at (2.90,0.42) {prior $-$ matched\ \ \cciIDdelta};

\gcard{9.00}{1.25}{domain shift}{TEST}{gTestChip}
  {\cciDomainbase}{\cciDomainprior}{\cciDomainmatched}
\node[delta] at (9.00,0.42) {prior $-$ matched\ \ \cciDomaindelta};

\gcard{2.90}{-0.85}{structure shift}{TEST}{gTestChip}
  {\cciStructurebase}{\cciStructureprior}{\cciStructurematched}
\node[delta] at (2.90,-1.68) {prior $-$ matched\ \ \cciStructuredelta};

\gcard{9.00}{-0.85}{joint shift}{TEST}{gTestChip}
  {\cciJointbase}{\cciJointprior}{\cciJointmatched}
\node[delta] at (9.00,-1.68) {prior $-$ matched\ \ \cciJointdelta};
\end{tikzpicture}%
}
\caption{\textbf{Qwen3-8B Coin City transfer at round 300 under five-draw
stochastic decoding, eight training seeds.} Bars show task-averaged response MAE
over four evidence depths; trained arms average eight seeds. Each cell prints the
paired seed-first prior-minus-matched contrast, bold when its 95\% interval
excludes zero. At three seeds three of the four contrasts excluded zero; at eight
only the in-distribution cell does. The structure-shift row is a control rather than
a transfer test: training contains no structural variation, so the trained policy
emits no delayed response anywhere (Section~\ref{sec:exp3}).}
\label{fig:exp3-transfer-results}
\vspace{10pt}
\end{figure}
%============================================================================
% END fig_exp3_results_grid.tex
%============================================================================

The finer contrast between episode-matched and shuffled-pair training does not
survive replication. Estimated at three seeds and extended to eight under a
pre-registered amendment, the Qwen3-8B joint-shift contrast falls from
\ePubEightJointEst{} \ePubEightJointCi{} to \eEightJointEst{} \eEightJointCi{};
all four of Qwen3-4B's cells excluded zero at three seeds and all span it at
eight. A wrapper crossover retains a similar three-seed mean ($.251$ versus $.255$;
App.~\ref{app:exp3-eight-seed}). One cell survives every test: Qwen3-8B in distribution,
\eEightIdEst{} \eEightIdCi{}, \eEightIdSign{} seeds, clearing the conservative
$t$ \eEightIdTci{}, a wild cluster bootstrap ($p=\eEightIdWild{}$) and an exact
sign test ($p=\eEightIdSignP{}$). It does not replicate on Qwen3-4B, and
Llama-3.1-8B's estimate has the opposite sign.

\noindent\textbf{Response-form selection is learnable, but transfer is limited.}
The original structure-shift row is a control: training contains only the direct
response, and every trained checkpoint emits zero delayed response. We therefore
rebuilt the task with both forms in training: two references have identical
horizon-1 responses but differ in persistence, and a cue names which the target
follows. Three-seed Qwen3-8B RL learns marginal persistence without reliable cue
following (App.~\ref{app:exp3-structure-selection}).

In a fixed component diagnostic crossing both domains and semantic or arbitrary
labels, all eight checkpoints identify the named reference perfectly.
Supplying exact response patterns does not rescue the 8B forecasts; an untrained
32B model shows partial numerical ability but often omits the resting baseline
(App.~\ref{app:exp3-components}). An exploratory single-seed supervised pilot on
4,800 prompts yields $97.9\%$ structure accuracy in Coin City/semantic versus
$50.0\%$ for matched RL, falling to $39.6\%$--$50.0\%$ under arbitrary labels
or in Coin Harbor. Trained natively it reaches $100.0\%$ in Harbor/semantic yet
stays at chance under arbitrary labels: the barrier is the label mapping, not
the domain. Rate and budget differ, limiting this to in-domain learnability,
not a general SFT advantage or transferable induction (App.~\ref{app:exp3-learnability}).

\noindent\textbf{Composition accuracy and evidence use are distinct.}
On the mechanism-family benchmark, episode-matched Qwen3-4B training beats a
\emph{fixed population-prior} control on held-out compositions: response-MAE
advantages are \eMechDiscEst{} \eMechDiscCi{} with the mechanism described and
\eMechUndiscEst{} \eMechUndiscCi{} with it inferred, both clearing all four
estimators at \eMechG{} seeds. A coherent gain intervention (two retained seeds) tracks better
than the prior on seen mechanisms but not on held-out ones
(App.~\ref{app:exp3-evidence-use}). A distribution-matched shuffled-target
control at \eShufG{} seeds now replaces the fixed prior: matched training keeps a
smaller advantage on the pre-registered stochastic endpoint
(\eShufUndiscStochEst{} \eShufUndiscStochCi{} inferred). The shuffle leaves
episode level near-intact while destroying response shape, so the gain falls
where the manipulation bites; held-out gain tracking stays unestablished
(App.~\ref{app:exp3-shuffled-control}).

\begin{figure}[!b]
\centering
\includegraphics[width=\linewidth]{exp4_qwen_scale_results.pdf}
\caption{\textbf{Training improves the tested bases, but scale alone does not
guarantee a better trained forecaster.} Across Qwen3 checkpoints from 1.7B to
32B, several trained models reach crowd-level performance, while the gain narrows
at the strongest base. Local model size is shown on a log parameter axis; hosted
systems at right are reference points. The Llama replication is in App.
Figure~\ref{fig:exp4-llama-scale-results}.}
\label{fig:sealed-market-grounding}
\label{fig:exp4-scale-results}
\end{figure}

\subsection{Exp.~4: external grounding in unseen historical markets}
\label{sec:exp4}

Historical markets test external forecasting accuracy without labeled mechanisms. We train on resolved markets using only
information available before each outcome, then evaluate 318 unseen questions
from different event families. Lower Brier scores mean less squared error
between forecast probabilities and realized outcomes
(App.~\ref{app:exp3b-protocol}).

Training lowers
error at every tested size, but the Qwen base-relative gap shrinks from 1.7B to
32B without steady improvement in trained accuracy (Fig.~\ref{fig:sealed-market-grounding}). At 32B, trained mean Brier is \(.1265\)
versus \(.1266\) for the base, the smallest observed gain; the trained-minus-base
interval \([-.00045,.00029]\) includes zero. A Llama 3B/8B
replication shows the same qualitative convergence pattern
(App. Figure~\ref{fig:exp4-llama-scale-results}). Several trained checkpoints
reach contemporaneous Polymarket performance. Hosted systems provide context,
not additional points on the local scaling curve.

\noindent\textbf{What Exp.~4 establishes.}
Outcome-based training improves mean forecasting accuracy on held-out event
families, with uncertain gains at the strongest Qwen base. This comparison is
trained versus base; without a shuffled-outcome control, it does not isolate the
contribution of learning episode-specific relationships. It complements the controlled transfer tests with external predictive validity.
%============================================================================
% END transfer_and_grounding.tex
%============================================================================
%============================================================================
% BEGIN related_work.tex
%============================================================================
\section{Related work: why accuracy is not enough}
\label{sec:related-work}

\noindent\textbf{Judgmental forecasting.}
Strong human forecasters combine base rates, decomposition, diverse evidence,
and repeated revision
\citep{tetlock2005expert,mellers2014psychological,tetlock2016superforecasting}.
LM benchmarks increasingly use time-indexed or leakage-controlled questions, and
retrieval or ensembling can bring systems near crowd accuracy
\citep{zou2022forecasting,karger2025forecastbench,halawi2024approaching,schoenegger2024wisdom}.
These results establish what probabilities models can produce. They do not tell
us what would make those probabilities change.

\noindent\textbf{Induction and relational reasoning.}
Transformers can infer functions from examples, and work on analogy and world
models asks whether learned representations capture relations or dynamics
\citep{garg2022what,chan2022data,webb2023emergent,yasunaga2024large,ha2018world,lecun2022path,li2023emergent}.
At the same time, shortcut research shows that surface success can fail under
structural change
\citep{geirhos2020shortcut,mitchell2023ai,vafa2025foundation}. A forecast can
therefore be right for the wrong reason: two models may report the same number
while relying on different populations, evidence, or pathways.

\noindent\textbf{Social agents and conditional validity.}
Generative agents can reproduce plausible individual and collective behavior
\citep{park2023generative,aher2023simulate,argyle2023out}, yet matching a
population average is not the same as preserving the relationships within that
population \citep{bisbee2024synthetic}. For systems intended to inform decisions,
we care about both the \emph{probability reported} and the \emph{relationship
that makes the evidence predictive}. Our experiments therefore move through
increasingly stringent tests: relevance and direction, contextual relationship
selection, revision from local evidence, induction of a novel convention, and
transfer across worlds. Each step rules out a weaker account of the same
forecast without assuming a general internal world model.
%============================================================================
% END related_work.tex
%============================================================================
%============================================================================
% BEGIN limitations.tex
%============================================================================
\section{Limitations: what these studies do and do not establish}
\label{sec:limitations}

These controlled tasks demonstrate context-dependent predictive relationships,
not human-like understanding or a general world model. Real settings contain
strategic interaction, unreliable cues, changing variables, and longer feedback
horizons absent from our market packets and constructed domains.

The internal intervention concerns one checkpoint and one convention; the
historical study covers one market cohort and training recipe. Additional
architectures, later cohorts, interactive settings, and misleading evidence
would test generality.

Performance is heterogeneous. Smaller checkpoints often struggle to induce the
KIV/ZOR meaning. The episode-matching advantage under joint domain and mechanism
shift does not survive the eight-seed replication; the benefit for compositions
of familiar mechanisms is checkpoint-specific and survives a
distribution-matched shuffled-target control only under stochastic decoding. Structure diagnostics use no target
observations; the supervised pilot has one seed and differs from RL in update
budget and, at the successful setting, learning rate. Semantic cue-to-pattern
learning could explain its in-domain success; native training places its
boundary at labels, not domain. The five-seed gain intervention establishes
evidence use on seen but not held-out mechanisms.
Historical-market gains narrow at the strongest Qwen base. The nonmonotonic
scale comparison uses one recipe; larger models may require different
optimization or data budgets.
%============================================================================
% END limitations.tex
%============================================================================
%============================================================================
% BEGIN conclusion.tex
%============================================================================
\section{Conclusion}
\label{sec:conclusion}

Forecast accuracy alone does not establish inductive reasoning. Models select
relevant evidence and revise contextual relationships, while training improves
forecasts across domains and held-out markets. A supervised pilot learns
response-form selection in distribution; its boundary is arbitrary labels, not
domain. Composition accuracy exceeds a shuffled-target control on seen
mechanisms without establishing transferred evidence use. These contrasts distinguish predictive accuracy from
reliable relationship use without presuming a general world model.
%============================================================================
% END conclusion.tex
%============================================================================
%============================================================================
% BEGIN submission_statements.tex
%============================================================================
\subsection*{AI use statement}

Language models produced the forecasts and experimental responses analyzed in this
work and assisted with software and manuscript preparation. They also generated
the context-reversal development vignettes. The authors verified
numerical claims against saved experimental records and executable analyses, checked
citations against original sources, reviewed the code, and take responsibility for
the submission and its artifacts.

\subsection*{Ethics statement}

The reported analyses use public market metadata and price histories, model
outputs, synthetic worlds, and an unpaid adult review of synthetic materials. The
review collected no identifying or sensitive information. Under our institutional
guidelines, it does not constitute human-subjects research; consequently, no IRB
review was sought. Counterfactual evidence packets were evaluated offline and were
not posted to markets or used for trading. We report aggregate model behavior
rather than recommend deployment.

\subsection*{Reproducibility statement}
Appendices~A--G specify the prompts, data-generating processes, sample construction,
censoring, estimators, deployment settings, missingness, hyperparameters, baselines,
and fail-closed checks. Builders, manifests, checksums, and saved generations support
each result. Experiment~3 reports the original three-seed Qwen3-8B, Qwen3-4B,
and Llama-3.1-8B roster, eight-seed Qwen extensions, a five-seed mechanism-family
study with its five-seed shuffled-target control, a three-seed structure-selection pilot, and exploratory component,
single-seed and native-cell supervised, and five-seed gain-intervention diagnostics. Parent-roster endpoints use
round 300 and five draws under a common stochastic decoder. The appendix
discloses the pilot's interim look at training monitors and the outcome-informed
supervised-launch amendment.
Experiment~4 evaluates Qwen and Llama checkpoints on family-disjoint historical
markets and reports hosted deployments as categorical references. Evaluation
examples are excluded from training; selection rules and exploratory analyses
are documented in the appendix.
%============================================================================
% END submission_statements.tex
%============================================================================
%============================================================================
% BEGIN acknowledgements.tex
%============================================================================
\subsubsection*{Acknowledgments}
Acknowledgments are redacted for double-blind review.
%============================================================================
% END acknowledgements.tex
%============================================================================

\clearpage
\bibliographystyle{iclr2027_conference}
\bibliography{references}

% All supplementary material lives in appendix.tex.
\clearpage
% Supplementary material for "Can language models forecast inductively?".
%
% This file holds every appendix section. Prose is edited here directly; the
% \input{tables/...} lines pull in generated analysis outputs, which their
% generators (see README.md) rewrite in place and which must not be edited here.
\appendix
\raggedbottom
% Keep appendix lists visually aligned with the text block rather than using the
% article class's comparatively deep default margins.
\setlist[itemize,1]{leftmargin=1.35em,labelsep=.45em}
\setlist[enumerate,1]{leftmargin=1.55em,labelsep=.45em}
\newcommand{\includeContextReversalPilot}{1}
%============================================================================
% BEGIN appendix_guide.tex
%============================================================================
\section*{Appendices}
\pdfbookmark[1]{Appendices}{appendix-guide}

The supplement is organized by experiment. A descriptive coin-flip elicitation
precedes the Experiment~1 evidence-response details. Experiment~2 documents
contextual relationship selection and revision, including a linear-probe
representation analysis; Experiments~3--4 document controlled and historical
transfer, replication, structure-selection diagnostics, supervised learnability,
and mechanism evidence use. Shared training details
and the artifact map close the appendix.

% Every entry below is derived from its \label: \ref* supplies the number,
% \nameref* the heading text, and \pageref* the page. Nothing here is typed by
% hand, so the guide cannot drift out of sync with the appendix headings the way
% a hard-coded list does. Adding or reordering an appendix section only requires
% adding its label to the list.
\begingroup
\hypersetup{hidelinks}
\small
\setlength{\parskip}{0pt}
\makeatletter
\newcommand{\appendixguidechapter}[1]{%
  \par\addvspace{0.45\baselineskip}%
  \noindent\hyperref[#1]{\textbf{\makebox[2.8em][l]{\ref*{#1}}\nameref*{#1}}}%
  \nobreak\hfill\nobreak
  \hyperref[#1]{\textbf{\pageref*{#1}}}\par
}
\newcommand{\appendixguidesection}[1]{%
  \@dottedtocline{2}{1.3em}{3.4em}%
    {\hyperref[#1]{\numberline{\ref*{#1}}\nameref*{#1}}}%
    {\hyperref[#1]{\pageref*{#1}}}%
}
\makeatother


\appendixguidechapter{app:carnival-coin-probe}

\appendixguidechapter{app:exp1}
\appendixguidesection{app:exp1-sample}
\appendixguidesection{app:exp1-protocol}
\appendixguidesection{app:metrics}
\appendixguidesection{app:extended-results}
\appendixguidesection{app:exp1-limitations}
\appendixguidesection{app:fig-code}
\ifdefined\includeContextReversalPilot
\appendixguidesection{app:exp1-context-reversal}
\fi

\appendixguidechapter{app:coin-city}
\appendixguidesection{app:coin-city-design}
\appendixguidesection{app:coin-city-arms}
\appendixguidesection{app:coin-city-execution}
\appendixguidesection{app:coin-city-analysis}
\appendixguidesection{app:coin-city-results}
\appendixguidesection{app:coin-city-mechanistic-probe}
\appendixguidesection{app:coin-city-limitations}

\appendixguidechapter{app:exp3}
\appendixguidesection{app:exp3-coin-structural}
\appendixguidesection{app:exp3-eight-seed}
\appendixguidesection{app:exp3-structure-selection}
\appendixguidesection{app:exp3-components}
\appendixguidesection{app:exp3-learnability}
\appendixguidesection{app:exp3-evidence-use}

\appendixguidechapter{app:exp4}
\appendixguidesection{app:exp3b-results}
\appendixguidesection{app:exp3b-update}

\appendixguidechapter{app:exp3-complete-training-setup}
\appendixguidesection{app:exp3-coin-structural-hparams}

\appendixguidechapter{app:exp3-repro}

\endgroup
%============================================================================
% END appendix_guide.tex
%============================================================================
\clearpage
%============================================================================
% BEGIN teaser_probe_appendix.tex
%============================================================================
\section{Descriptive coin-flip forecast elicitation}
\label{app:carnival-coin-probe}

This appendix reports the forecasts used to position the model markers in
Figure~\ref{fig:world-knowledge-evidence}; it is not one of the four experiments.
We queried GPT-5.4, Claude Opus~4.8, and DeepSeek V4-Pro eight times each
($24$ calls total) at temperature $0.7$ and a 600-token output ceiling. All
responses parsed. The figure places each model's logo at its eight-sample mean;
the green $\approx .68$ annotation rounds the pooled mean of $.679$.

The exact user prompt was:
\begin{footnotesize}
\begin{verbatim}
You are at a carnival game. A coin you have never seen before is flipped
twice in front of you, and it comes up heads both times. These two flips
are everything you know about this coin.

What is the probability that the next flip comes up heads?

Think about it intuitively, the way you actually would in the
moment -- do NOT use any formula, rule, or calculation, and do not
work out any numbers in your reasoning. Just reason in plain words
about what you'd expect and why. Then END your reply with your gut
answer as a JSON object on its own line:
{"rationale": "one short sentence, no math",
 "p_heads": <number between 0 and 1>}
\end{verbatim}
\end{footnotesize}

\begin{center}
\begin{minipage}{0.78\linewidth}
\centering
\captionof{table}{Forecasts underlying the three model markers in
Figure~\ref{fig:world-knowledge-evidence}. Each row summarizes eight independent
temperature-$0.7$ calls.}
\label{tab:carnival-coin-probe}
\small
\begin{tabular}{lccc}
\toprule
Model & Mean & Median & Range \\
\midrule
GPT-5.4 & $.644$ & $.670$ & $[.60,.67]$ \\
Claude Opus~4.8 & $.681$ & $.675$ & $[.65,.75]$ \\
DeepSeek V4-Pro & $.712$ & $.750$ & $[.60,.75]$ \\
\midrule
Pooled ($n=24$) & $.679$ & $.670$ & $[.60,.75]$ \\
\bottomrule
\end{tabular}
\end{minipage}
\end{center}

\noindent\textbf{What this probe does and does not separate.}
It has one condition. The $0.5$ and $1.0$ endpoints in
Figure~\ref{fig:world-knowledge-evidence} are reference policies rather than
measurements: no fair-coin arm and no no-carnival arm were collected, so the
probe cannot attribute the departure from $0.5$ to the carnival framing. It also
does not separate a prior about carnival coins from ordinary sequential updating
on an unknown bias. Laplace's rule of succession gives $3/4$ after two heads from
a coin with an unknown bias, using no world knowledge at all, and the observed
$.679$ sits below that rather than above it. Finally, the instruction not to use
any formula or calculation discourages explicitly deriving either the $0.5$ or
$0.75$ reference policy, changing the elicitation mode and complicating comparison
to both. We report the probe as a descriptive
illustration of the distinction the paper draws, not as evidence for it; the
four experiments carry the argument. \path{eval/run_coin_probe.py} implements
three further arms---a stipulated fair coin, an unfamiliar coin with no carnival,
and the carnival prompt without the no-calculation instruction---which would
supply the missing contrast.

The prompt, temperature, output limit, parsed probabilities, rationales, and raw
completions are stored in
\path{exp2_simulated_worlds/biased_news/data/coin_probe/coin_probe_2026-06-30.jsonl};
the collection and parsing code is
\path{exp2_simulated_worlds/biased_news/eval/run_coin_probe.py}. The prompt names
neither a fair nor a rigged coin. It requests qualitative reasoning while retaining
a numeric forecast for the plotted quantity.
%============================================================================
% END teaser_probe_appendix.tex
%============================================================================
%============================================================================
% BEGIN experiment1_appendix.tex
%============================================================================
% ============================================================
%  APPENDIX — EXPERIMENT 1 FULL DETAILS
%  Place after the main body, e.g. \appendix\section{...}
% ============================================================

\section{Experiment 1: prospective forecast updating}
\label{app:exp1}

Experiment~1 tests whether
forecast revisions follow evidence direction and whether directional packets
produce larger changes than orthogonal packets.
All reported counts refer to parseable records in the frozen consistency report
unless stated otherwise.

The main comparison reports the nine deployments with scale-valid archived
revisions. This appendix retains Qwen~2.5-7B for directional-compliance diagnostics
and documents why its magnitude-based statistics are excluded.

\subsection{Study design, sample, and agent roster}
\label{app:exp1-sample}

\noindent\textbf{Protocol summary.}
Experiment~1 crosses ten forecasting agents with a frozen set of unresolved binary
Polymarket questions.  The study has four stages: (1)~freeze the market sample;
(2)~obtain one or more initial forecasts from each agent; (3)~generate nine synthetic
evidence packets for each of the 100 core markets (three pro-$H_1$, three
anti-$H_1$, and three outcome-orthogonal); and (4)~ask the same agent to revise each
parseable initial forecast after receiving one packet at a time.  The unit of analysis is an \emph{agent $\times$ market $\times$ initial run
$\times$ counterfactual packet} update.  Requested repeat
counts were five initial runs for hosted agents and three for local agents; realized
valid counts are uneven because of provider failures, parsing failures, and targeted
top-ups.  Section~\ref{app:metrics} specifies how those records are normalized and
aggregated.

\subsubsection{Sampling frame and market selection}
\label{app:market-selection}

The primary sampling frame is a Polymarket API snapshot taken on June~9, 2026.
Selection was deterministic given that snapshot.  We retained markets that were
binary, active and unresolved, scheduled to resolve 14--180 days after sampling,
priced strictly between 0.10 and 0.90 on the YES contract, and supported by at least
\$1,000 in USD volume.  A keyword-based relevance filter retained politics,
government, elections, economics, technology, and geopolitics while excluding
sports, cryptocurrency, entertainment, and other out-of-scope topics.  When several
contracts belonged to the same event, we retained the highest-volume contract.

The recorded filter funnel makes the sampling loss explicit.  The snapshot contained
60,989 binary records; 37,386 were active, 8,639 passed the topic filter, 3,570 passed
the resolution-window filter, 1,403 passed the price filter, 839 passed the volume
filter, and event-level deduplication left 436 candidates.  A diversity pass selected
the 100-market core using a cosine-similarity threshold of 0.25 and a maximum of
16 markets per topical bucket.  The frozen core contains 16 global-election,
16 other, 15 Iran-nuclear, 12 U.S.-federal-policy, 10 AI/technology, 9 U.S.
state-election, 7 Russia--Ukraine, 5 Brazil, 5 U.S.-midterm, 3 macroeconomic,
and 2 China--Taiwan markets.

Counterfactual generation and Stage~3 updating target this 100-market core for
every model, but per-model coverage of it is not uniform. Llama~3.1-70B has 99
analyzable markets because one market has no parseable initial forecast, and
Claude Opus~4.8 has usable updates on all three packet arms in 73 markets;
Section~\ref{app:exp1-alternatives} reports the coverage table and the resulting
robustness checks.

\begin{center}
\begin{minipage}{\linewidth}
\centering
\captionof{table}{Representative markets from the frozen 100-market core.
Snapshot mid is the stored YES midpoint on June~9, 2026 rather than an eventual
outcome.}
\label{tab:markets}
\footnotesize
\resizebox{\linewidth}{!}{%
\begin{tabular}{llc}
\toprule
\textbf{Question} & \textbf{Topic} & \textbf{Snapshot mid} \\
\midrule
Will Renan Santos win the 2026 Brazilian presidential election? & Brazil & 0.17 \\
US--Iran nuclear deal by June~30? & Middle East & 0.20 \\
US $\times$ Iran permanent peace deal by July~31, 2026? & Iran & 0.32 \\
Will a new Gemini flagship be released by June~30, 2026? & AI/technology & 0.89 \\
\bottomrule
\end{tabular}%
}
\end{minipage}
\end{center}

\subsubsection{Agent roster and execution modes}
\label{app:agent-roster}

The roster comprises three hosted API systems and seven open-weight local systems
(Table~\ref{tab:agent-roster}).  The hosted models ran through Azure AI~Foundry and used
its native search-supported Responses workflow \citep{microsoft2026foundry}.  Local
models were served by Ollama \citep{ollama2026} on cluster GPUs; their research turn was supplied by Tavily search.  All agents received the same substantive three-step contract: build an event
model, gather current evidence, and emit a structured forecast. Provider-specific
search orchestration and model runtimes were not identical.

\begin{center}
\begin{minipage}{\linewidth}
\centering
\captionof{table}{Experiment~1 agent roster and requested initial-forecast repeats.  Requested
$k$ is the collection target rather than an analysis denominator; exact realized coverage
appears in Table~\ref{tab:model-summary}.}
\label{tab:agent-roster}
\small
\begin{tabular}{llllc}
\toprule
\textbf{Model} & \textbf{Runtime} & \textbf{Research search} &
\textbf{Nominal params} & \textbf{Requested $k$} \\
\midrule
GPT-5.4          & Azure Foundry & Foundry native & --- & 5 \\
Claude Opus~4.8  & Azure Foundry & Foundry native & --- & 5 \\
DeepSeek V4-Pro  & Azure Foundry & Foundry native & --- & 5 \\
Qwen~2.5-7B      & Ollama & Tavily & 7B  & 3 \\
Qwen~2.5-14B     & Ollama & Tavily & 14B & 3 \\
Qwen~2.5-32B     & Ollama & Tavily & 32B & 3 \\
Qwen~2.5-72B     & Ollama & Tavily & 72B & 3 \\
Llama~3.1-8B     & Ollama & Tavily & 8B  & 3 \\
Llama~3.1-70B    & Ollama & Tavily & 70B & 3 \\
Llama~3.3-70B    & Ollama & Tavily & 70B & 3 \\
\bottomrule
\end{tabular}
\end{minipage}
\end{center}

For local models, reasoning calls used the frozen configuration temperature of 0.1 and
the final JSON-formatting call used temperature 0.0.  Search queries, tool responses,
token counts, parse errors, and provider errors were retained with each run.  Some
smaller models returned probabilities on a 0--100 rather than 0--1 scale.  The
initial-forecast parser converted those values, but the frozen Stage~3 data retain a
mixed-scale Qwen~2.5-7B subset documented in
Section~\ref{app:record-processing}.

\subsection{Forecast, intervention, and validation protocol}
\label{app:exp1-protocol}

\subsubsection{Initial forecast elicitation}
\label{app:prompts}

\noindent\textbf{Design rationale.}
The elicitation asks each agent to restate the resolution criterion, construct
explicit YES and NO hypotheses, identify a status-quo trajectory, research
current evidence, and map its final YES probability to the posterior assigned to
$H_1$. Question restatement and status-quo anchoring follow prior
forecasting-prompt evidence \citep{halawi2024approaching,wilson2025automated}.

\noindent\textbf{Executed three-turn sequence.}
Every frozen initial forecast followed the same substantive sequence, with
provider-specific tool plumbing:

\begin{enumerate}
  \item \textbf{Event model.}  The agent received the question, full resolution
        description, days to resolution, and category.  It restated what would count
        as YES versus NO; identified actors, mechanisms, and latent variables; formed
        $H_1$ (YES) and $H_2$ (NO); and described the status-quo outcome.
  \item \textbf{Current-evidence research.}  The agent gathered evidence for and
        against each hypothesis, checked whether the required chain of events was
        plausible within the remaining time, and revised its assessment.  Hosted
        systems used Azure Foundry's native search workflow; local systems generated
        queries that were executed through Tavily and received the returned evidence
        in context.
  \item \textbf{Structured forecast.}  A final call converted the accumulated analysis
        to the common JSON contract below.  A run entered downstream processing only
        when this object was parseable.
\end{enumerate}

The research turn included this calibration instruction:

\begin{small}
\begin{verbatim}
Given the days remaining, what chain of events is required for H1, and
is that pace of change plausible? Give extra weight to the status quo;
depart substantially only when concrete directional evidence warrants it.
\end{verbatim}
\end{small}

The common output contract was:

\begin{small}
\begin{verbatim}
{
  "event_model": {
    "key_actors": [str],
    "key_mechanisms": [str],
    "latent_variables": [str]
  },
  "hypotheses": [
    {"id": "H1", "description": str,
     "prior_probability": float, "posterior_probability": float,
     "supporting_evidence": [str], "contradicting_evidence": [str]},
    {"id": "H2", ...}
  ],
  "hypothesis_to_forecast_mapping": str,
  "yes_prob": float,
  "confidence": "low"|"medium"|"high",
  "rationale": str,
  "evidence_sources": [str]
}
\end{verbatim}
\end{small}

The contract makes the intended invariant explicit: \(\hat p=P(H_1)\).  The ICS
metric tests whether the returned fields obey that invariant after an intervention.
The substantive prompt contract was common, but prompts were not byte-identical across
runtimes because the hosted and local search interfaces require different orchestration.

\subsubsection{Counterfactual packet construction and validation}
\label{app:cf-taxonomy}

For each of the 100 core markets, GPT-5.4 generated nine synthetic evidence packets:
three pro-$H_1$, three anti-$H_1$, and three orthogonal.  Each generation call received
the market question, resolution description, and a canonical parseable initial event
model and forecast.  It produced:

\begin{itemize}
  \item \textbf{evidence\_headline}: one sentence that states the new development;
  \item \textbf{evidence\_text}: a 2--4 sentence synthetic wire-style item;
  \item \textbf{mechanism\_targeted}: the causal pathway by which the item should
        affect, or remain irrelevant to, resolution;
  \item \textbf{direction}: \texttt{pro\_H1}, \texttt{anti\_H1}, or
        \texttt{orthogonal};
  \item \textbf{slot\_type}: one of the six roles in
        Table~\ref{tab:slot-taxonomy}; and
  \item \textbf{plausibility\_rating}: the generator's 1--5 self-rating.
\end{itemize}

Within each direction, sequential exclusion required a new packet to target a different
actor and mechanism from earlier packets.  Search was disabled so the packets remained
controlled synthetic premises rather than paraphrases of live reporting.

\begin{center}
\begin{minipage}{\linewidth}
\centering
\captionof{table}{Counterfactual slot taxonomy.  Each market receives one packet of every
direction--slot combination shown, for nine packets total.}
\label{tab:slot-taxonomy}
\small
\begin{tabular}{p{2.5cm} p{1.8cm} p{7.5cm}}
\toprule
\textbf{Slot type} & \textbf{Direction} & \textbf{Role in the intervention} \\
\midrule
DATA        & pro, anti & A measurement, price signal, statistic, or poll that
               shifts the likelihood of $H_1$. \\
DECISION    & pro, anti & An actor decision, official statement, policy announcement,
               commitment, or refusal bearing on $H_1$. \\
STRUCTURAL  & pro, anti & A technical, logistical, legal, or contextual change that
               makes $H_1$ more or less likely. \\
PROCEDURAL  & orthogonal & A same-domain administrative or process update without a
               direct causal path to $H_1$ versus $H_2$. \\
PERSONNEL   & orthogonal & A staffing or leadership change involving a relevant actor
               but an issue unrelated to $H_1$ versus $H_2$. \\
BACKDROP    & orthogonal & Contextual economic, logistical, or social information that
               sets the scene without shifting the outcome odds. \\
\bottomrule
\end{tabular}
\end{minipage}
\end{center}

\noindent\textbf{Packet-level integrity checks.}
The frozen packet file contains 900/900 parseable records: 100 markets multiplied by
nine packets, with exactly 300 pro-$H_1$, 300 anti-$H_1$, and 300 orthogonal records.
Every packet has a nonempty headline, and the logs contain zero search calls.  All 900
generator-assigned plausibility ratings equal 4/5. That field documents the
generator's compliance claim but provides neither variation nor independent
validation; the separate human materials review is reported below.
These checks verify structural completeness; they do not establish that directional
strength or difficulty is perfectly matched across packets.

As an illustration, for the Hormuz market the pro-$H_1$ headline
\textit{Oman says Iran agrees daily protected windows for commercial Hormuz transits}
elicited the largest single update reported in the viewer: Claude Opus~4.8 revised
from 10\% to 45\% (\(\Delta\hat p=+35\)~pp).  The orthogonal IMO
incident-reporting item produced negligible movement.

\subsubsection{Stage~3 update and pairing}
\label{app:update-protocol}

Each packet was presented separately with the agent's own parseable initial structured
forecast.  The agent was told to reason only from that forecast and the packet and to
return the same structured forecast with revised $H_1$, $H_2$, and YES probabilities.
Search was disabled.

\noindent\textbf{One deployment received an extra instruction.}
The Stage~3 template was not common across the roster.  GPT-5.4 alone was run with the
runner's \texttt{novelty} option, which appends a rule asking for a
\texttt{novelty\_assessment}: \emph{a single sentence stating whether this evidence was
likely already reflected in your initial forecast as background knowledge (suggesting a
small update) or is genuinely new and surprising (warranting a larger revision)}.
The frozen records show the asymmetry exactly: 3{,}838 of 3{,}921 GPT-5.4 updates carry
the field and no other deployment carries it in any record.  The instruction tells a
model to modulate revision size by perceived novelty, which is close to the construct
Section~\ref{app:exp1-alternatives} estimates, and GPT-5.4 is the deployment that
behaves differently on exactly that dimension: it returns the prior unchanged on
$8.4\%$ of orthogonal packets against $71.4\%$ and $44.4\%$ for the other two hosted
systems, writing a small revision where they abstain.  We cannot separate that
difference from the prompt difference, so GPT-5.4's abstention rate is not comparable to
the rest of the roster and we do not read the hosted ordering on abstention as a model
difference.  Its conditional magnitude, directional correctness, and AUC use the same
scoring as every other deployment.

For an agent--market cell with \(k\) parseable initial runs, the intended crossed design
contains \(9k\) update records.  Each update is conditionally independent in prompt
context: it contains one initial run and one packet, never another packet or another
agent's answer.  The novelty field is not a reported endpoint and should not be
viewed as a validated correction.

Provider interruptions, parse failures, and subsequent top-ups make realized cells
uneven.  We therefore report exact valid record and market denominators by model in
Table~\ref{tab:model-summary} rather than treating requested \(k\) or the theoretical
\(9k\) as observed sample sizes.

\subsection{Estimands, record processing, and aggregation}
\label{app:metrics}

\subsubsection{Record construction, deduplication, and missingness}
\label{app:record-processing}

The frozen analysis artifact is
\path{exp1_prospective/data/results/consistency_report_2026-06-17.json}; tables in
this paper use its stored denominators rather than a fresh scan of the evolving
collection directories.  During report construction, date suffixes are stripped from
task identifiers so daily collection files join to the same market.  Counterfactuals
are deduplicated by \texttt{cf\_id}; updates are deduplicated by
\texttt{update\_id}, with a parseable record preferred over an error record when both
exist.  Files are read newest first for the initial-forecast and market-price lookup.
Malformed JSON lines are skipped.  Missing values are handled per estimand: a record
may therefore contribute to ICS but not EHC, or vice versa, depending on which fields
parse.

Stage~3 update fields were evaluated on their stored scale in the frozen report.
For Qwen~2.5-7B, 207 of 8,636 deduplicated records store
\texttt{updated\_yes\_prob} above 1.5; no other model does. Its directional sign
metrics remain usable, but its absolute revision magnitudes and sensitivity
ratio are scale-contaminated and are omitted from the paper tables.

\subsubsection{Consistency estimands}
\label{app:metric-definitions}

For update record \(i\), let \(d_i\in\{+1,-1\}\) denote the intended direction
(\(+1\) for pro-$H_1$, \(-1\) for anti-$H_1$).  Orthogonal packets do not have an
intended sign and are excluded from EHC and HFC.  The 3~pp movement threshold excludes
cases in which a model effectively leaves the forecast unchanged.

\noindent\textbf{Evidence--hypothesis consistency (EHC).}
Let \(\Delta p(H_1)_i=p(H_1)_i^{\mathrm{post}}-\hat p_i^{\mathrm{initial}}\).
For directional records with \(|\Delta p(H_1)_i|\geq0.03\),
\begin{equation}
  \mathrm{EHC}_i=\mathbf{1}\!\left[\Delta p(H_1)_i d_i>0\right].
  \label{eq:ehc}
\end{equation}
EHC asks whether the updated hypothesis field moves in the direction implied by the
packet; it does not score the size of that movement.

\noindent\textbf{Hypothesis--forecast consistency (HFC).}
Let \(\Delta\hat p_i=\hat p_i^{\mathrm{post}}-\hat p_i^{\mathrm{initial}}\).
For directional records with \(|\Delta\hat p_i|\geq0.03\),
\begin{equation}
  \mathrm{HFC}_i=\mathbf{1}\!\left[\Delta\hat p_i d_i>0\right].
  \label{eq:hfc}
\end{equation}
HFC asks the same directional question of the headline YES probability. Because
the output contract requires this probability to equal $P(H_1)$, we retain HFC in
the appendix as a compliance diagnostic rather than as a separate main-text result.

\noindent\textbf{Internal coherence score (ICS).}
ICS is defined whenever both updated fields parse and is unconditional on movement:
\begin{equation}
  \mathrm{ICS}_i=\mathbf{1}\!\left[
  \left|\hat p_i^{\mathrm{post}}-p(H_1)_i^{\mathrm{post}}\right|<0.02\right].
  \label{eq:ics}
\end{equation}
The inequality is strict: a discrepancy of exactly 2~pp does not pass.

\noindent\textbf{Directional sensitivity.}
Packet selectivity is summarized by
\begin{equation}
  \mathrm{Sensitivity}=
  \frac{\overline{|\Delta\hat p|}_{\mathrm{pro}+\mathrm{anti}}}
       {\overline{|\Delta\hat p|}_{\mathrm{orthogonal}}}.
  \label{eq:sensitivity}
\end{equation}
A value near one indicates comparable movement for directional and orthogonal
packets; larger values indicate preferential response to directional packets. This
ratio is descriptive.  It cancels a uniform multiplicative scale within a model, but
does not repair mixed scales across records, hence the Qwen~2.5-7B caveat above.

\subsubsection{Aggregation and uncertainty}
\label{app:aggregation}

For EHC, HFC, and ICS, the evaluator first pools eligible records within each
model--market cell to obtain a market-level rate.  The reported model score is the
unweighted mean of those market rates, so a market with more successful repeat runs
does not receive greater weight.  Reported standard errors are the sample standard
deviation of market-level rates divided by the square root of the number of markets.

Revision magnitudes and sensitivity retain the record-pooled point estimands.
Uncertainty uses 20,000 percentile cluster-bootstrap draws over represented markets.
All nine packets and every repeated initial-run update from a selected market remain
together, and sensitivity is recomputed from pooled directional and orthogonal
movement within each draw. This accounts for dependence among records from the same
market while preserving the published estimand.

\subsubsection{Output compliance and threshold robustness}
\label{app:exp1-threshold-robustness}

\noindent\textbf{Output-field coupling.}
The final-output schema explicitly instructs the agent that
$\hat p=P(H_1)$, and the Stage~3 prompt requests the same structure after
updating. EHC and HFC are therefore not independent tests of an inferred
hypothesis-to-forecast link. ICS is a schema-adherence diagnostic, while close
EHC--HFC agreement reflects a mixture of directional updating and compliance
with the prompted equality. Among the 26,010 relevant records with both updated
fields present, 25,984 (99.90\%) satisfy the equality exactly after scale
normalisation and 25,987 (99.91\%) satisfy it within the 2-pp ICS tolerance. We
therefore do not treat EHC--HFC parity as independent evidence for a latent
representation or for spontaneous propagation between separately elicited
quantities.

\noindent\textbf{Unconditional estimand.}
For threshold $t$, Table~\ref{tab:exp1-threshold-robustness} retains a valid
directional update when $|\Delta|\geq t$ and reports the fraction for which
$d_i\Delta_i>0$. At $t=0$, every valid pro- and anti-$H_1$ update enters the
denominator and a zero change is scored incorrect. We first compute correctness
within market and then average markets equally, matching the primary analysis.
The \emph{Below 3 pp} column is the record-weighted fraction and count excluded
by the published 3-pp conditional metric. The table's $n$ is the complete
pre-threshold denominator: records require a parseable initial probability and
the applicable updated field, but are not selected on the direction or magnitude
of movement.

\begin{center}
\begin{minipage}{\linewidth}
\centering
\captionof{table}{Unconditional directional correctness and robustness to 0-,
1-, 3-, and 5-percentage-point movement thresholds. The 0-pp column is
unconditional over all valid directional updates and scores zero movement as
incorrect; positive-threshold columns condition on the stated magnitude.
\emph{Below 3 pp} gives the percentage and count omitted from the corresponding
published 3-pp EHC or HFC denominator. Rates are macro-averaged across markets.}
\label{tab:exp1-threshold-robustness}
\scriptsize
\setlength{\tabcolsep}{4pt}
\resizebox{\linewidth}{!}{%
\input{tables/exp1_threshold_robustness}
}
\end{minipage}
\end{center}

\noindent\textbf{Threshold-sweep results.}
The threshold sweep preserves the qualitative conclusion but narrows its
wording. Hosted-model unconditional correctness is 97.3--98.4\%; at the
published 3-pp threshold it is 97.7--99.0\%. The 3-pp rule excludes 6.9\% of
GPT-5.4, 12.6--12.7\% of DeepSeek V4-Pro, and 18.1\% of Claude Opus~4.8
directional updates. Thus, the near-1.0 EHC/HFC values do not result from
discarding a majority of cases, and the conclusion is stable across 0, 1, 3,
and 5~pp. The precise statement is \emph{99\% directionally correct among updates
of at least 3~pp}. It is not \emph{99\% coherent updating}. Small local
models have fewer sub-3-pp updates (0.6--2.9\%) but lower unconditional
correctness (88.1--96.4\%), so their errors are not hidden primarily by the
movement filter.

\noindent\textbf{Normalisation and reproducibility.}
Before recomputing deltas, the analysis independently normalises the initial
headline probability, updated headline probability, and updated $H_1$
posterior from the documented $[0,1]$ or $[0,100]$ formats. This matters for
166 relevant Llama~3.1-8B records in which $H_1$ was emitted as a percentage
while \texttt{yes\_prob} was emitted as a proportion. The complete eligible
record and market counts at every threshold, invalid-field counts, market-level
standard errors, input-file SHA-256 hashes, and full-precision estimates are in
\path{exp1_prospective/data/results/threshold_robustness.json}. The table is
regenerated by
\path{exp1_prospective/agent/evaluate_threshold_robustness.py}; neither script
nor table filters items using realized outcomes.

\subsection{Extended results}
\label{app:extended-results}

\subsubsection{Analysis coverage}
Table~\ref{tab:model-summary} gives the exact Stage~3 record and market counts.
The market column counts markets with at least one valid update. Coverage of all
three packet arms is not uniform across models: Llama~3.1-70B has 99 analyzable
markets because one market has no parseable initial forecast, and Claude Opus~4.8
has all three arms in 73 of 100. Section~\ref{app:exp1-alternatives} gives the
per-arm coverage and assesses whether missingness is arm-selective.

\begin{center}
\begin{minipage}{\linewidth}
\centering
\captionof{table}{Valid Stage~3 update records and represented core markets by
agent.  Hosted systems first; open-weight rows in descending parameter count.}
\label{tab:model-summary}
\small
\begin{tabular}{llrr}
\toprule
\textbf{Model} & \textbf{Type} & \textbf{Records} & \textbf{Markets} \\
\midrule
Claude Opus~4.8 & Hosted  & 3{,}125 & 100 \\
GPT-5.4         & Hosted  & 3{,}913 & 100 \\
DeepSeek V4-Pro & Hosted  & 4{,}395 & 100 \\
Qwen~2.5-72B    & Local 72B & 2{,}692 & 100 \\
Llama~3.3-70B   & Local 70B & 2{,}682 & 100 \\
Llama~3.1-70B   & Local 70B & 2{,}490 &  99 \\
Qwen~2.5-32B    & Local 32B & 2{,}693 & 100 \\
Qwen~2.5-14B    & Local 14B & 2{,}700 & 100 \\
Llama~3.1-8B    & Local 8B  & 5{,}404 & 100 \\
Qwen~2.5-7B     & Local 7B  & 8{,}636 & 100 \\
\midrule
\textbf{Total}  &           & \textbf{38{,}730} & \\
\bottomrule
\end{tabular}
\end{minipage}
\end{center}

\subsubsection{Prospective settlement audit}
\label{app:exp1-outcome-followup}

The frozen-cohort evaluation includes a dated settlement audit that leaves the
directional-updating estimand unchanged. On August~24, 2026, we queried the public
Polymarket market endpoint for the exact 100 frozen IDs. We admitted a label only
when the market was closed and the YES/NO token prices were exactly
\([1,0]\) or \([0,1]\).  This strict rule yielded 60 scoreable markets
(21 YES and 39 NO). One additional closed market had an invalid \(0.5/0.5\)
settlement and 39 remained unresolved at the audit cutoff; all 40 were excluded.

For each agent--market cell, the forecast is the median of its parseable repeated
initial YES probabilities, matching the maintained Experiment~1 viewer.  The crowd
benchmark is the YES price stored in the frozen June~9 selection record.  We report
threshold accuracy (YES iff \(p>0.5\)), Brier score, and natural-log loss.  Proper-score
contrasts are paired agent-minus-market differences; negative values favor the
agent.  Percentile intervals use 20,000 market-level bootstrap resamples.

\begin{center}
\begin{minipage}{\linewidth}
\centering
\captionof{table}{Descriptive accuracy on the first 60 strictly settled markets.
A/M denotes agent/June~9 market on the same rows; Correct is the number of
thresholded calls matching the outcome.  Lower Brier and log loss are better, and
each \(\Delta\) is agent minus market.  Llama~3.1-70B lacks a parseable initial
forecast for one settled market, so its paired denominator is 59.}
\label{tab:exp1-resolved-outcomes}
\scriptsize
\resizebox{\linewidth}{!}{%
\input{tables/exp1_resolved_outcomes}
}
\end{minipage}
\end{center}

The raw correct-call counts are close for the hosted agents and crowd:
Claude Opus~4.8 gets 41/60 correct, GPT-5.4 and DeepSeek V4-Pro each get
40/60, and the June~9 market side gets 39/60.  GPT-5.4 has the lowest Brier
score (\(.225\) versus \(.246\) for the market), a paired difference of
\(-.021\;[-.074,.030]\); its log-loss difference is
\(-.037\;[-.178,.103]\).  Qwen~2.5-7B has the lowest log loss (\(.651\)
versus \(.700\)), but gets 38/60 thresholded calls correct.  No agent's paired
Brier or log-loss interval excludes zero.

This audit describes the outcome-ready slice of the frozen cohort. Settlement by
the audit date censors the cohort toward shorter-horizon and early-resolving
questions, so the 60 are not a random subset of the frozen 100.  Moreover, the
benchmark is the June~9 snapshot whereas agent forecasts were collected from
June~10 through June~17; later-running agents may have had additional public
information.  The result therefore does not establish superiority to the crowd.
It documents the prospective check; the directional-update estimates do not, by
themselves, imply a detectable accuracy advantage.

\subsubsection{Revision magnitudes and packet selectivity}
Table~\ref{tab:anchoring} reports mean absolute forecast revision
\(\overline{|\Delta\hat p|}\) by packet direction.  For the nine models whose updates
are consistently stored on 0--1, hosted-system sensitivity is $8.2$--$24.0\times$,
compared with $1.9$--$6.4\times$ for local systems. These are the unconditional
ratios; Table~\ref{tab:exp1-abstention} separates the abstention and magnitude
factors behind them, and the main text reports those two rather than the product.

Qwen~2.5-7B has mixed 0--1 and 0--100 stored updates
(Section~\ref{app:record-processing}); we therefore omit its absolute magnitudes
and sensitivity ratio rather than present them as comparable estimates.

\begin{center}
\begin{minipage}{\linewidth}
\centering
\captionof{table}{Mean absolute forecast revision by counterfactual direction.
Hosted systems first; open-weight rows in descending parameter count.
Sensitivity is mean directional divided by mean orthogonal movement; \(n\) is the
valid Stage~3 record count. Values are percentage points for models with
consistently stored 0--1 updates.}
\label{tab:anchoring}
\small
\begin{tabular}{llccccc}
\toprule
\textbf{Model} & \textbf{Type} & pro-$H_1$ & anti-$H_1$ & orthogonal
  & \textbf{Sensitivity} & $n$ \\
\midrule
Claude Opus~4.8 & Hosted  & 12.7\,pp & 9.6\,pp  & 0.5\,pp & $24.0\times$ & 3{,}125 \\
DeepSeek V4-Pro & Hosted  & 13.9\,pp & 11.4\,pp & 1.4\,pp & $8.9\times$  & 4{,}395 \\
GPT-5.4         & Hosted  & 10.6\,pp & 8.5\,pp  & 1.2\,pp & $8.2\times$  & 3{,}913 \\
\midrule
Qwen~2.5-72B    & Local 72B & 15.0\,pp & 10.2\,pp & 3.0\,pp & $4.2\times$  & 2{,}692 \\
Llama~3.3-70B   & Local 70B & 14.8\,pp & 14.3\,pp & 2.3\,pp & $6.4\times$  & 2{,}682 \\
Llama~3.1-70B   & Local 70B & 14.8\,pp & 15.5\,pp & 3.3\,pp & $4.6\times$  & 2{,}490 \\
Qwen~2.5-32B    & Local 32B & 15.0\,pp & 10.3\,pp & 4.0\,pp & $3.2\times$  & 2{,}693 \\
Qwen~2.5-14B    & Local 14B & 14.8\,pp &  8.6\,pp & 4.1\,pp & $2.9\times$  & 2{,}700 \\
Llama~3.1-8B    & Local 8B  & 11.1\,pp & 12.0\,pp & 5.9\,pp & $1.9\times$ & 5{,}404 \\
Qwen~2.5-7B     & Local 7B  & --- & --- & --- & --- & 8{,}636 \\
\bottomrule
\end{tabular}
\end{minipage}
\end{center}
Qwen~2.5-7B magnitude cells are omitted because 207 deduplicated updates are
stored above 1.5, mixing probability scales within the row.
Table~\ref{tab:anchoring} documents the magnitude contrast; Figure~\ref{fig:causal-coherent} then places directional consistency and packet selectivity on a common visual scale.

\begin{center}
\begin{minipage}{\linewidth}
\centering
\includegraphics[width=\linewidth]{figures/exp1_ehc_sensitivity_combined.pdf}
\captionof{figure}{\textbf{Directional consistency and packet selectivity by
deployment.} \emph{Left:} EHC is $87$--$90\%$ for 7--8B local models and
$98$--$99\%$ for hosted systems. \emph{Right:} sensitivity is
$1.9$--$6.4\times$ for local models and $8.2$--$24.0\times$ for hosted systems;
Qwen~2.5-7B is omitted from this panel because its archived absolute revisions mix
0--1 and 0--100 scales and therefore do not yield a comparable ratio. The panel
plots the unconditional ratio, which Table~\ref{tab:exp1-abstention} decomposes.}
\label{fig:causal-coherent}
\end{minipage}
\end{center}






\subsection{Study limitations}
\label{app:exp1-limitations}


\subsubsection{Sampling and outcome scope}
The sample is a filtered slice of liquid, medium-horizon Polymarket questions
rather than a probability sample of forecasting tasks. Topic exclusions, the
0.10--0.90 price window, the 14--180 day horizon, event deduplication, and the
diversity cap shape its difficulty distribution.  The primary Experiment~1
estimand evaluates conditional updating on questions that were unresolved at
collection.  The settlement audit in
Section~\ref{app:exp1-outcome-followup} is naturally censored to the first 60
strict settlements and is descriptive.  Experiment~4 separately evaluates
accuracy on a locked sample of resolved historical markets with synchronized
forecast cutoffs and market-price baselines.

\subsubsection{What the update prompt withholds}
\label{app:exp1-withholding}
The Stage~3 template fills in one packet field, \texttt{evidence\_text}. The packet
record also holds the intended \texttt{direction}, \texttt{target\_hypothesis},
\texttt{expected\_hypothesis\_shift}, \texttt{slot\_type}, \texttt{evidence\_headline},
\texttt{mechanism\_targeted}, and a free-text \texttt{rationale}. None of these
reaches the agent; all are read only when the update is scored. The agent sees its
own prior forecast, the snippet, and the fixed instruction.

We checked this two ways. Over the 12{,}728 frozen records that kept their prompt
string, spanning all three hosted systems, the evidence block is identical to the
packet's \texttt{evidence\_text} and contains no private packet field. Prompt
persistence was a per-runner logging choice, so this covers a subset of records by
logging coverage rather than by outcome; the runners that stored nothing use the
same one-field template. Turning to the snippets, none of the 900 names $H_1$,
$H_2$, the market, or its resolution. Twenty-two (2.4\%, 10 pro-$H_1$ and 12
anti-$H_1$) use any likelihood wording---\emph{odds}, \emph{chance of},
\emph{more}/\emph{less likely}, \emph{probability}, \emph{likelihood}---and each
time it describes something in the world, a simulation percentage or a betting
line, not the market's resolution. The sign is not stated; it must be inferred from the relation between the
reported event and the market outcome. Reproduced by
\path{exp1_prospective/agent/audit_packet_polarity.py}.

\subsubsection{Alternative explanations for the selectivity gap}
\label{app:exp1-alternatives}
Here we state the competing accounts of the movement gap and what we find for each.
All checks come from \path{exp1_prospective/agent/audit_selectivity_robustness.py},
\path{audit_orthogonal_abstention.py}, \path{audit_lexical_sign_baseline.py}, and
\path{audit_group_contrast.py}.

\textbf{Where the numbers come from.} The checks read the frozen June~17 report,
which enumerates every scored run, so they use exactly the data behind the
published numbers. The enumeration matches the report's own per-arm counts for
seven of the nine models; Claude Opus~4.8 and GPT-5.4 carry 14 and 4 extra runs in
markets absent from the per-market join, 18 of 38{,}730. Reproduced sensitivity
ratios agree with the published ones to within $0.2$. Three checks need packet
text and so join runs to the June~10 packet set; 346 runs (1.1\%) do not join
because they come from a June~7--8 pilot against a superseded 63-packet
generation, and they are dropped from those three checks only.

\textbf{Attrition is much higher for the hosted systems.} Claude Opus~4.8 made
9{,}313 calls and yielded 3{,}232 usable updates, losing 3{,}656 to broken pipes,
914 to HTTP~401s after a key expired mid-run, and 1{,}481 to unparseable
responses. Of those, 1{,}461 come from one malfunctioning June~15 retry pass that
failed at 91\%; the main June~10 collection parsed at 99.3\%, so the published
compliance figures are not survivorship. DeepSeek~V4-Pro made 6{,}338 calls for
4{,}598 usable updates, losing 1{,}629 to HTTP~401s; GPT-5.4 made 4{,}269 for
4{,}057. The six local models lose almost nothing---four record no failed calls,
the other two lose 8 and 15 responses to parsing. Content-based exclusion is
confined to 28 GPT-5.4 calls across 2 markets. Recorded failures are predominantly
transport or parsing failures; the coverage check below assesses arm and market
imbalance, not content-specific selection. Separately, two
update streams were corrupted by concurrent writes and repaired in August~2026,
after the freeze; the original bytes and unrecoverable fragments sit in
\path{exp1_prospective/data/_recovery_audit/}. Nothing in this section reads those
files.

\textbf{The generator saw one agent's world model.} All packets were built from a
single canonical initial forecast, and it is DeepSeek~V4-Pro's: its June~10 YES
probability matches the stored generation prior on 100 of 100 markets, where no
other agent matches on more than 28. For the other eight models the evidence came
from a different system's event model. The seed model gains nothing by it---among
hosted systems it is second on sensitivity ($8.9\times$ against Claude's
$24.0\times$), third on AUC ($.967$), and last on EHC ($.977$)---but its estimates
are the least independent in the roster.

\textbf{Coverage is uneven across arms.} Table~\ref{tab:exp1-arm-coverage} gives
per-arm market coverage. Seven models cover 99 or 100 markets on every arm. Claude
Opus~4.8 covers 82 on the orthogonal arm and 77 on the directional arms, 73 in
common, so its pooled ratio divides means taken over slightly different market
sets. Restricting to its 73 common markets gives $23.7\times$ against $24.2\times$
pooled. As a proxy check, every other model has similar sensitivity on the markets
Claude covers and those it misses ($8.4$ vs.\
$7.6$ for GPT-5.4, $9.7$ vs.\ $7.2$ for DeepSeek~V4-Pro, $4.3$ vs.\ $4.2$ for
Qwen~2.5-72B, $1.9$ vs.\ $2.1$ for Llama~3.1-8B).

\textbf{The ratio has an unstable denominator.} Claude's orthogonal mean is
$0.5$~pp, so its $24\times$ comes from a small divisor rather than unusually large
directional movement, and its market-clustered interval is correspondingly wide,
$[17.3,36.9]$. Table~\ref{tab:exp1-selectivity-robustness} reports three
statistics that avoid that division. The difference of means is $5.6$--$12.3$~pp
and does not order the models the way the ratio does. The within-market
probability that a directional packet moves the forecast more than an orthogonal
one runs from $.791$ for Llama~3.1-8B to $.994$ for GPT-5.4, putting GPT-5.4 above
Claude Opus~4.8 and reversing the ratio ranking. Every model moves further on
directional packets in every market it covers.

\textbf{The denominator is mostly zeros.} That instability has a specific cause,
and naming it replaces the ratio with two statistics that are easier to read. A
revision of exactly zero---the model reproduces the probability shown to it---enters
the mean like any other record, so the ratio silently mixes how often a
deployment declines to move with how far it moves when it does. Writing $A$ for
the share of runs with $\Delta\hat p=0$ and $M$ for the mean absolute revision
among the rest, $\overline{|\Delta\hat p|}=(1-A)M$, and the ratio factors exactly:
\begin{equation}
  \mathrm{Sensitivity}
  =\underbrace{\frac{M_{\mathrm{dir}}}{M_{\mathrm{orth}}}}_{\text{magnitude}}
  \times\underbrace{\frac{1-A_{\mathrm{dir}}}{1-A_{\mathrm{orth}}}}_{\text{abstention}}.
  \label{eq:sensitivity-decomposition}
\end{equation}
Table~\ref{tab:exp1-abstention} gives both factors. Directional abstention is
$0.0$--$1.4\%$ for every deployment, so the second factor is governed entirely by
the orthogonal arm, where abstention runs from $1.1\%$ for Llama~3.1-8B to
$71.4\%$ $[63.8,78.5]$ for Claude Opus~4.8. Claude's $24\times$ is therefore
$7.0\times$ $[5.5,9.1]$ in magnitude multiplied by $3.48$ in abstention;
DeepSeek~V4-Pro's $8.9\times$ is $4.9\times$ times $1.79$; GPT-5.4's $8.2\times$
is $7.5\times$ times $1.09$, because GPT-5.4 writes a small revision where the
other two hosted systems return the prior unchanged. Across the roster the
published $1.9$--$24.2\times$ becomes $1.9$--$7.5\times$ once the zeros are
separated out, and the ordering changes: GPT-5.4 rather than Claude Opus~4.8 is
highest, and Llama~3.3-70B at $5.7\times$ sits above DeepSeek~V4-Pro at
$4.9\times$.

We report both factors rather than the product. Declining to revise at all is an
appropriate response to an outcome-irrelevant premise, so the abstention rate is
a selectivity measure in its own right and a better-behaved one: it is a bounded
proportion, it is paired within model across arms, and it needs no division by a
small number. It is also the more fragile of the two. An exact copy of a
displayed prior is a discrete act adjacent to formatting, and the update prompt
does instruct the agent to leave unaffected fields unchanged, so we do not read a
high rate as evidence about internal computation. The no-information condition
described in Section~\ref{app:exp1-limitations} is what would separate declining
to revise from copying the anchor, and it was not run. Reproduced by
\path{exp1_prospective/agent/audit_orthogonal_abstention.py}.

\begin{center}
\begin{minipage}{\linewidth}
\centering
\captionof{table}{Decomposition of the sensitivity ratio into magnitude and
abstention factors (Equation~\ref{eq:sensitivity-decomposition}). \emph{Orthogonal
abstention} is the share of orthogonal runs with $\Delta\hat p$ exactly zero; the
directional-arm rate is $0.0$--$1.4\%$ throughout. The two \emph{given a move}
columns condition on a nonzero revision, and their ratio is the magnitude factor.
Brackets are 95\% market-clustered percentile intervals from 20,000 draws, the
scheme used for the published movement intervals. Qwen~2.5-7B is omitted for
mixed probability scales.}
\label{tab:exp1-abstention}
\scriptsize
\setlength{\tabcolsep}{4pt}
\resizebox{\linewidth}{!}{%
\input{tables/exp1_abstention_decomposition}
}
\end{minipage}
\end{center}

\textbf{The deployment contrast depends on the statistic.} Comparing the three
hosted systems against the six scale-valid open-weight models by market-clustered
bootstrap (20{,}000 replicates, macro-averaged within group so a model with more
runs does not dominate), the statistics disagree. On the ratio the
difference is $+9.8$ $[7.2,14.2]$, group means $13.7\times$ and $3.9\times$, and
the per-model ranges do not overlap. On AUC it is $+.070$ $[.062,.079]$, which
excludes zero, but the ranges do overlap: Llama~3.3-70B at $.989$ and
Llama~3.1-70B at $.978$ both beat DeepSeek~V4-Pro at $.967$. On raw movement there
is no difference at all, $+0.8$ $[-0.0,1.7]$~pp around means of $10.1$ and $9.3$.

The decomposition explains why the ratio is the one statistic that separates the
groups completely. Both of its factors favour the hosted systems on average and
neither separates them on its own: orthogonal abstention differs by $+.286$
$[.244,.329]$ around group means of $.413$ and $.126$, with four open-weight
deployments at or above GPT-5.4's $8.4\%$; the magnitude factor differs by
$+3.10$ $[2.31,4.08]$ around $6.5\times$ and $3.4\times$, with Llama~3.3-70B
above DeepSeek~V4-Pro. Multiplying two partially separating factors yields a
product whose ranges happen not to overlap, so we do not report that non-overlap
as a finding. The two classes move about equally far on directional evidence; the
hosted systems more often decline to move at all on controls, and move somewhat
less far when they do. We state the deployment result that way rather than as a
general selectivity ordering. Group contrasts come from
\path{exp1_prospective/agent/audit_group_contrast.py}.

\textbf{The arms differ in wording.} By construction directional packets take the
DATA, DECISION, and STRUCTURAL slots and orthogonal packets take PROCEDURAL,
PERSONNEL, and BACKDROP, and the text differs accordingly: directional snippets
average $96.8$ words and $7.1$ digits, orthogonal ones $85.7$ and $3.7$. Movement
could track density rather than content. Two checks argue otherwise. Within the
directional arm, where length cannot indicate sign, Spearman correlation between
word count and $|\Delta\hat p|$ is $-.06$ to $.00$ and between digit count and
$|\Delta\hat p|$ is $.01$ to $.10$. And the arms' word counts almost fully overlap
(orthogonal 66--114, directional 66--121); restricting to that band keeps 889 of
900 packets and moves every sensitivity estimate by at most $0.1$
(Table~\ref{tab:exp1-selectivity-robustness}, last column). The arms are not
matched on surface form. These checks show that word and digit counts do not
explain the effect; other surface differences remain. Directional correctness is
not explained by length, since length cannot say which way to move.

\textbf{A shallow classifier might be enough.} We tested it. Multinomial naive
Bayes and L2-regularised logistic regression were fit on the 900 packet texts with
five-fold cross-validation grouped by market, so no market appears on both sides
of a fold. Telling directional from orthogonal packets is nearly free: naive Bayes
reaches $.989$ against a $.667$ majority rate. That follows from the slot-type
plan rather than revealing anything---the most distinctive orthogonal tokens are
\emph{appointed}, \emph{appointment}, \emph{oversee}, and \emph{accreditation},
which name the roles the generator was told to write. So a model could route on
that marking without weighing content, and the sensitivity ratio alone cannot rule
it out. Matching the control arm on evidence type would need a different packet
design; the checks above bound only the length part of the problem.

Recovering the \emph{sign} is another matter. On pro-$H_1$ versus anti-$H_1$ the
same classifiers reach $.688$ and $.675$ against $.500$ chance, and prepending the
market question does not help ($.665$ and $.680$): these bag-of-words models do
not combine question polarity with event valence. The tokens it keys on carry no
consistent sign---\emph{vessel}, \emph{venue}, \emph{signature} for pro-$H_1$,
\emph{table}, \emph{materials}, \emph{findings} for anti-$H_1$. Against this $.69$ lexical baseline, the human panel reaches $.762$ on its
smaller categorical task and the hosted agents $.973$--$.984$ on probability
revision. These are not matched performance comparisons. The agent--classifier gap
is evidence against this bag-of-words account, not proof that no surface heuristic
can recover sign; unlike the arm contrast, it stays within pro- versus anti-$H_1$
packets.

\textbf{Does the stated reasoning track the packet?} If movement were routed on
surface marking, a model's written rationale would not need to bear any particular
relation to what it does with the forecast. We tested two versions of this
(\path{exp1_prospective/agent/audit_mechanism_tracking.py}).

The first asks whether a rationale identifies the specific mechanism its packet was
built to target, ranking the true \texttt{mechanism\_targeted} against the other
eight from the same market so that question, actors, and domain vocabulary are held
fixed. Rationales rank it first on $.278$--$.430$ of records against $1/9$ chance,
but the snippet alone scores $.516$, so most of that is echo of the text the model
just read. Removing every token the snippet supplied leaves $.128$--$.195$ under
matched query length, still above chance---but the generator's own rationale for
the packet, scored the same way, lands at $.119$, essentially chance, because it
paraphrases the mechanism rather than restating it. With the reference writer at
chance the metric has no usable dynamic range, and we draw no conclusion from it.

Models often declare orthogonal packets irrelevant: explicit non-relevance wording appears in
$4.9$--$61.9\%$ of orthogonal rationales against $0.1$--$2.8\%$ of directional
ones. What matters is whether saying it predicts doing it. Within the orthogonal
arm, records whose rationale declares irrelevance move less than those that do not,
and the market-clustered gap excludes zero for eight of the nine scale-valid
deployments: $0.88$~pp $[0.56,1.22]$ for Claude Opus~4.8, $1.70$
$[1.45,1.96]$ for DeepSeek~V4-Pro, $2.73$ $[2.28,3.19]$ for Qwen~2.5-72B, down to
$0.36$ $[0.21,0.54]$ for GPT-5.4, whose unflagged orthogonal movement is already
$1.3$~pp. The exception is Llama~3.1-8B at $0.26$ $[-0.64,1.06]$, where no coupling is
detectable; it is also the weakest deployment on sensitivity and directional
correctness, though we do not test whether the two are related. Qwen~2.5-7B is
omitted for mixed probability scales.

For eight of nine deployments, then, the revision a model makes is not independent
of the relevance it states, and the one deployment where that link is undetectable
is the weakest on every other measure. Two limits: the flag is a hand-built
lexicon, so it captures what a rationale declares rather than what the model
understood; and consistent self-description is compatible with surface routing, so
this does not show relevance was derived from the mechanism.

\textbf{One deployment was prompted differently.} GPT-5.4, and only GPT-5.4,
received the \texttt{novelty\_assessment} rule at Stage~3
(Section~\ref{app:update-protocol}).  Because that rule speaks directly to revision
size, its abstention rate is not comparable to the other deployments'. The
within-GPT contrast uses one template, but the extra instruction could interact
with packet arm, so its contrast should not be attributed to model differences
alone.

\textbf{Two remaining caveats.} A single generator wrote all packets and scored its
own work at a uniform 4/5 for plausibility, so automated validation establishes
structural completeness but not equal difficulty or independent content validity.
And the local models ran through Ollama at the default tag for each name, with the
served weight precision not recorded per run; rows ordered by parameter count also
vary in quantization and inference stack, so we do not attribute the open-weight
ordering to parameter count alone.

\begin{center}
\begin{minipage}{\linewidth}
\centering
\captionof{table}{Per-arm market coverage in the frozen June~17 report. Markets
counts entries in the model's per-market block, which for three Qwen rows
includes 20 carrying no scored update; the arm columns count markets with at
least one.}
\label{tab:exp1-arm-coverage}
\small
\begin{tabular}{lcccc}
\toprule
\textbf{Model} & \textbf{Markets} & \textbf{pro-$H_1$/anti-$H_1$} & \textbf{orthogonal} & \textbf{all three} \\
\midrule
Claude Opus~4.8 & 100 & 77 & 82 & 73 \\
GPT-5.4         & 100 & 99 & 99 & 99 \\
DeepSeek V4-Pro & 100 & 100 & 100 & 100 \\
\midrule
Qwen~2.5-72B    & 120 & 100 & 100 & 100 \\
Llama~3.3-70B   & 100 & 100 & 100 & 100 \\
Llama~3.1-70B   & 100 & 99 & 99 & 99 \\
Qwen~2.5-32B    & 120 & 100 & 100 & 100 \\
Qwen~2.5-14B    & 120 & 100 & 100 & 100 \\
Llama~3.1-8B    & 100 & 100 & 100 & 100 \\
\bottomrule
\end{tabular}
\end{minipage}
\end{center}

\begin{center}
\begin{minipage}{\linewidth}
\centering
\captionof{table}{Selectivity under statistics that do not divide by the orthogonal
mean. \emph{Sens.} is the frozen ratio; \emph{Diff.} is mean directional minus mean
orthogonal absolute revision; \emph{AUC} is the within-market probability that a
directional packet moves the forecast more than an orthogonal one, averaged over
markets; \emph{Markets} counts markets where mean directional movement exceeds
mean orthogonal movement; \emph{Length-m.} recomputes the ratio on the 66--114-word
band the arms share. Qwen~2.5-7B is omitted for mixed probability scales.}
\label{tab:exp1-selectivity-robustness}
\small
\input{tables/exp1_selectivity_robustness}
\end{minipage}
\end{center}

\subsubsection{Human materials review}
\label{app:exp1-human-review}
Nine adult reviewers submitted the 18-item materials protocol. It samples one
report from each of 18 distinct markets, balanced across six pro-YES, six
anti-YES, and six orthogonal packets. One reviewer met the disclosed
response-pattern exclusion below, leaving eight quality-eligible reviewers.
We treat the study as descriptive materials evidence. Before the synthetic
report, every item shows the market question,
resolution criteria, and matching summary from the frozen June~10 initial
forecast. The consent screen and every item instruct reviewers to treat
June~10, 2026 as the current date and use no external information or assistance.
Intended directions, source IDs, generator
metadata, and model responses remain hidden.

Reviewers record the report's conditional direction, confidence, clarity,
real-world plausibility, and usability as an explicit hypothetical premise.
We report exact agreement with the frozen direction and Cohen's $\kappa$
between the panel-majority labels and the frozen key. Reviewer-level Cohen's
$\kappa$, the component ratings, and every item-level judgment appear below.
Ambiguous judgments remain in every denominator.

Practice responses are stored separately and excluded from the analysis.

The analyzed sample contains eight quality-eligible reviewers (R1 and R3--R9),
exceeding the six-reviewer target for this descriptive materials check.
Table~\ref{tab:stage3-human} includes every eligible reviewer.

Exact-direction agreement is 108/144 (75.0\%). The panel-majority label matches
the frozen direction on 17 of 18 items (94.4\%; Cohen's $\kappa=.919$): all six
anti-$H_1$ items, all six orthogonal items, and five of six pro-$H_1$ items.
Across individual judgments, exact agreement is 35/48 for pro-$H_1$, 37/48 for
anti-$H_1$, and 36/48 for orthogonal packets. Reviewers marked the premise usable
on 124/144 judgments (86.1\%). The corresponding all-reviewer sensitivity is
reported below. These descriptive results indicate that readers usually recovered
the intended conditional direction and could use the reports as stipulated premises.
They also bound how legible that direction is: 75.0\% exact agreement and mean
pairwise Cohen's $\kappa=.47$ indicate substantial but imperfect recovery across
the four response labels. The panel is not a matched baseline for the agents---reviewers
emit a categorical label on 18 items, agents emit a revised probability against
their own prior on all 900---so we read it as evidence that sign had to be
inferred, not as a human--model comparison.
\begin{center}
\begin{minipage}{\linewidth}
\centering
\small
\captionof{table}{\textbf{Stage~3 human materials judgments for eight quality-eligible reviewers.}
Direction is exact agreement with the frozen conditional direction and
$\kappa$ is Cohen's $\kappa$ against that key over the four response options.
Confidence, clarity, and plausibility (C/C/P) are reviewer-level medians on 1--5
scales. The pooled row is descriptive because reviewers rated the same items; the
majority row takes the modal label per item, resolving ties to \emph{ambiguous}.}
\label{tab:stage3-human}
\begin{tabular}{lcccc}
\toprule
Reviewer & Direction agreement & $\kappa$ vs. key & Premise usable & C/C/P \\
\midrule
R1 & 12/18 (66.7\%) & .54 & 10/18 & 4/4/4 \\
R3 & 12/18 (66.7\%) & .51 & 13/18 & 4/4/3 \\
R4 & 18/18 (100.0\%) & 1.00 & 18/18 & 5/5/5 \\
R5 & 10/18 (55.6\%) & .38 & 18/18 & 4/4/4 \\
R6 & 13/18 (72.2\%) & .61 & 12/18 & 4/4/5 \\
R7 & 15/18 (83.3\%) & .76 & 17/18 & 5/4/3 \\
R8 & 12/18 (66.7\%) & .55 & 18/18 & 4/5/4 \\
R9 & 16/18 (88.9\%) & .84 & 18/18 & 5/5/5 \\
\midrule
Pooled & 108/144 (75.0\%) & --- & 124/144 & --- \\
Majority label & 17/18 (94.4\%) & .92 & --- & --- \\
\bottomrule
\end{tabular}
\end{minipage}
\end{center}

\begin{center}
\begin{minipage}{\linewidth}
\centering
\scriptsize
\setlength{\tabcolsep}{3pt}
\captionof{table}{\textbf{Item-level Stage~3 judgments for eight quality-eligible reviewers.}
Key is the frozen conditional direction; Majority is the modal panel response,
with ties assigned to \emph{ambiguous}. Agree and Usable give counts out of eight;
C/L/P gives median direction confidence, clarity, and plausibility on 1--5 scales.
Item numbers follow the frozen public packet order.}
\label{tab:stage3-items}
\input{tables/exp1_stage3_item_validation}
\end{minipage}
\end{center}


The response-pattern quality rule excludes any reviewer who selects one direction
on at least 17 of 18 balanced items. R2 met this rule by selecting \emph{more
likely} on 17 of 18 balanced items, including all six frozen negative items
and five of six frozen null items. Raw responses are retained, and no item is
removed or relabeled. Retaining R2 gives 115/162 (71.0\%) exact agreement across
all nine reviewers.


The task requests no identifying or sensitive information. It was conducted as
an internal materials review under the authors' institutional guidance; no IRB
protocol was opened.

Figure~\ref{fig:exp1-review-interface} records the reviewer-facing workflow. The
capture script at
\path{exp1_prospective/stage3_materials_annotation/capture_review_screenshots.py}
loads the frozen item packet against an empty temporary annotation database.
Item order was independently randomized by reviewer.


\begin{figure}[p]
\centering
\begin{subfigure}[t]{0.36\linewidth}
  \centering
  \includegraphics[width=\linewidth]{figures/exp1_stage3_review_consent.png}
  \caption{Information, temporal frame, and consent.}
\end{subfigure}\hfill
\begin{subfigure}[t]{0.60\linewidth}
  \centering
  \includegraphics[width=\linewidth]{figures/exp1_stage3_review_item.png}
  \caption{A complete item with June~10 context, synthetic report, and five judgments.}
\end{subfigure}
\caption{\textbf{Stage~3 materials-review interface.} Reviewers enter a private
code before seeing the information-and-consent screen. Each item repeats the
June~10 temporal frame, presents the question and pre-report background, labels
the fictional evidence as a synthetic report, and collects conditional
direction, confidence, clarity, plausibility, and premise usability. Item order
was independently randomized by reviewer. Intended direction and private packet
metadata are hidden.}
\label{fig:exp1-review-interface}
\end{figure}

Orthogonal packets are a useful outcome-relevance control, but are not equivalent
to receiving no new text. A literal no-information update condition would separate
generic prompt-induced revision from response to topically adjacent content.  The very
small hosted-system orthogonal revisions (0.5--1.4~pp) do not substitute for
that control.

\subsubsection{Prompt, retrieval, and runtime confounding}
Agents shared a substantive task contract, while execution was not byte-identical. Hosted
models used Azure Foundry search, while local models used Tavily results in an
Ollama workflow; temperatures and provider behavior also differ. Hosted--local
gaps therefore combine model capability with retrieval and runtime differences.

\subsubsection{Scale, missingness, and dependence}
Realized repeat counts are uneven because failures and top-ups differ by provider.
Equal-weight market aggregation prevents high-coverage markets from dominating
EHC, HFC, and ICS. Movement and sensitivity intervals use 20,000
market-clustered bootstrap draws that retain every packet and repeated initial-run
update from each selected market. The mixed-scale Qwen~2.5-7B row described in
Section~\ref{app:record-processing} remains an appendix compliance diagnostic and is
excluded from the main magnitude comparison.

\subsubsection{Premise acceptance at Stage~3}
Search is intentionally unavailable during updating, and agents are told to
accept the synthetic headline as true. This design measures conditional premise
use rather than source verification or epistemic refusal. The human materials review
uses the same instruction and therefore assesses comprehension and premise
usability. These results do not establish whether a model should trust or verify
the synthetic report.

\subsection{Reproducibility and artifact map}
\label{app:fig-code}

The paper treats the June~17 model-update report and the August~23 final
materials-review summary as the frozen numerical authorities. The relevant
artifacts are:

\begin{itemize}
  \item \textbf{Sampling:}
        \path{exp1_prospective/fetch_markets/select_markets.py} and the June~9
        selection/diversity JSON and manifests in
        \path{exp1_prospective/data/selected_markets/}.
  \item \textbf{Initial elicitation:}
        \path{exp1_prospective/agent/forecast_agent.py} for hosted runs and
        \path{exp1_prospective/agent/run_local_forecast.py} for Ollama runs, with
        records and manifests in \path{exp1_prospective/data/initial_forecasts/}.
  \item \textbf{Interventions:}
        \path{exp1_prospective/agent/build_counterfactuals.py} and the frozen
        \path{exp1_prospective/data/counterfactuals/counterfactuals_2026-06-10.jsonl}.
  \item \textbf{Updates:}
        hosted and local update runners under \path{exp1_prospective/agent/}, with
        update records in \path{exp1_prospective/data/updated_forecasts/}.
  \item \textbf{Evaluation:}
        \path{exp1_prospective/agent/evaluate_consistency.py} and
        \path{exp1_prospective/data/results/consistency_report_2026-06-17.json};
        \path{exp1_prospective/agent/evaluate_clustered_uncertainty.py}
        reconstructs the 20,000-draw market-clustered intervals in
        \path{exp1_prospective/data/results/clustered_movement_uncertainty.json};
        \path{exp1_prospective/agent/audit_orthogonal_abstention.py} writes the
        abstention/magnitude decomposition of Table~\ref{tab:exp1-abstention} to
        \path{exp1_prospective/data/results/orthogonal_abstention.json} and its
        table source, under the same resampling scheme.
  \item \textbf{Human materials review:}
        the frozen protocol, blinded packet, browser application, and tests in
        \path{exp1_prospective/stage3_materials_annotation/};
        \path{make_final_summary.py} reconstructs
        \path{data/exports/final_summary_current.json}, while
        \path{data/exports/agreement_current.json} is the paper-facing agreement
        export.
  \item \textbf{Figures:}
        \path{paper/generate_figures.py} writes the direction-selective panel;
        \path{paper/generate_paper_figures.py} writes the combined
        EHC/sensitivity panel.
\end{itemize}

The model-update tables are defined by the June~17 frozen JSON. The mutable collection
directories contain later records, so evaluating their present contents would
define a different dataset. Bit-for-bit reconstruction begins from the frozen
JSON; a full pipeline reconstruction uses only the manifests and records included
in that freeze.
%============================================================================
% END experiment1_appendix.tex
%============================================================================
%============================================================================
% BEGIN experiment1_context_reversal_appendix.tex
%============================================================================
\subsection{Exploratory context-reversal development pilot}
\label{app:exp1-context-reversal}

The original packet design leaves a relevance shortcut: a word-only classifier
separates directional from orthogonal packets with $.989$ accuracy. We therefore
ran a development pilot in which the news is identical within each family and
only its relationship to the forecast outcome changes. These are 20
assistant-authored fictional families, without independent human editing or
validation; the pilot develops a measurement instrument rather than providing
confirmatory evidence about real-world forecasting. The families span five reused
mechanism classes, not 20 distinct mechanisms.

\paragraph{Design and execution.}
Each family supplies a question, resolution rule, background, prior information,
and three relational contexts. Identical news is intended to favor the outcome,
oppose it, or have no bearing on it when the connecting relationship is broken.
A fourth condition omits the relationship description; it has no prespecified
sign or null expectation. Instructions ask models to condition on the supplied
fictional premises, with no web retrieval. Condition labels and private
annotations are excluded from prompts.

For each context, the model first gives a separate baseline probability. Three
independent branches then receive the same baseline and context with either
new news, a no-news message, or a verbatim repeat of information already present
at baseline. The repeated-information control is not a second presentation of
the genuinely new evidence. All branches use the same instructions about novelty,
and no branch sees another branch's answer. This separates contextual changes
from news-induced revision.

Qwen2.5-72B-Instruct, Llama-3.1-70B-Instruct, and Qwen3-32B ran from locally cached
checkpoints with three repeats, temperature $0.7$, top-$p=1$, and a 128-token
output limit; Qwen3 thinking was disabled. Syntax-constrained JSON required a
probability in $[0,1]$ without specifying its value or direction. Each model
produced 240 baselines and 720 updates: all 960 records were valid, for 2,880
records in total. Infrastructure retries replaced no valid responses.

\paragraph{Analysis and results.}
Directional accuracy uses all 120 planned positive/negative trials per model;
paired reversal requires both signs to be correct in each of 60 planned pairs.
Zero movement, invalid responses, and missing responses count as failures; there
is no minimum-movement filter. Broken-link movement uses 60 trials per model.
Intervals use 2,000 whole-family bootstrap draws, preserving contexts, branches,
and repeats. Thus 60 repeated pairs do not represent 60 independent families.

\begin{center}
\begin{minipage}{\linewidth}
\centering
\captionof{table}{Exploratory context reversal on independently unvalidated
materials. Brackets give descriptive 95\% family-bootstrap intervals. Sign
accuracy and paired reversal include every planned trial; broken-link movement
is the mean absolute change in percentage points.}
\label{tab:exp1-context-reversal}
\small
\input{tables/exp1_context_reversal}
\end{minipage}
\end{center}

Table~\ref{tab:exp1-context-reversal} shows only 30.0--36.7\% paired reversal.
The prespecified broken-link stability criterion requires the signed-mean 90\%
interval strictly within $\pm2$ percentage points and the upper endpoint of the
95\% interval for mean absolute movement below 2 points. All three models fail:
mean absolute movement is 10.5--13.3 points. All 240 no-news branches per model
exactly reproduce their priors, making raw and drift-adjusted results identical.
Mean absolute movement after repeated information, pooled across contexts, is
$.375$, $.021$, and $.521$ points for Qwen2.5-72B, Llama-3.1-70B, and Qwen3-32B,
respectively. Under this prompting protocol, ordinary re-elicitation drift does
not explain the much larger broken-link response to news.

\paragraph{Exploratory frontier extension.}
After inspecting the completed local pilot, we added GPT-5.6 Sol
(\texttt{gpt-5.6-sol}) on all 20 original families and all four contexts, using
only repeat zero. The 80-unit plan was frozen before its API responses were
collected; no family, label, or prompt was revised. The Responses API used low
reasoning effort, a 2,048-token output limit, strict JSON-schema output, no tools,
and no supplied temperature or seed. All 320 calls returned valid probabilities.
This is a post-local-results development extension, not a prespecified fourth
member of the original local roster. Local models used different decoding and
output constraints, so the comparison concerns these deployments rather than an
isolated effect of model identity.

\begin{center}
\begin{minipage}{\linewidth}
\centering
\captionof{table}{Matched single-repeat development comparison on the same 20
families. Each model contributes 40 directional trials, 20 reversal pairs, and
20 broken-link trials, with 320/320 valid records. Brackets are descriptive 95\%
family-bootstrap intervals; inference settings differ between deployments.}
\label{tab:exp1-context-reversal-frontier}
\small
\input{tables/exp1_context_reversal_frontier}
\end{minipage}
\end{center}

GPT-5.6 obtains 37/40 correct directions and 17/20 successful reversal pairs
(Table~\ref{tab:exp1-context-reversal-frontier}). Broken-link absolute movement
averages 0.5 points; its signed-mean 90\% interval is $[-1.5,0.0]$ and its upper
95\% absolute-movement bound is 1.5, meeting the original 2-point stability
criterion on this development set. Nineteen broken-link trials have zero change;
one changes by $-10$ points. All 80 no-news and all 80 repeated-news branches
exactly reproduce their own baselines. None of the 40 directional baselines
mechanically blocks the expected sign. The three failed directional updates
remain in the denominator, including the ambiguous reservoir family. These
results support testing relational relevance in a fresh validated cohort, while
the local failures show that performance is not uniform across deployments.

\paragraph{Exploratory direction-only companion.}
After inspecting the probability results, we elicited standalone direction
classifications from the same three local checkpoints, retaining all 20
families and four contexts at repeat zero. Each model separately classified 240
scenario--message combinations: new news, no news, and repeated information for
each context. No numerical prior was shown or forecast elicited. Constrained
JSON selected increase, decrease, unchanged, or unclear, with temperature $0.7$,
top-$p=1$, and 128 output tokens; Qwen3 thinking remained disabled. All 720
planned outputs were valid. Scored denominators retain every planned trial;
unclear, invalid, and missing outputs fail determinate targets. Masked new-news
judgments have no target and remain descriptive. The completed human materials
review covers the original 18 evidence packets; its documented scope is unchanged.

\begin{center}
\begin{minipage}{\linewidth}
\centering
\captionof{table}{Exploratory standalone direction judgments on the unchanged
development families. Entries are correct/planned counts. Paired reversal
requires both signed contexts to be correct. Broken-link news and the 160
no-news/repeated-news controls target unchanged. Each model has 240/240 valid
outputs; unclear is valid syntax but fails a determinate target.}
\label{tab:exp1-context-direction}
\small
\input{tables/exp1_context_direction}
\end{minipage}
\end{center}

Direction-only paired reversal is 6/20, 11/20, and 9/20 for Qwen2.5, Llama,
and Qwen3, respectively (Table~\ref{tab:exp1-context-direction}), compared with
5/20, 7/20, and 6/20 in their original repeat-zero probability responses.
For both Llama and Qwen3, four of five signed classification-only successes
occur where the original numerical prior mechanically blocked the required
update: a positive target starting at 1 or a negative target starting at 0.
On unblocked trials, Llama scores 27/34 categorically versus 26/34 numerically;
Qwen3 scores 24/36 versus 25/36. Qwen2.5 has no blocked endpoints and scores
25/40 in both tasks. The positive net gains for Llama and Qwen3 are therefore
concentrated in endpoint cases. Their separate elicitations and single repeat
do not isolate the cause of those differences. All three models continue to
miss many paired reversals and usually fail to classify broken-link news as
unchanged, despite largely correct no-news and repeated-information controls.
The original probability outcomes and interval-based stability criterion remain
the primary numeric analysis.

\paragraph{Exploratory Qwen3 reasoning follow-up.}
The same Qwen3-32B checkpoint ran with thinking enabled and disabled on all
20 families and four contexts at repeat zero. Each arm elicited 80
context-specific priors and three separately generated update branches per
prior, for 320 planned responses. Updates received only their own arm's prior.
Both modes used unconstrained decoding, identical request seeds, temperature
$0.7$, top-$p=1$, a 4,096-token output limit, and an 8,192-token context limit.
The original constrained, 128-token non-thinking deployment remains a separate
reference. After observing non-thinking outputs, a recorded parsing amendment
permitted exactly one optional Markdown fence around otherwise strict final
probability JSON. This post-output rule applied symmetrically to both modes;
original format compliance is separately reported. Every received raw
answer was preserved and reparsed without regeneration. Format recovery
generated only previously unattempted update branches whose baselines became
parseable. Truncation, unfinished reasoning, and other remaining failures were
retained; the parser never extracted a probability from unfinished reasoning.
All planned trials remain in scored denominators, with the original sign,
paired-reversal, and broken-link stability criteria.

\begin{center}
\begin{minipage}{\linewidth}
\centering
\captionof{table}{Exploratory Qwen3 reasoning follow-up on the same 20 families.
Correct-sign and paired-reversal entries retain all planned trials. The two
unconstrained arms use the symmetric format amendment; the original constrained
arm is a descriptive reference. Broken-link mean absolute movement is observed
only, with its missing trial shown explicitly.}
\label{tab:exp1-context-reasoning}
\small
\resizebox{\linewidth}{!}{\input{tables/exp1_context_reasoning}}
\end{minipage}
\end{center}

All 320 records per arm are present. The disabled arm has 320 valid probabilities
after recovery: 103 previously rejected fenced answers were reparsed, and 198
previously blocked update branches were generated for the first time. Before
that amendment, its strict-format record counts were 19 valid, 103 parse
failures, and 198 blocked branches. The enabled arm retains 300 valid records,
five truncated baselines, and their 15 blocked updates; format recovery changes
none of these outcomes. Four truncated baselines belong to the radio family
(\texttt{dev\_16}), and one to the positive carton-transfer context
(\texttt{dev\_06}). These failures remain in all planned binary denominators.

Thinking enabled obtains 30/40 correct signs and 12/20 paired reversals, compared
with 25/40 and 6/20 for matched unconstrained thinking disabled
(Table~\ref{tab:exp1-context-reasoning}). Using 2,000 paired whole-family
bootstrap draws, the enabled-minus-disabled contrast is $+12.5$ percentage
points for signs (descriptive 95\% interval $[-2.5,27.5]$) and $+30.0$ points
for paired reversal ($[5.0,55.0]$). Each paired binary contrast retains all
20 families, including failures. Direction-blocking baseline endpoints occur
in 5/40 valid disabled baselines and 1/37 valid enabled baselines; three enabled
directional baselines are missing.

Broken-link absolute movement averages 14.05 points for disabled thinking
(20/20 observed) and 7.63 for enabled thinking (19/20 observed). On their
19 common complete pairs, the respective means are 14.74 and 7.63 points,
giving an enabled-minus-disabled contrast of $-7.11$ points
($[-16.32,1.32]$). The missing radio-family measurement is never replaced
by zero. The original stability criterion is not met by disabled thinking
and is indeterminate for enabled thinking because a planned measurement is
missing. Enabled outputs contain nonempty reasoning in all 305 received
completions, with 300 closing the reasoning segment and five exhausting the
token limit. Mean output length is 966.4 tokens per received enabled completion
versus 12.5 for disabled thinking. These results describe the matched deployment
change on this development sample; uncertainty, changed baseline judgments,
and remaining failures limit conclusions about the mechanism of improvement.

\paragraph{Controlled robustness development on the original families.}
A separately versioned development run retains all 20 original parent families,
including previous successes and failures. Each family has a repaired base,
consistent entity renaming, a paraphrase, and a resample with identical visible
text but different request IDs and seeds. Explicit transfer destinations repair
underspecified routes; signed and broken contexts share the same inventory of
positive, negative, and side-channel roles, with actor assignments changed and
row order counterbalanced. These revisions do not replace the original results.
Repair changes the supplied premises and can change elicited priors, so an
old-versus-repaired contrast is not a causal estimate of clarification.

All four versions retain positive, negative, broken, and masked contexts.
Each numerical context receives its own baseline and three independent branches:
new news, no news, and information already present at baseline. The separate
categorical task presents no numerical prior. The three original local
checkpoints use constrained JSON, temperature $0.7$, top-$p=1$, and a 128-token
output limit; Qwen3 thinking is disabled. Each model has all 1,280 numerical and
960 categorical records. Valid outputs are Qwen2.5-72B: 1280/1,280 numeric and 960/960 categorical; Llama-3.1-70B: 1280/1,280 numeric and 960/960 categorical; Qwen3-32B: 1280/1,280 numeric and 960/960 categorical.
Invalid responses remain scored failures, and magnitude estimates use observed
measurements with missingness shown explicitly. The independent audit verifies
execution, provenance, and scoring; it is not independent human validation of
these revised materials.

Table~\ref{tab:exp1-robustness-v2-performance} reports every model and version.
Zero numerical movement fails a signed target. Paired reversal requires both
signed contexts to be correct; categorical unclear fails a determinate target.
All intervals are descriptive, using 2,000 whole-parent-family bootstrap draws
(seed 20260921), carrying contexts, branches, and versions together. There are
20 sampling families; the four versions do not create 80 independent families.

\begin{center}
\begin{minipage}{\linewidth}
\centering
\captionof{table}{Controlled robustness development. Binary entries are
correct/planned counts: 40 signed trials, 20 reversal pairs, and 20 broken-link
judgments per version. Broken movement is mean absolute numerical revision in
percentage points, with its descriptive 95\% family-bootstrap interval.}
\label{tab:exp1-robustness-v2-performance}
\small
\resizebox{\linewidth}{!}{\input{tables/exp1_context_robustness_v2_performance}}
\end{minipage}
\end{center}

For Qwen2.5-72B, Llama-3.1-70B, and Qwen3-32B, respectively,
repaired-base numerical paired reversal is 8/20, 8/20, 7/20, and broken-link
mean absolute movement is 15.50, 6.46, 10.00 points. The corresponding
interval-based stability statuses are not met, not met, not met.

The numerical broken-link stability diagnostic retains the original criterion:
the signed-mean 90\% interval must lie strictly within $\pm2$ points and the
upper 95\% bound on mean absolute movement must be below 2 points, with all
planned measurements observed. Detailed statuses are retained in the source
summaries; a small point estimate alone does not establish stability.

\paragraph{Editing and stochastic variation.}
Let $\Delta p_v$ be a version's new-news update relative to its own context's
baseline. Table~\ref{tab:exp1-robustness-v2-edits} separates mean absolute
baseline discrepancy $D_0$ from mean absolute update discrepancy
$D_\Delta=|\Delta p_v-\Delta p_{\mathrm{base}}|$, each averaged across the
80 planned family--context pairs. The within-2-point count describes individual paired observations.
The excess column averages
$E_\Delta=|\Delta p_v-\Delta p_{\mathrm{base}}|
-|\Delta p_{\mathrm{resample}}-\Delta p_{\mathrm{base}}|$
on common observed measurements. The same-text resample is a stochastic
reference at the same temperature; its excess is zero by construction.
Renaming and paraphrase discrepancies must therefore be read alongside this
reference. A single resample does not isolate a causal effect of editing,
and an interval including zero does not establish equivalence or invariance.

\begin{center}
\begin{minipage}{\linewidth}
\centering
\captionof{table}{Each edited version or same-text resample compared with repaired
base. Numerical discrepancies are in percentage points; brackets give 95\%
parent-family-bootstrap intervals for excess over resampling. Agreement uses
80 planned categorical pairs, including masked new news. Joint correctness uses
60 determinate pairs, excluding masked new news. Stable wrong answers count as
agreement but fail joint correctness. Missing/invalid binary pairs remain failures;
magnitudes are observed only.}
\label{tab:exp1-robustness-v2-edits}
\small
\resizebox{\linewidth}{!}{\input{tables/exp1_context_robustness_v2_edits}}
\end{minipage}
\end{center}

\paragraph{Controls and lexical checks.}
No-news and repeated-information controls remain separate
(Table~\ref{tab:exp1-robustness-v2-controls}); the repeated information was already
present at baseline and is not a second presentation of genuinely new evidence.
Each control has 80 planned contexts per model and version. Their performance
does not substitute for directional or broken-link performance.

\begin{center}
\begin{minipage}{\linewidth}
\centering
\captionof{table}{Separate no-news and repeated-information controls. Within-2-point
and categorical unchanged entries use all 80 planned contexts. Mean absolute
numerical movement is in percentage points; any missing magnitude measurements
are shown explicitly. Full descriptive family-bootstrap intervals are retained
in the source summaries.}
\label{tab:exp1-robustness-v2-controls}
\small
\resizebox{\linewidth}{!}{\input{tables/exp1_context_robustness_v2_controls}}
\end{minipage}
\end{center}

Whole-parent-family-held-out news-unigram, full-text unigram, full-text bigram,
and character-ngram baselines obtain directional accuracies of 33.3\%, 33.3\%, 33.3\%, 32.5\%,
and balanced relevance accuracies of 50.0\%, 50.0\%, 50.0\%, 47.8\%, respectively. The chance
references are $1/3$ and $1/2$. These measurements concern the tested lexical
baselines; they do not establish that all shortcuts are removed. The revised
materials and all three checkpoints remain development evidence. No independent
fresh-cohort result is asserted here.

\paragraph{Interpretive limits of the original pilot.}
In the original three-repeat pilot, baseline probabilities at the relevant endpoint mechanically prevent a correct
signed update in 0/120 Qwen2.5-72B, 19/120 Llama, and 12/120 Qwen3 trials; these
remain failures. The original pilot materials retain ambiguities: the carton-transfer family
(\texttt{dev\_06}) does not specify where cartons leaving the buffer go, and the
reservoir family (\texttt{dev\_07}) does not explicitly say that outgoing water
bypasses the turbines. The resin-conveyor (\texttt{dev\_10}) and silt-pump
(\texttt{dev\_18}) families likewise leave outlet destinations unspecified. No family was excluded or relabeled after observing
responses. The development results do not establish reliable relational relevance
judgments on independently validated or real-world materials. Any material revision requires a new development version; a fresh,
independently validated and frozen cohort is needed for confirmation.
%============================================================================
% END experiment1_context_reversal_appendix.tex
%============================================================================

% Insertion point for the fresh finite-rule context-reversal appendix. The
% guarded publication pipeline in
% exp1_prospective/context_reversal/fresh/publish.py stages a copy of this file
% and inserts an \input for that appendix directly below this marker. Keep the
% marker line verbatim and on a line of its own.
%%% FRESH-CONTEXT-APPENDIX %%%

%============================================================================
% BEGIN experiment2_appendix.tex
%============================================================================
% ============================================================
% EXPERIMENT 2 -- COIN CITY FULL DETAILS
% ============================================================

\section{Experiment 2: contextual relationship selection and revision in Coin City}
\label{app:coin-city}

Experiment~2 tests whether context
selects which response relationship a model applies and whether direct target
observations revise that selection. The six primary deployments and symbol-control
comparisons reuse the same 250 seeded episodes, so prompt-arm differences are
paired rather than comparisons between independently sampled tests.

\subsection{Design and data-generating process}
\label{app:coin-city-design}

\subsubsection{Design summary}
The design identifier is
\path{coin_city_stable_relationship_claude_n250_v4}. It contains 250 independent
episodes generated from seeds 930000--930249, five target-city prefixes
\(k\in\{0,1,2,3,4\}\), and three semantic-selection arms. Cue-inversion and
arbitrary-symbol modules reuse those episode identities and numerical inputs; each
has a hashed task manifest (Section~\ref{app:coin-city-execution}). Every model received the same
episode records, queries, and arm-specific prompt text.

\subsubsection{Balanced latent regimes}
Each episode crosses two binary factors: whether City A or City B is the strongly
responsive reference city, and whether target City C is strongly or weakly
responsive. The closest-balanced allocation contains 63 A-strong/C-strong,
62 A-strong/C-weak, 62 B-strong/C-strong, and 63 B-strong/C-weak episodes.
Within an episode, each city receives an independently sampled stable slope
\[
  g_{ec}\sim
  \begin{cases}
    \mathcal N(0.90,0.03^2), & \text{strong regime},\\
    \mathcal N(0.25,0.03^2), & \text{weak regime}.
  \end{cases}
\]
Thus the cue identifies a response regime but not the realized coefficient.

\subsubsection{Observed cases and forecast query}
Every city has four reference cases arranged as two matched pairs. Rows in a pair
share a starting poll and news magnitude but have opposite signs; the first sign is
random. For reference row \(r\),
\[
  p^{\mathrm{end}}_{ecr}
    = \operatorname{round}_{0.1}
      \left(p^{\mathrm{start}}_{ecr}+g_{ec}x_{ecr}
      +\epsilon_{ecr}\right),
  \qquad \epsilon_{ecr}\sim\mathcal N(0,5^2).
\]
Starting polls are uniform on \([35,65]\), and reference-news magnitudes are drawn
from \(\{6,7,8,9,10\}\). The query shock is drawn from
\(\{-10,\ldots,-4,4,\ldots,10\}\). Its scoring target is the latent conditional
expectation \(p_0+g_Cx\), with no additional query-time observation noise.

The prefix \(k\) exposes the first \(k\) City C rows in their stored order.
Consequently, \(k=2\) completes the first matched pair and \(k=4\) completes both;
\(k=1\) and \(k=3\) deliberately expose partial pairs. City A and B always
contribute all four rows in the ABC arms.

\begin{center}
\begin{minipage}{\linewidth}
\centering
\captionof{table}{Coin City design and evaluation dimensions. Final analytic
record counts retain one record per task cell; raw transport-level calls are
discussed in Section~\ref{app:coin-city-limitations}.}
\label{tab:coin-city-design}
\small
\begin{tabular}{p{0.31\linewidth} p{0.62\linewidth}}
\toprule
Quantity & Value \\
\midrule
Independent episodes & 250 \\
Factorial cells & 63, 62, 62, 63 across the four A/B-by-C regimes \\
Target-city prefixes & \(k\in\{0,1,2,3,4\}\) \\
Primary prompt arms & target only; ABC/no C context; ABC/correct C context \\
Intervention modules & ABC/misleading C context; ABC/episode-randomized arbitrary labels \\
Primary deployment roster & Claude Opus~4.8; GPT-5.6; DeepSeek-V4-Pro; Kimi K3; Gemini 3.6 Flash; Claude Opus~5 \\
Symbol-comparison roster & Qwen3 4B/8B/14B/32B; Qwen2.5 7B/14B/32B/72B;
Llama~3.1 8B/70B; DeepSeek-V4-Pro; Kimi K3; Claude Opus~4.8/5; GPT-5.6 Sol \\
Primary task cells / final records & 22,500 / 22,500 (22,497 parseable) \\
Misleading task cells / final records & 7,500 / 7,500 (all parseable) \\
Symbol-comparison task cells / response records & 56,250 / 56,250 (15 models) \\
Strong / weak slope means & 0.90 / 0.25 poll points per news point \\
Slope standard deviation & 0.03 within each regime \\
Reference-case noise & 5 poll points \\
Reference cases per city & 4 (two matched pairs) \\
\bottomrule
\end{tabular}
\end{minipage}
\end{center}

\subsubsection{Preflight validation and task integrity}
The GPT-5.6, DeepSeek-V4-Pro, Kimi K3, Gemini 3.6 Flash, and Claude Opus~5
replications reuse the exact task files and prompt hashes created for the original
Claude Opus~4.8 evaluation; no episode was regenerated. Before model calls, the
manifest hashed the engine, task files, answer and scoring keys, validator, runner,
renderer, and launcher. Automated validation confirmed the intended dimensions and
factorial balance, prompt contracts and hashes, correct cue-to-regime mapping,
visible matched pairs, and benchmark reconstruction. No model response or
performance statistic was used by this preflight validator.

\subsection{Prompt arms and information available}
\label{app:coin-city-arms}

\subsubsection{Shared task instructions}
All prompts state that rows are separate polling cases rather than a time series,
positive news favors the candidate, poll change is end minus start, and
within-city responsiveness is stable even though individual polls are noisy. Rows
sharing a pair label are described as having the same starting poll and
equal-magnitude, opposite-sign news. Every displayed row includes the starting
poll, news value, change, and ending poll.

The required answer is a JSON object with exactly a one-sentence
\texttt{rationale} and a numeric \texttt{predicted\_poll} between 0 and 100. The
prompt asks for no step-by-step derivation.

\subsubsection{Five information conditions}
\begin{description}
  \item[Target only.]
  The model sees the first \(k\) City C cases and is explicitly told that no
  additional background information is available. This arm measures recovery
  from target-city evidence alone.

  \item[ABC, no City C context.]
  The model also sees four cases from each reference city, each introduced by a
  background sentence. The strongly responsive reference city is described
  verbatim as receiving ``national news coverage, where campaign stories receive
  sustained attention and tend to produce larger immediate poll movements'', and
  the weakly responsive city as receiving ``local news coverage, where campaign
  stories compete with local issues and tend to produce smaller immediate poll
  movements''. The direction of the coverage-to-response mapping is therefore
  stated in words, though no numerical value is. City C receives no background
  description. This arm tests numeric transfer without a cue that aligns C to
  either reference regime.

  \item[ABC plus correct City C context.]
  The numeric rows and query are identical to the preceding arm. The sole added
  information is that City C receives national or local coverage, correctly
  aligned with its latent strong or weak regime. This is the primary cue-transfer
  condition.

  \item[ABC plus misleading City C context.]
  This intervention arm is identical to the correct-context prompt except that its
  City C sentence names the opposite coverage regime. Every number, case row,
  query, and other instruction is unchanged. This arm tests whether the model
  follows the semantic cue when it conflicts with the truth.

  \item[ABC plus an arbitrary City C label.]
  Every national/local description is replaced by one of two arbitrary labels,
  \texttt{KIV} or \texttt{ZOR}. The prompt states only that cities sharing a label
  share a typical response regime and that the labels have no inherent meaning or
  ordering. Their higher/lower-response mapping reverses on alternating episodes;
  City A and B's numerical cases are therefore the only evidence from which the
  episode-local mapping can be induced before applying City C's label.
\end{description}

No prompt contains a fitted coefficient, an estimator, the strong or weak
numerical means, the case-noise scale, a rule for pooling cities, or any analyst
benchmark. The correct-versus-no-context contrast therefore isolates the one City
C background sentence while holding all displayed numeric evidence fixed; the
correct-versus-misleading comparison isolates the sentence's regime assignment.
The symbol-versus-no-context contrast asks whether the same assignment can be
induced without lexical knowledge of national and local coverage.

Because the reference sentences state the direction of the mapping, the semantic
arms are not by themselves a test of whether that direction can be induced. A
policy that reads no displayed number and assigns slope $1.0$ to the
nationally covered regime and $0.2$ to the locally covered one attains $k=0$
forecast MAE $0.523$ and implied-slope correlation $.996$, better than any
deployment and better than the analogical-regression benchmark of
Table~\ref{tab:coin-city-analogical}; $42$ of $108$ prior pairs on the grid
$\{0.4,\dots,1.5\}\times\{0.0,\dots,0.8\}$ beat every deployment's $k=0$ MAE.
The semantic arms therefore measure how a stated regime assignment is converted
into a numerical response. Two separate controls carry the induction claim: the
arbitrary-symbol arm, whose labels have no stated direction and whose mapping
reverses across episodes, and the reference-selection analysis of
Section~\ref{app:coin-city-reference-selection}, which shows that the forecast
tracks the named reference city's own sampling error rather than any fixed prior.

\subsection{Model execution and response processing}
\label{app:coin-city-execution}

\subsubsection{Deployment configurations}
Table~\ref{tab:coin-city-execution} reports the effective API settings used for
production. Temperature refers to the value sent to the provider: GPT-5.6 and
Claude Opus~5 omit it even though the local GPT run manifest records a default of
zero. No deployment used tools or retrieval.

\begin{center}
\begin{minipage}{\linewidth}
\centering
\captionof{table}{Production deployment configuration and primary-arm parse counts.}
\label{tab:coin-city-execution}
\scriptsize
\setlength{\tabcolsep}{3.5pt}
\begin{tabular}{p{0.16\linewidth} p{0.21\linewidth} p{0.27\linewidth} @{\hspace{1.5em}} r r}
\toprule
Deployment & Interface & Reasoning / sampling & Output ceiling & Parsed \\
\midrule
Claude Opus~4.8 & Anthropic Messages via Azure & temperature 0 & 512 & 3,750 \\
GPT-5.6 & OpenAI Responses & low reasoning; temperature omitted & 4,096 & 3,749 \\
DeepSeek-V4-Pro & OpenAI Chat Completions & temperature 0 & 4,096 & 3,750 \\
Kimi K3 & OpenAI Chat Completions & low reasoning; temperature 0 & 16,384 & 3,750 \\
Gemini 3.6 Flash & Gemini OpenAI-compatible Chat & low reasoning; temperature 0 & 16,384 & 3,750 \\
Claude Opus~5 & Foundry Anthropic Messages & default extended thinking; temperature omitted & 4,096 & 3,748 \\
\bottomrule
\end{tabular}
\end{minipage}
\end{center}

Claude Opus~5 returns extended thinking before its answer; the runner therefore
selects its first text block. Gemini is the sole fast-tier deployment, so the six
systems should be read as a replication set rather than a tier-matched leaderboard.
Authentication material was supplied at runtime and was not written to source,
manifests, logs, or response files.

\subsubsection{Deployment attempts, retries, and pilots}
Every primary model--prompt pair received one application-level attempt. A failed
or answerless task was not re-asked. HTTP-client retries were disabled for GPT,
DeepSeek, Kimi, and Opus~5. The original Opus~4.8 client allowed up to two
transport retries. Gemini's runner allowed adaptive pre-model transport retries
but recorded none in the completed shards. These transport policies do not alter
the one-answer-per-task analytic rule.

Before production, 119 Kimi, 15 Gemini, and 17 Opus~5 task-shaped calls were used
only to select configurations that reliably returned the required JSON schema.
Forecast accuracy was not inspected. The calls are archived separately and are
excluded from every denominator and analysis; each production namespace began
empty. Production tasks were sharded only for throughput.

\subsubsection{Response parsing and record selection}
The parser scans the raw reply for JSON objects, retains the last object with an
exact \texttt{predicted\_poll} key whose value converts to a number in
\([0,100]\), truncates the retained rationale to 500 characters, and stores at
most 8,000 raw characters. When duplicate task records exist, the analysis keeps
the latest parseable response, allowing a valid response to supersede an earlier
error. A single offline normalization removed whitespace from JSON keys in one
otherwise valid Opus~4.8 reply; this recovered one primary record without another
model call. Other record-level corrections are documented in
Section~\ref{app:coin-city-limitations}.

\subsubsection{Misleading-context design}
The cue-inversion module begins from the frozen semantic-selection tasks. Before
any inversion call, its builder verifies that every prompt equals its
correct-context counterpart under
exactly one sentence substitution and that all numeric content is unchanged. The
cue is inverted for all 1,250 episode-prefix tasks: 625 strong-city tasks are
described as weak and 625 weak-city tasks as strong. Each deployment uses the
same client and production configuration as its correct-context arm.

\subsubsection{Arbitrary-symbol design}
The symbol builder generates 1,250 matched prompts from the same frozen
episodes and verifies that their complete sequence of numeric tokens is identical
to the correct-context arm. The higher-response label is \texttt{KIV} in 125
episodes and \texttt{ZOR} in 125; the target regime-by-label cells contain 62--63
episodes each. A prompt validator rejects national/local and strong/weak semantic
descriptions. At \(k=0\), City C supplies no numerical cases, so its correct regime
can be recovered only by decoding the label from City A and B's examples.

For DeepSeek-V4-Pro, we collected all 1,250 symbol tasks and contemporaneously
re-collected all 2,500 no-context and semantic-context tasks on the same reachable
Azure resource, using one deployment identifier and configuration. The released
15-model comparison combines that triplet with Qwen3-4B, Qwen3-8B,
Qwen3-14B, Qwen3-32B, Qwen2.5-7B, Qwen2.5-14B, Qwen2.5-32B, Qwen2.5-72B,
Llama~3.1-8B, Llama~3.1-70B, Kimi K3, Claude Opus~4.8, Claude Opus~5, and
GPT-5.6 Sol.

The eight Qwen and two Llama checkpoints received all 3,750 prompts through vLLM
in bfloat16, with temperature-zero decoding, a 3,072-token model context, and a
512-token output ceiling. Qwen uses the no-thinking chat template; Llama uses its
native instruction template. Qwen2.5-14B uses two-way and Qwen3-32B and
Qwen2.5-32B four-way tensor parallelism, while Qwen2.5-72B and Llama~3.1-70B use
eight-way tensor parallelism.

For Kimi K3, Claude Opus~5, and GPT-5.6 Sol, the comparison reuses each
deployment's preserved no-context and semantic-context responses and adds only
the 1,250 symbol prompts under the same model-specific production protocol.
Claude Opus~4.8's symbol arm instead ran on a second Azure resource serving the
same deployment and used a 1,024-token rather than 512-token ceiling; all 1,250
responses returned valid JSON. We therefore treat its arm contrast as a
cross-run robustness result and do not use it to rank providers.

Every runner verifies the frozen task hash before generation and writes a
timestamped manifest. Each released model has 1,250 unique response records in
every arm; this count does not imply that every prediction parsed. No primary
response file or frozen task is overwritten. Original outputs are preserved, and
the final tables use matched complete cases after the no-call arithmetic reparse
described in Section~\ref{app:coin-city-limitations}.

\subsection{Estimands, benchmarks, and statistical inference}
\label{app:coin-city-analysis}

\subsubsection{Forecast and latent-response estimands}
Forecast error is measured against the latent conditional expectation rather than
a newly sampled noisy ending poll:
\[
  \operatorname{MAE}
  =\frac{1}{n}\sum_{e=1}^{n}
    \left|\widehat p_e-(p_{0e}+g_{Ce}x_e)\right|.
\]
For every nonzero query shock, a forecast also identifies an implied slope
\(\widehat g_e=(\widehat p_e-p_{0e})/x_e\). We report its mean absolute error from
\(g_{Ce}\) and its Pearson correlation with \(g_{Ce}\) across episodes.
Correlation is needed because a nearly constant prior-mean slope can attain
moderate MAE while failing to distinguish strong from weak cities.

Within each deployment and prefix, the three primary-arm metrics use the common
intersection of episodes with parseable responses in all three arms. Arm
contrasts are paired by episode and prefix. The misleading-arm table uses its
final unique parseable record for every task. The arbitrary-symbol analysis uses
each model's corresponding common intersection across its three matched arms.

\subsubsection{Analyst-only benchmarks}
The ABC OLS benchmark fits poll change on news through the origin using every
displayed A, B, and available C row. The target-only OLS benchmark uses only
displayed City C rows and is undefined at \(k=0\). Both are reconstructed from
the visible data by the analyst; neither the estimator nor its output appears in
a model prompt.

The label-conditioned analogical regression implements the minimal structured
inductive policy defined in Methods. It
fits separate through-origin slopes to City A and City B, uses City C's coverage
label to select the reference assigned to the same response regime, and predicts
with that slope. At $k>0$, it pools the selected reference rows with the visible
City C rows before refitting. In the misleading arm, the inverted label selects
the opposite reference. The implementation uses only displayed rows; neither the
estimator nor its output appears in a model prompt.
\begin{center}
\begin{minipage}{\linewidth}
\centering
\captionof{table}{Zero-shot behavior of the analogical policy and the six model
deployments. At $k=0$, City C has no numerical cases. Model columns give the
range across deployments; forecast MAE is in poll points.}
\label{tab:coin-city-analogical}
\small
\begin{tabular}{l r r c c}
\toprule
Cue & Policy MAE & Policy $\rho$ & Model MAE range & Model $\rho$ range \\
\midrule
Correct & 1.842 & .656 & 1.596--2.474 & [.52,\,.70] \\
Inverted & 4.561 & $-.667$ & 4.204--4.361 & [$-.68$,\,$-.50$] \\
\bottomrule
\end{tabular}
\end{minipage}
\end{center}


\subsubsection{Uncertainty and comparison families}
Primary arm contrasts use percentile intervals from 2,000 bootstrap resamples of
the paired episode-level absolute-error difference. Deterministic, model-specific
seeds make these intervals reproducible. The \(k=0\) pairwise model table uses
4,000 paired bootstrap resamples with seed 20260812 and the common episode
intersection for each model pair. Negative differences favor the first-listed
condition or model.

The episodes are the resampling units. Intervals are pointwise descriptions of
this fixed seeded synthetic evaluation: they are not simultaneous bands across
prefixes, and no multiple-comparison correction is applied. Accordingly, the
number of pairwise intervals excluding zero is not used as an estimand; each
interval supports only its displayed cell.

\subsection{Response-inference diagnostics}
\label{app:coin-city-results}

The tables below retain a common arm and prefix ordering. Within a deployment and
prefix, all three primary arms use the same complete-case episode set; \(n\) is
repeated so that every exception is visible. Boldface identifies the lowest
forecast MAE or slope MAE, and the highest slope correlation, within that
deployment and prefix.

\subsubsection{Per-deployment primary-arm recovery}

\begingroup
\setlength{\topsep}{3pt}
\setlength{\partopsep}{0pt}
\setlength{\parsep}{0pt}
\setlength{\itemsep}{0pt}
\begin{center}
\begin{minipage}{\linewidth}
\centering
\captionof{table}{Claude Opus~4.8 forecast and latent-response diagnostics for every prompt
arm. Forecast MAE is in poll points; slope MAE and $\rho$ compare the implied
response with the true City C slope.}
\label{tab:coin-city-recovery}
\small
\begin{tabular}{c l r r r r}
\toprule
$k$ & Arm & $n$ & Forecast MAE & Slope MAE & $\rho(\hat g,g_C)$ \\
\midrule
0 & Target only & 250 & 2.321 & .332 & .019 \\
  & ABC, no C context & 250 & 2.651 & .379 & -.063 \\
  & ABC + C context & 250 & \textbf{1.817} & \textbf{.265} & \textbf{.700} \\
\midrule
1 & Target only & 250 & 3.730 & .549 & .371 \\
  & ABC, no C context & 250 & 2.438 & .358 & .322 \\
  & ABC + C context & 250 & \textbf{2.147} & \textbf{.316} & \textbf{.584} \\
\midrule
2 & Target only & 250 & 2.668 & .403 & .383 \\
  & ABC, no C context & 250 & \textbf{2.041} & \textbf{.296} & \textbf{.492} \\
  & ABC + C context & 250 & 2.079 & .308 & .491 \\
\midrule
3 & Target only & 250 & 2.485 & .358 & .449 \\
  & ABC, no C context & 250 & 2.128 & .305 & .513 \\
  & ABC + C context & 250 & \textbf{1.977} & \textbf{.291} & \textbf{.627} \\
\midrule
4 & Target only & 250 & 2.209 & .324 & .498 \\
  & ABC, no C context & 250 & 1.929 & .281 & .542 \\
  & ABC + C context & 250 & \textbf{1.760} & \textbf{.257} & \textbf{.658} \\
\bottomrule
\end{tabular}
\end{minipage}
\end{center}

\begin{center}
\begin{minipage}{\linewidth}
\centering
\captionof{table}{GPT-5.6 forecast and latent-response diagnostics on the same frozen
prompts. Metrics use episodes parseable in all three arms at each prefix.}
\label{tab:coin-city-recovery-gpt56}
\small
\begin{tabular}{c l r r r r}
\toprule
$k$ & Arm & $n$ & Forecast MAE & Slope MAE & $\rho(\hat g,g_C)$ \\
\midrule
0 & Target only & 250 & 3.596 & .511 & .054 \\
  & ABC, no C context & 250 & 2.518 & .352 & -.041 \\
  & ABC + C context & 250 & \textbf{1.596} & \textbf{.224} & \textbf{.678} \\
\midrule
1 & Target only & 249 & 2.993 & .439 & .394 \\
  & ABC, no C context & 249 & 2.768 & .399 & .302 \\
  & ABC + C context & 249 & \textbf{1.586} & \textbf{.225} & \textbf{.669} \\
\midrule
2 & Target only & 250 & 2.560 & .387 & .459 \\
  & ABC, no C context & 250 & 2.050 & .299 & .534 \\
  & ABC + C context & 250 & \textbf{1.854} & \textbf{.268} & \textbf{.606} \\
\midrule
3 & Target only & 250 & 2.240 & .328 & .528 \\
  & ABC, no C context & 250 & 1.967 & .286 & .584 \\
  & ABC + C context & 250 & \textbf{1.773} & \textbf{.256} & \textbf{.600} \\
\midrule
4 & Target only & 250 & 2.153 & .318 & .577 \\
  & ABC, no C context & 250 & 1.746 & .253 & .653 \\
  & ABC + C context & 250 & \textbf{1.594} & \textbf{.228} & \textbf{.664} \\
\bottomrule
\end{tabular}
\end{minipage}
\end{center}

\begin{center}
\begin{minipage}{\linewidth}
\centering
\captionof{table}{DeepSeek-V4-Pro forecast and latent-response diagnostics on the same
frozen prompts. All 250 episodes parsed in every arm and prefix.}
\label{tab:coin-city-recovery-deepseek}
\small
\begin{tabular}{c l r r r r}
\toprule
$k$ & Arm & $n$ & Forecast MAE & Slope MAE & $\rho(\hat g,g_C)$ \\
\midrule
0 & Target only & 250 & \textbf{2.447} & .353 & -.027 \\
  & ABC, no C context & 250 & 2.841 & .401 & -.027 \\
  & ABC + C context & 250 & 2.474 & \textbf{.344} & \textbf{.523} \\
\midrule
1 & Target only & 250 & 3.585 & .525 & .334 \\
  & ABC, no C context & 250 & 3.130 & .458 & .302 \\
  & ABC + C context & 250 & \textbf{2.456} & \textbf{.347} & \textbf{.578} \\
\midrule
2 & Target only & 250 & 3.380 & .501 & .319 \\
  & ABC, no C context & 250 & 3.423 & .502 & .365 \\
  & ABC + C context & 250 & \textbf{2.687} & \textbf{.391} & \textbf{.492} \\
\midrule
3 & Target only & 250 & 2.908 & .424 & .454 \\
  & ABC, no C context & 250 & 3.224 & .463 & .403 \\
  & ABC + C context & 250 & \textbf{2.797} & \textbf{.403} & \textbf{.504} \\
\midrule
4 & Target only & 250 & 2.812 & .415 & .430 \\
  & ABC, no C context & 250 & 3.040 & .459 & .366 \\
  & ABC + C context & 250 & \textbf{2.809} & \textbf{.409} & \textbf{.470} \\
\bottomrule
\end{tabular}
\end{minipage}
\end{center}

\begin{center}
\begin{minipage}{\linewidth}
\centering
\captionof{table}{Kimi K3 forecast and latent-response diagnostics on the same frozen
prompts. All 250 episodes parsed in every arm and prefix.}
\label{tab:coin-city-recovery-kimi}
\small
\begin{tabular}{c l r r r r}
\toprule
$k$ & Arm & $n$ & Forecast MAE & Slope MAE & $\rho(\hat g,g_C)$ \\
\midrule
0 & Target only & 250 & 3.352 & .476 & -.016 \\
  & ABC, no C context & 250 & 2.548 & .357 & -.021 \\
  & ABC + C context & 250 & \textbf{1.886} & \textbf{.266} & \textbf{.593} \\
\midrule
1 & Target only & 250 & 3.362 & .494 & .365 \\
  & ABC, no C context & 250 & 3.370 & .490 & .345 \\
  & ABC + C context & 250 & \textbf{2.382} & \textbf{.341} & \textbf{.579} \\
\midrule
2 & Target only & 250 & 3.194 & .491 & .387 \\
  & ABC, no C context & 250 & 3.114 & .487 & .352 \\
  & ABC + C context & 250 & \textbf{2.306} & \textbf{.349} & \textbf{.455} \\
\midrule
3 & Target only & 250 & 2.816 & .412 & .430 \\
  & ABC, no C context & 250 & 2.636 & .407 & .410 \\
  & ABC + C context & 250 & \textbf{2.005} & \textbf{.292} & \textbf{.609} \\
\midrule
4 & Target only & 250 & 2.378 & .357 & .444 \\
  & ABC, no C context & 250 & 1.957 & .284 & \textbf{.622} \\
  & ABC + C context & 250 & \textbf{1.840} & \textbf{.263} & .601 \\
\bottomrule
\end{tabular}
\end{minipage}
\end{center}


\begin{center}
\begin{minipage}{\linewidth}
\centering
\captionof{table}{Gemini 3.6 Flash forecast and latent-response diagnostics on the same
frozen prompts. All 250 episodes parsed in every arm and prefix.}
\label{tab:coin-city-recovery-gemini}
\small
\begin{tabular}{c l r r r r}
\toprule
$k$ & Arm & $n$ & Forecast MAE & Slope MAE & $\rho(\hat g,g_C)$ \\
\midrule
0 & Target only & 250 & 2.534 & .370 & .002 \\
  & ABC, no C context & 250 & 2.552 & .356 & .078 \\
  & ABC + C context & 250 & \textbf{2.174} & \textbf{.302} & \textbf{.559} \\
\midrule
1 & Target only & 250 & 3.733 & .550 & \textbf{.403} \\
  & ABC, no C context & 250 & 3.717 & .548 & .361 \\
  & ABC + C context & 250 & \textbf{3.300} & \textbf{.500} & .292 \\
\midrule
2 & Target only & 250 & \textbf{3.564} & \textbf{.558} & \textbf{.348} \\
  & ABC, no C context & 250 & 3.703 & .579 & .335 \\
  & ABC + C context & 250 & 3.600 & .569 & .330 \\
\midrule
3 & Target only & 250 & 3.370 & .513 & .314 \\
  & ABC, no C context & 250 & 3.092 & .460 & .397 \\
  & ABC + C context & 250 & \textbf{3.047} & \textbf{.443} & \textbf{.471} \\
\midrule
4 & Target only & 250 & 3.308 & .498 & .389 \\
  & ABC, no C context & 250 & 2.911 & .447 & .433 \\
  & ABC + C context & 250 & \textbf{2.793} & \textbf{.419} & \textbf{.447} \\
\bottomrule
\end{tabular}
\end{minipage}
\end{center}


\begin{center}
\begin{minipage}{\linewidth}
\centering
\captionof{table}{Claude Opus~5 forecast and latent-response diagnostics on the same
frozen prompts. Two responses exhausted the output budget on extended thinking and
returned no answer, one in each ABC arm.}
\label{tab:coin-city-recovery-opus5}
\small
\begin{tabular}{c l r r r r}
\toprule
$k$ & Arm & $n$ & Forecast MAE & Slope MAE & $\rho(\hat g,g_C)$ \\
\midrule
0 & Target only & 250 & 2.359 & .333 & -.018 \\
  & ABC, no C context & 250 & 2.619 & .371 & -.007 \\
  & ABC + C context & 250 & \textbf{1.764} & \textbf{.247} & \textbf{.617} \\
\midrule
1 & Target only & 249 & 3.549 & .520 & .391 \\
  & ABC, no C context & 249 & 2.303 & .338 & .370 \\
  & ABC + C context & 249 & \textbf{1.741} & \textbf{.254} & \textbf{.601} \\
\midrule
2 & Target only & 250 & 3.804 & .596 & .309 \\
  & ABC, no C context & 250 & 2.680 & .412 & .389 \\
  & ABC + C context & 250 & \textbf{2.060} & \textbf{.314} & \textbf{.524} \\
\midrule
3 & Target only & 250 & 3.022 & .461 & .379 \\
  & ABC, no C context & 250 & 2.363 & .356 & .459 \\
  & ABC + C context & 250 & \textbf{2.075} & \textbf{.309} & \textbf{.550} \\
\midrule
4 & Target only & 249 & 2.700 & .418 & .443 \\
  & ABC, no C context & 249 & 2.369 & .366 & .459 \\
  & ABC + C context & 249 & \textbf{2.201} & \textbf{.335} & \textbf{.499} \\
\bottomrule
\end{tabular}
\end{minipage}
\end{center}




\subsubsection{Misleading-cue results}

\begin{center}
\begin{minipage}{\linewidth}
\centering
\captionof{table}{Misleading-cue arm: forecast MAE and latent-slope correlation at each
City C prefix. The prompt is byte-identical to the correct-context arm except that the
City C sentence names the opposite response regime. Negative $\rho$ means the implied
response is anti-correlated with the truth, i.e. the deployment has adopted the stated
regime. Each of the six deployments contributes all 1,250 prompts.}
\label{tab:coin-city-wrong-context}
\small
\setlength{\tabcolsep}{4pt}
\begin{tabular}{l rr rr rr rr rr}
\toprule
& \multicolumn{2}{c}{$k{=}0$} & \multicolumn{2}{c}{$k{=}1$} & \multicolumn{2}{c}{$k{=}2$}
& \multicolumn{2}{c}{$k{=}3$} & \multicolumn{2}{c}{$k{=}4$} \\
\cmidrule(lr){2-3}\cmidrule(lr){4-5}\cmidrule(lr){6-7}\cmidrule(lr){8-9}\cmidrule(lr){10-11}
Model & MAE & $\rho$ & MAE & $\rho$ & MAE & $\rho$ & MAE & $\rho$ & MAE & $\rho$ \\
\midrule
Claude Opus~4.8 & 4.20 & $-.66$ & 3.88 & $-.51$ & 3.04 & $.09$ & 2.61 & $.22$ & 2.31 & $.34$ \\
GPT-5.6 & 4.23 & $-.68$ & 4.15 & $-.63$ & 2.78 & $.17$ & 2.71 & $.25$ & 2.10 & $.46$ \\
DeepSeek-V4-Pro & 4.36 & $-.62$ & 4.46 & $-.48$ & 3.61 & $.07$ & 3.39 & $.14$ & 3.04 & $.32$ \\
Kimi K3 & 4.31 & $-.62$ & 4.12 & $-.25$ & 3.03 & $.15$ & 2.72 & $.14$ & 2.18 & $.44$ \\
Gemini 3.6 Flash & 4.29 & $-.50$ & 4.40 & $-.06$ & 3.93 & $.24$ & 3.52 & $.25$ & 2.93 & $.36$ \\
Claude Opus~5 & 4.25 & $-.62$ & 3.15 & $-.21$ & 2.59 & $.27$ & 2.41 & $.41$ & 2.36 & $.47$ \\
\bottomrule
\end{tabular}
\end{minipage}
\end{center}

\subsubsection{Arbitrary-symbol induction results}
\label{app:coin-city-symbol-context}

\begin{center}
\begin{minipage}{\linewidth}
\centering
\captionof{table}{\textbf{Fifteen-system episode-randomized arbitrary-symbol
comparison at $k=0$.} Panel A reports forecast calibration; Panel B reports
response-regime discrimination. All arms share the frozen numeric task content;
Section~\ref{app:coin-city-execution} documents model-specific execution and the
Opus~4.8 cross-resource deviation. MAE is in poll points;
$\rho$ correlates implied and true latent slopes; Acc. classifies the response
regime at slope 0.575. Each $\Delta$ is symbol minus no context, with a paired
95\% episode-bootstrap interval. With no City C examples, its label can be decoded
only from City A/B examples.}
\label{tab:coin-city-symbol-context}
\scriptsize
\setlength{\tabcolsep}{2.2pt}
\input{tables/exp2_symbol_context_results}
\end{minipage}
\end{center}

\begin{center}
\begin{minipage}{\linewidth}
\centering
\captionof{table}{\textbf{Matched open-weight scaling across target-city
observations.} Symbol-minus-no-context effects are reported for Qwen2.5
32B/72B and Llama~3.1 8B/70B at each City C prefix. Negative
$\Delta$MAE favors the induced symbol mapping; positive $\Delta\rho$ and
$\Delta$accuracy indicate better regime discrimination.}
\label{tab:coin-city-symbol-scaling}
\scriptsize
\setlength{\tabcolsep}{2.2pt}
\resizebox{\linewidth}{!}{\input{tables/exp2_symbol_context_scaling}}
\end{minipage}
\end{center}

At $k=0$, Qwen3 provides a within-family contrast.
Symbol-minus-no-context $\Delta\rho$ is $.095$ $[-.060,.230]$ at 4B and $.023$
$[-.128,.171]$ at 8B, but $.398$ $[.253,.541]$ at 14B and $.302$ $[.177,.433]$
at 32B; the regime-accuracy contrasts follow the same pattern. The 4B and 8B
intervals span zero, whereas the 14B and 32B intervals do not. The 32B point
estimate sits below the 14B one, so this is a boundary among the tested Qwen3
checkpoints rather than a monotone or universal parameter threshold. The smaller checkpoints are not wholesale task
failures: their semantic-context slopes correlate $.370$ and $.286$ with truth,
versus $-.011$ and $-.014$ without context.

All four tested Qwen2.5 checkpoints discriminate under arbitrary labels
($\Delta\rho=.264/.236/.372/.381$ at 7B/14B/32B/72B), reinforcing the family
qualification. Their symbol-minus-no-context MAE changes are
$+.338/+.105/+.529/+.924$: mapping discrimination does not ensure calibrated
numerical use. Direct City C examples
make the large-Qwen symbol advantage redundant; by $k=3$--4, the 32B and 72B
$\Delta\rho$ and $\Delta$accuracy intervals span zero.

Llama~3.1 provides a second within-family contrast: $\Delta\rho=.045$
$[-.123,.250]$ at 8B and $.424$ $[.314,.543]$ at 70B. The 70B
regime-accuracy contrast is $.092$ $[.036,.152]$, although its MAE rises by
$.448$ $[.030,.886]$. In Qwen3 the mapping-recovery interval excludes zero only
from 14B upward, and in Llama only at 70B. The positive Qwen2.5-7B and
Qwen2.5-14B estimates, despite Qwen3-8B's interval spanning zero, indicate that
the tested transition depends on model family rather than parameter count. Mapping recovery and numerical calibration remain
separate outcomes.

DeepSeek improves discrimination ($\Delta\rho=.436$ $[.323,.558]$) without a
detectable zero-shot MAE change. Kimi, both Opus deployments, and GPT combine
mapping recovery with lower forecast error: their $\Delta\rho$ values are
$.633/.623/.591/.696$, while MAE falls by $.497/.584/.603/.774$ points.
Cross-provider magnitudes remain descriptive because runtimes, reasoning settings,
and output ceilings differ.

\subsubsection{Paired cue effects and interpretation}

\begin{center}
\begin{minipage}{\linewidth}
\centering
\captionof{table}{Paired forecast-MAE difference between the correct-context arm
and the context-blind ABC OLS benchmark, which the analyst fits through the origin
on every displayed A, B and available City C row. Negative values favour the
deployment. Boldface marks intervals excluding zero.}
\label{tab:coin-city-vs-ols}
\small
\setlength{\tabcolsep}{5pt}
\begin{tabular}{l ccccc}
\toprule
Model & $k{=}0$ & $k{=}1$ & $k{=}2$ & $k{=}3$ & $k{=}4$ \\
\midrule
Claude Opus~4.8 & \textbf{$-.73$} & $-.18$ & $-.05$ & $-.01$ & $-.08$ \\
GPT-5.6 & \textbf{$-.95$} & \textbf{$-.74$} & \textbf{$-.28$} & $-.21$ & \textbf{$-.24$} \\
DeepSeek-V4-Pro & $-.07$ & $.13$ & \textbf{$.56$} & \textbf{$.81$} & \textbf{$.97$} \\
Kimi K3 & \textbf{$-.66$} & $.06$ & $.18$ & $.02$ & $.01$ \\
Gemini 3.6 Flash & \textbf{$-.37$} & \textbf{$.98$} & \textbf{$1.47$} & \textbf{$1.06$} & \textbf{$.96$} \\
Claude Opus~5 & \textbf{$-.78$} & \textbf{$-.57$} & $-.07$ & $.09$ & \textbf{$.36$} \\
\bottomrule
\end{tabular}
\end{minipage}
\end{center}
\endgroup

At $k=0$ every deployment matches or beats the context-blind benchmark, and five
of the six intervals exclude zero: with no City C cases, the only way to separate
City C from the pooled A/B fit is to use its description. The advantage does not
persist once City C supplies its own rows. By $k=2$ DeepSeek-V4-Pro and Gemini 3.6
Flash are significantly worse than the benchmark, Claude Opus~5 becomes so at
$k=4$, and the remaining deployments are indistinguishable from it. The cue effect
this experiment measures is therefore a zero-shot phenomenon, and a plain
through-origin estimator remains competitive once direct target evidence
accumulates.

In the zero-shot matched-reference comparison, the target-city description changes
the influence of the A/B cases on forecasts. All six $k=0$ bootstrap intervals
exclude zero. A label-conditioned benchmark that selects a reference, estimates
its slope, and transfers that estimate reproduces the aggregate pattern; the
comparison does not identify the model's internal algorithm.
More target rows do not monotonically reduce error: DeepSeek and Opus~5 drift by
$0.34$ and $0.44$ points from $k=0$ to $k=4$, and Gemini's context advantage is
significant only at $k\in\{0,1\}$. Short, noisy prefixes can temporarily mislead
either arm.

\subsection{Does the forecast track the cue-named reference?}
\label{app:coin-city-reference-selection}

Two accounts explain the correct-versus-no-context contrast equally well. Under a
lexical-prior account the model maps the City C coverage sentence straight to a
remembered response magnitude and never reads City A or City B's numbers. Under a
reference-selection account it uses the cue to choose a reference city, estimates
that city's slope from its four displayed cases, and transfers the estimate. Both
predict a regime-level correlation, because the cue and the latent regime coincide.

The design separates them. A through-origin fit to four noisy reference cases has
a within-regime standard deviation of $.332$ against a generator standard
deviation of $.03$, so the matched reference's fitted slope correlates only $.656$
with the true City C slope where the regime indicator correlates $.996$. That
sampling error is the discriminating signal: a forecaster that is deterministic
given the cue cannot reproduce it, whereas one that transfers a fitted reference
slope inherits it. For every deployment, arm and prefix we therefore report
\[
  r\!\left(\widehat g_e,\;\widehat\beta^{\mathrm{ref}}_e \;\middle|\; R_e\right),
\]
the partial correlation between the forecast-implied slope and a reference city's
through-origin OLS slope, controlling for the response regime $R_e$ that the City
C label names. Arms with no City C label are scored against the true regime and
report the regime-matched and regime-mismatched reference instead. Intervals are
2,000-draw episode bootstraps with deterministic per-cell seeds, and the analysis
makes no model call.

\begin{figure}[!ht]
\centering
\includegraphics[width=\linewidth]{figures/exp2_reference_selection.pdf}
\caption{\textbf{The cue selects one of the two displayed references.} Partial
correlation between each deployment's forecast-implied City C slope and a
reference city's own fitted slope at $k=0$, controlling for the response regime
the City C sentence names; bars are 95\% episode-bootstrap intervals. Orange is
the reference that sentence points to, gray the other. With no reference cases
displayed neither loads; with both displayed and no cue they load equally, the
signature of pooling; the correct cue concentrates loading on the reference it
names; and an inverted cue moves that loading onto the reference it falsely names.
A prior attached to the words \emph{national} and \emph{local} predicts no
loading in any panel. Per-deployment values and the arbitrary-symbol version are in
Appendix~\ref{app:coin-city-reference-selection}.}
\label{fig:coin-city-reference-selection}
\end{figure}

\begin{center}
\begin{minipage}{\linewidth}
\centering
\captionof{table}{Partial correlation between the forecast-implied City C slope and
each reference city's fitted slope at $k=0$, controlling for the regime the City C
label names. ``Named'' is the reference the label points to, or the regime-matched
reference in the two arms without a label. Boldface marks named-reference intervals
excluding zero.}
\label{tab:coin-city-reference-selection}
\small
\setlength{\tabcolsep}{4pt}
\input{tables/exp2_reference_selection}
\end{minipage}
\end{center}

The four arms make four different predictions and all four hold. With no reference
cases displayed at all, the target-only arm loads on neither reference: eleven of
its twelve intervals span zero. With both references displayed but no City C label,
every deployment loads on the two references almost equally, the pooling signature
of a model that cannot tell which reference applies. Adding the correct label
concentrates the loading on the named reference and removes it from the other,
every ``other'' interval spanning zero. Inverting the label moves the loading
wholesale onto the reference that the false cue names, leaving the truth-matched
reference near zero.

The forecast therefore tracks the particular reference city the label names,
including that city's sampling error, whether or not the label is correct. No
prior attached to the words \emph{national} and \emph{local} reproduces this. It
also explains why implied-slope correlations plateau near $.66$ rather than near
$.996$: faithful transfer of a noisy reference estimate is capped at that
reference's own $.656$ correlation with the truth. GPT-5.6 and Claude Opus~4.8
slightly exceed the cap at $k=0$, which is what blending a reference estimate with
a regime prior produces.

\noindent\textbf{The estimator only detects an inefficiency.} The signal this analysis
reads is noise the forecast did not have to import. Each city draws its slope
independently, so a reference city's deviation from its regime mean is
uncorrelated with City C's: simulating the generator gives
$\operatorname{corr}(\widehat\beta^{\mathrm{ref}}-\mu,\;g_C-\mu)=-.0000$
(Section~\ref{app:coin-city-probe-ceilings}). Within a regime the reference's
fitted slope is therefore pure sampling error about a mean the regime already
supplies, and the response-recovery-optimal policy is to read the cue, ignore the
reference numbers, and predict the regime mean---which attains $\rho=.996$ against
the $.656$ a faithful transfer is capped at. A forecaster that did this would show
no loading on either reference and would be indistinguishable, under this
estimator, from the lexical-prior account the section argues against. The
estimator thus identifies reference selection only in forecasters that fail to
shrink, and has least power on those closest to the shrinking policy: the two
deployments above the transfer cap are the two with the highest slope
correlation, one of which also has the lowest $k=0$ forecast error of the six. We read a
positive loading as evidence for reference selection, and a null under this
estimator as uninformative rather than as evidence for a prior.

\begin{center}
\begin{minipage}{\linewidth}
\centering
\captionof{table}{The same partial correlation in the arbitrary-symbol arm at
$k=0$, where \texttt{KIV} and \texttt{ZOR} carry no stated direction and their
mapping reverses across episodes. The no-label columns repeat each model's pooling
control. Boldface marks symbol named-reference intervals excluding zero.}
\label{tab:coin-city-reference-selection-symbol}
\small
\setlength{\tabcolsep}{4pt}
\input{tables/exp2_reference_selection_symbol}
\end{minipage}
\end{center}

Applying the estimator to the arbitrary-symbol arm removes the lexical account
entirely. Twelve of the fifteen released systems have a named-reference interval
excluding zero, and the three that do not---Qwen3-4B, Qwen3-8B and
Llama~3.1-8B---are exactly the three whose symbol-minus-no-context $\Delta\rho$
intervals span zero in Table~\ref{tab:coin-city-symbol-context}. The two
statistics are computed from different quantities and agree on all fifteen
systems.

%============================================================================
% BEGIN experiment2_mechanistic_probe_appendix.tex
%============================================================================
\subsection{Linear-probe analysis of cue-conditioned representations}
\label{app:coin-city-mechanistic-probe}

\subsubsection{Behavioral screening criteria}
We screen the same held-out prompts with Qwen3-8B, Qwen3-14B, dense Qwen3-32B, and
Qwen2.5-72B-Instruct. Each checkpoint received 16 test episodes crossed with
$k\in\{0,4\}$ and correct, absent, or inverted City C context (96 prompts), with
96/96 parseable forecasts.

\begin{center}
\begin{minipage}{\linewidth}
\centering
\captionof{table}{\textbf{Behavioral screening results for the linear-probe analysis.}
Entries are paired mean differences with 95\% episode-bootstrap intervals
($n=16$). Evidence is no-context $k=4$ minus $k=0$ aligned slope; selection is
correct minus absent context at $k=0$; reversal is inverted minus correct context
at $k=0$; forecast penalty is inverted minus correct poll MAE at $k=0$. Positive
evidence and selection, negative reversal, and a positive forecast penalty are
the predicted directions.}
\label{tab:coin-city-behavioral-gate}
\scriptsize
\setlength{\tabcolsep}{2.4pt}
\resizebox{\linewidth}{!}{\input{tables/exp2_qwen_behavior_gate}}
\end{minipage}
\end{center}

The 14B and 32B intervals exclude zero in all four predicted directions. For 72B,
the evidence, reversal, and forecast-penalty intervals exclude zero, while the
correct-versus-absent selection interval spans zero. All four 8B intervals span
zero. We use 32B for the dense Qwen3 scaling comparison and 72B for a
cross-generation comparison; neither is interpreted as another point on a
monotonic size curve. Qwen3-4B and Llama-3.1-8B use the identical task file,
held-out split, and model, layer, target, and endpoint selection rules.

\subsubsection{Activation and probe protocol}
We selected 64 episodes without replacement, balanced over which reference city
was strongly responsive and whether City C's latent regime was strong or weak.
Forty-eight episodes form a development set and 16 form a sealed test set. Each
episode is rendered at $k=0$ and $k=4$ under correct, absent, and inverted context,
giving 384 activation records. Within an episode-prefix pair the three prompts
have identical numeric content. The only intervention is the City C relationship
description.

For the final input token, we save the embedding output and every transformer-block
output. The resulting tensor dimensions are $[384,37,2560]$ for Qwen3-4B;
$[384,37,4096]$ for Qwen3-8B; $[384,33,4096]$ for Llama-3.1-8B;
$[384,41,5120]$ for Qwen3-14B; $[384,65,5120]$ for Qwen3-32B; and
$[384,81,8192]$ for Qwen2.5-72B-Instruct. We fit separate dual-ridge probes at
each layer and $k$, using only correct-context development episodes. Four
cell-balanced development folds select the ridge strength from
$\{10^{-4},3\!\times\!10^{-3},10^{-1},3,100\}$ and then select one reported
layer. The untouched test episodes are evaluated in all three context arms
($n=16$ per arm). A 10,000-repeat held-out label permutation test keeps the
development-selected layer, ridge strength, and fitted probe fixed. Thus neither
the test labels nor test-arm scores select the reported layer.

We probe three increasingly strict targets: the binary strong/weak regime; the
full realized slope $g_C$; and the within-regime residual
$g_C-.90$ for strong episodes or $g_C-.25$ for weak episodes. The first two are
recoverable from the displayed evidence and the third is not, by construction:
Section~\ref{app:coin-city-probe-ceilings} derives what a Bayes-optimal reader of
the same prompt could score on each, and the residual is included as a
specificity control rather than as a target a model could hit. Mean-pooled
embedding and final-layer states are additional pooling controls.

\begin{center}
\begin{minipage}{\linewidth}
\centering
\captionof{table}{\textbf{Six-model correct-context probe summary.}
Each entry gives $k=0/k=4$ on the 16 sealed episodes. Layers and ridge strengths
were selected on development episodes only. Residual $p$ is the one-sided
10,000-repeat held-out label-permutation value.}
\label{tab:coin-city-mechanistic-summary}
\small
\setlength{\tabcolsep}{3.0pt}
\resizebox{\linewidth}{!}{\input{tables/exp2_open_model_mechanistic_summary}}
\end{minipage}
\end{center}

\begin{center}
\begin{minipage}{\linewidth}
\centering
\captionof{table}{\textbf{Held-out layer-wise open-weight probe.} AUC is reported
for the regime target and $R^2$ for both slope targets. Correct, none, and wrong
score the same development-selected probe against the true City C target under
the three context arms; wrong/cue instead scores the inverted arm against the
relationship named by its cue. $p$ is the one-sided held-out label-permutation
value on the correct-context test episodes.}
\label{tab:coin-city-mechanistic-probe}
\scriptsize
\setlength{\tabcolsep}{2.4pt}
\resizebox{\linewidth}{!}{\input{tables/exp2_qwen3_mechanistic_probe}}
\end{minipage}
\end{center}

\begin{figure}[ht]
  \centering
  \includegraphics[width=.62\linewidth]{figures/exp2_qwen3_mechanistic_main.pdf}
  \caption{\textbf{Semantic-arm slope decoding follows cue assignment.}
  Each line is a sealed model--$k$ cell across the six open checkpoints of the
  16-episode roster; purple is the median. Probes are fit on correct-context
  development episodes. Under inversion, held-out $R^2$ collapses against the true
  slope and recovers against the cue-assigned slope. The input-embedding control
  below shows this signal is available before any computation, which is why the
  arbitrary-label replication rather than this figure carries the
  representational claim.}
  \label{fig:coin-city-mechanistic-main}
\end{figure}

\begin{figure}[ht]
  \centering
  \includegraphics[width=\linewidth]{figures/exp2_qwen3_mechanistic_probe.pdf}
  \caption{\textbf{Development layer trajectories for the six-model probe.}
  Curves show four-fold development CV rather than test performance; stars mark
  the development-selected layer. Layer depth is normalized
  because the checkpoints have 32--80 transformer blocks. The regime and
  full-slope signals are broad, whereas residual-slope CV is weak and does not
  survive held-out evaluation.}
  \label{fig:coin-city-mechanistic-probe}
\end{figure}

\subsubsection{Regime and full-slope decodability}
Across the six checkpoints and two evidence depths, correct-context regime AUC
is $.844$--$1.000$ and full-slope $R^2$ is $.627$--$.975$
(Table~\ref{tab:coin-city-mechanistic-summary}). This information is not a
context-invariant encoding of the true coefficient: removing the City C
description drives selected-layer full-slope $R^2$ below zero in all 12
model--$k$ cells. Under inversion, true-target regime AUC falls to
$.000$--$.109$ while cue-target AUC is $.891$--$1.000$; the full-slope probe
likewise fits the cue-assigned relationship better than the true one. The Llama
replication shows that this pattern is not specific to the Qwen checkpoints.
Across both model families, the linearly available signal therefore tracks the
response family that context assigns to City C.

\subsubsection{What each probe target can be scored against}
\label{app:coin-city-probe-ceilings}

A probe is scored against a latent generator quantity, so a null is only a fact
about a model if the displayed evidence determines that quantity. We therefore
derive, for each target and evidence depth, the score a Bayes-optimal reader of
the same prompt would obtain. This is an analytic property of
Equation~\ref{eq:coin-city-dgp} and its row construction, computed in closed form
and re-derived by simulating the generator's own row logic; the two agree to
within $.005$ and no model output enters; the derivation is
\path{exp2_v2/biased_news/analysis/compute_bayes_ceilings.py}.
Table~\ref{tab:coin-city-probe-ceilings} gives the ceilings. The same module
reports forecast-MAE floors: $2.28$ poll points at $k=0$ with no cue, $0.17$ for
an oracle told the regime means, and $1.78$ for a reader who must instead estimate
the named regime's slope from its four displayed rows.

The same simulation records one quantity the reference-selection analysis of
Section~\ref{app:coin-city-reference-selection} depends on. Cities draw their
slopes independently, so a reference city's deviation from its regime mean
carries no information about City C's:
$\operatorname{corr}(\widehat\beta^{\mathrm{ref}}-\mu,\;g_C-\mu)=-.0000$ over
400,000 simulated episodes, against a reference-fit standard deviation of $.318$
within regime. That is why a loading on the reference's own sampling error
identifies reference selection, and equally why a forecaster that shrank
optimally would show none.

\begin{center}
\begin{minipage}{\linewidth}
\centering
\captionof{table}{Maximum score a Bayes-optimal reader of the same prompt can
obtain on each probe target, under the frozen data-generating process. The regime
and the full slope are recoverable; the within-regime residual is not, because it
is generated with standard deviation $.03$ while a through-origin fit to City C's
displayed rows carries a standard deviation of $.318$. At $k=0$ City C displays no
rows at all and the residual ceiling is exactly zero.}
\label{tab:coin-city-probe-ceilings}
\small
\input{tables/exp2_probe_target_ceilings}
\end{minipage}
\end{center}

\subsubsection{Specificity: within-regime residuals}
The residual target is the specificity control, and it behaves as a target with
no available signal should. The 12 correct-context residual-slope $R^2$ values
range from $-.531$ to $.003$, none of their permutation tests rejects the null
($p=.221$--$.927$), and 33 of 36 arm-level residual scores are nonpositive.
Against a ceiling of $.000$ at $k=0$ and $.009$ at $k=4$, that is the result the
design permits: no reader of these prompts, and no representation computed from
them, could score materially above zero. We therefore read the residual null as
evidence that the probe does not manufacture structure where none is available,
and not as a limit on what these checkpoints encode. A probe that returned a
positive residual score here would be the finding, and it would be a warning
about the probe.

The two recoverable targets do carry information, and the input-embedding control
is what bounds their interpretation. Mean-pooled input embeddings for every Qwen
checkpoint already give perfect regime AUC and full-slope $R^2=.988$--$.993$ under
correct context---at the $.992$ ceiling---because the explicit description names a
regime and the regime means account for essentially all slope variance. Llama
input embeddings likewise give perfect regime AUC and full-slope $R^2=.756/.827$.
The layer-wise semantic result is therefore evidence that the cue-selected
response family is linearly represented, and reaches nothing further about the
continuous $g_C$: the part of $g_C$ beyond the family is not in the prompt to
begin with. The original behavioral inversion screen connects this representation
to output sensitivity in the three larger screened checkpoints: inverted context
raises $k=0$ forecast MAE by $1.66$, $1.96$, and $2.21$ points for 14B, 32B, and
72B, respectively.

\subsubsection{Power-matched and cross-model arbitrary-symbol replication}
\label{app:coin-city-probe-power}

The six-checkpoint roster above seals 16 test episodes per arm. Its exact
label-permutation null spans ROC AUC $[.203,.797]$, so it can separate a score
from chance only above about $.75$. That resolution suffices for the semantic arm,
whose scores sit at ceiling, but it cannot evaluate a weaker signal at all. We
therefore evaluate Qwen3-14B at the design's balanced maximum---248
episodes, 160 development and 88 sealed test, 1,488 activation records---once with
the semantic cue and once with the episode-randomized \texttt{KIV}/\texttt{ZOR}
labels. Anchor, targets, fold structure, ridge grid, layer-selection rule, pooling
controls and the 10,000-repeat held-out permutation test are unchanged; only the
number of episodes differs. The 88-episode null reaches AUC $.603$. The enlarged
development set is disjoint from every episode used by the 16-episode analysis,
providing an independent estimate rather than a re-slice.

\begin{center}
\begin{minipage}{\linewidth}
\centering
\captionof{table}{\textbf{Power-matched Qwen3-14B probe, 88 sealed episodes per
arm.} Correct, none and wrong score the development-selected probe against the true
City C target under the three cue arms; wrong/cue scores the inverted arm against
the relationship its cue names. ``Layer'' is the development-selected block of 40
and ``Embed.'' is the mean-pooled input-embedding control under correct context.
$p$ is the one-sided held-out label-permutation value.}
\label{tab:coin-city-probe-power}
\scriptsize
\setlength{\tabcolsep}{3.2pt}
\resizebox{\linewidth}{!}{\input{tables/exp2_probe_power_matched}}
\end{minipage}
\end{center}

The two arms separate sharply, and the input-embedding control is what separates
them. Under the semantic cue the development-selected layer is the first
transformer block for three of the four conditions, and mean-pooled input
embeddings alone already reach regime AUC $1.000$ and full-slope $R^2=.989$: the
decodable quantity is the cue token itself. Under arbitrary labels the same control
collapses to AUC $.555$ and $R^2=.003$, the selected layers move to blocks 26 and
40 of 40, and held-out regime AUC is $.678$ at $k=0$ ($p=.0016$) and $.713$ at
$k=4$ ($p<.001$). The induced label-to-regime mapping is therefore linearly
represented, but only after computation, far more weakly than a regime the model
was told outright, and below the $.979$ AUC a Bayes-optimal reader of the same two
reference cities could reach. Unlike the residual target, the full slope is
recoverable here---its ceiling is $.992$---so its symbol-arm score of $R^2=.016$ at
$k=0$, which does not reject its null, is a genuine gap rather than a property of
the design.

Representation and behavior agree closely once the cue must be induced. On the same
88 sealed episodes the model's own forecasts give regime AUC $.702$, $.518$ and
$.331$ under correct, absent and inverted labels, against probe AUCs of $.678$,
$.467$ and $.321$; inversion raises $k=0$ poll MAE from $2.49$ to $4.28$. With four
City C cases the inverted-arm representation reverses: true-target regime AUC rises
from $.321$ to $.747$, and the full-slope probe fits the cue-named relationship at
$R^2=-.855$ while fitting the true one at $+.126$. Direct target evidence overrides
the induced label in the representation as it does in the forecast.

Six further checkpoints received the identical symbol protocol at the same 88
sealed episodes: Qwen2.5-32B and Qwen2.5-72B, which change model generation;
Qwen3-8B and Qwen3-4B, which hold family, tokenizer and architecture fixed and
change only scale; and Llama~3.1-8B and Llama~3.1-70B, which do the same within a
second vendor and architecture. Held-out regime AUC at $k=0$ is $.778$ for
Qwen2.5-32B ($p<.001$, selected layer 50 of 64), $.728$ for Llama~3.1-70B
($p<.001$, layer 75 of 81), $.662$ for Qwen2.5-72B ($p=.003$, layer 63 of 80),
$.566$ for Qwen3-8B ($p=.15$), $.520$ for Qwen3-4B ($p=.38$) and $.423$ for
Llama~3.1-8B ($p=.89$). No null checkpoint rejects in any $k=0$
condition---regime, slope or residual slope. Mean-pooled input embeddings stay
between $.495$ and $.556$ in every symbol run, so no symbol arm carries the
mapping before computation.

Decodability tracks behavioral recovery across all seven checkpoints, four
positive and three null. The four whose forecasts recover the mapping are the four
in which it decodes (Llama~3.1-70B, $\Delta\rho=.424$; Qwen3-14B, $.398$;
Qwen2.5-72B, $.381$; Qwen2.5-32B, $.372$), and the three whose forecasts do not
are the three in which it does not (Qwen3-4B, $.095$; Llama~3.1-8B, $.045$;
Qwen3-8B, $.023$; all three intervals spanning zero). The Llama~3.1
contrast carries the weight, because it holds family, tokenizer and architecture
fixed, varies only scale, and is the one within-family contrast whose behavioral
classification survives the harder generator population of
Section~\ref{app:coin-city-robustness}: 8B is null on both populations and 70B
positive on both, while the probe runs $.423$ to $.728$ across them. Because the
null and the positive share an architecture and a tokenizer, the split is not an
artifact of either.

The Qwen3 sequence---$.520$, $.566$, $.678$ from 4B to 14B---is consistent with
this but does not add independent weight, because the population that would test
it is the one on which its behavioral boundary moves. That movement is smaller
than the $\Delta\rho$ column alone suggests. Symbol-arm regime accuracy is
$.512$, $.516$, $.660$, $.656$ originally against $.528$, $.554$, $.589$, $.606$
on the harder population. The larger pair remains above the smaller pair, but
14B and 32B exchange order and the between-group separation shrinks. The
$\Delta\rho$ contrast then subtracts a no-context
baseline whose standard error is about $.064$, and baselines of $-.036$ for
Qwen3-8B and $+.025$ for Qwen3-14B are enough to put one interval above zero and
the other $.004$ below it, with the two intervals overlapping by $78\%$. We
therefore read the harder population as preserving the broad size-group contrast
without locating the significance threshold at a particular checkpoint.
Section~\ref{app:coin-city-robustness} should be read alongside the
seven-for-seven agreement, which is evidence that the probe reads what the
forecasts express on the episodes both were measured on and not evidence for a
scale law. The positives span three model generations and two
vendors, and every positive shows the same shape: the decoded regime follows the
cue at $k=0$ and the truth at $k=4$.
Qwen3-4B and Qwen3-8B do reject at $k=4$ ($.685$, $p=.002$ and $.610$, $p=.04$),
which is expected rather than contrary: at $k=4$ the displayed target cases
exhibit the regime whatever the label says, and the mean-pooled input embeddings
rise to $.585$ and $.621$ accordingly, in Qwen3-8B's case above the probe itself.
The induction claim is a $k=0$ claim throughout. The sealed test is 88 episodes
per model, well above the 16 the roster began with but not a large sample, and the
permutation nulls are computed rather than assumed.

At 16 episodes, the same arbitrary-symbol procedure returns $k=0$ regime AUC
$.312$ with $p=.90$. Its $[.203,.797]$ null is too broad to resolve a signal of
the magnitude observed at 88 episodes.

\subsubsection{Causal test: patching the City C label representation}
\label{app:coin-city-causal-patch}

Decodability is not mediation. To test whether the representation the probe reads
is the one the forecast uses, we intervene on it directly. Under arbitrary labels, an episode's correct and swapped prompts are
positionally aligned through the City C label and differ at exactly two label-token
positions.
We replace the residual stream at those two positions with the states the opposite
arm produces, in both directions, and regenerate the forecast greedily. The
endpoint is the episode-level mean of the regime-aligned correct-into-swapped
shift and the sign-reversed swapped-into-correct shift, tested one-sided against
an episode sign-flip null.

The analysis uses three patch sites and two controls. The primary site is a
three-layer window centred on the development-selected relational layer;
patching at the embedding output only, and at every causally effective layer, are
the comparison sites; the unpatched arm and a self-patch that writes each state
back over itself are the controls. Layer, window and ridge were selected on
development episodes with the sealed test unread.

\begin{center}
\begin{minipage}{\linewidth}
\centering
\captionof{table}{\textbf{Causal effect of patching the two City C symbol
subtokens, Qwen3-14B.} Shift is the regime-aligned symmetric patch mean; positive
means the forecast moved toward the relationship the donor label names. The
16- and 88-episode analyses select their windows independently on
their own development sets. The null interval is the 88-episode sign-flip null.}
\label{tab:coin-city-causal-patch}
\small
\setlength{\tabcolsep}{4pt}
\resizebox{\linewidth}{!}{\input{tables/exp2_causal_patch}}
\end{minipage}
\end{center}

At $k=0$ the intervention moves the forecast: patching the selected window shifts
the regime-aligned prediction by $.333$ against a null interval of
$[-.177,.178]$ ($p<.0001$, 82 complete episodes), and patching the embedding
alone gives $.319$. All eight stability checks pass, including $347/347$ exact
self-patch matches.

Two readings of the table matter more than the significance. First, the
embedding-only and all-layer rows are identical, and necessarily so: the two
prefixes differ only at the patched label positions, so overwriting their embeddings
and letting the model recompute reproduces what overwriting every later state does. They are one intervention reported at two depths, not two
converging results. What distinguishes the mid-window row is that it reaches the
same magnitude while touching only three layers late in the network.

Second, the effect is confined to $k=0$. With four City C cases every patch site
returns a null ($p=.83$ and $p=.64$), so the induced label no longer controls the
forecast once direct target evidence is available. That is the revision claim of
Section~\ref{app:coin-city-results} in causal form, and it matches the
representational result: at $k=4$ the inverted-arm probe recovers the true regime
rather than the cue-assigned one.

At 16 episodes, the primary patch estimate is $-.033$ at $p=.75$ and the sign-flip
null spans $[-.387,.387]$. Its 64-episode development set selects layers
$33$--$35$, whereas the 160-episode development set selects $19$--$21$ at a
higher cross-validation score ($.903$ against $.806$).

\noindent\textbf{Permutation implementation.} The test enumerates all $2^{n}$ sign
assignments for 16 episodes; exhaustive enumeration is infeasible for 88
($2^{88}\approx3\times10^{26}$). Enumeration is retained for $n\le20$; above it
the same null is sampled with $200{,}000$ draws under a fixed seed, giving a
Monte Carlo standard error of $.0005$ at $p=.05$. The estimand, endpoint, and
direction are identical under both implementations.

\noindent\textbf{Scope.} This is one checkpoint under arbitrary labels. It shows that
the label representation Qwen3-14B computes is causally responsible for its
zero-shot regime assignment, and that direct evidence displaces it. It does not
establish the same for the hosted deployments, the semantic cue, or other
architectures.

\subsubsection{Reproducibility and scope}
The checkpoints are pinned to commits
\texttt{cdbee75f17c01a7cc42f958dc650907174af0554} (Qwen3-4B),
\texttt{b968826d9c46dd6066d109eabc6255188de91218} (Qwen3-8B),
\texttt{0e9e39f249a16976918f6564b8830bc894c89659} (Llama-3.1-8B),
\texttt{40c069824f4251a91eefaf281ebe4c544efd3e18} (Qwen3-14B),
\texttt{9216db5781bf21249d130ec9da846c4624c16137} (Qwen3-32B), and
\texttt{495f39366efef23836d0cfae4fbe635880d2be31} (Qwen2.5-72B).
The activation hashes are
\texttt{e1913859\allowbreak{}63d639aa\allowbreak{}5b8574a9\allowbreak{}abbeb91b\allowbreak{}3519aed9\allowbreak{}34ac491f\allowbreak{}86365563\allowbreak{}1b26dbbc}
(Qwen3-4B) and
\texttt{7f65aeec\allowbreak{}7a91e404\allowbreak{}cc600686\allowbreak{}18fedf25\allowbreak{}c679abde\allowbreak{}5d48c2b8\allowbreak{}477f7366\allowbreak{}70fbc54f}
(Llama-3.1-8B); all six hashes and tensor dimensions are enforced by the
fail-closed generator. Frozen tasks, activation tensors, extraction manifests,
generated behavior, selected predictions, complete layer curves, and probe
results are under
\path{exp2_v2/biased_news/data/coin_city_stable_relationship_claude_n250_v4/mechanistic_probe/}.
The analysis and asset generator are
\path{exp2_v2/biased_news/mechanistic_probe/analyze_probe.py} and
\path{paper/generate_coin_city_mechanistic_comparison.py}. The causal run is
\path{mechanistic_probe/qwen3_14b_symbol_relational_v2/}, frozen by
\path{freeze_symbol_relational_protocol.py} and executed by
\path{extract_symbol_label_states.py}, \path{patch_symbol_label_activations.py}
and \path{analyze_symbol_relational_probe.py}; every count in that pipeline is
derived from the frozen task manifest and checked by
\path{tests/test_symbol_relational_scaling.py}.
\path{paper/generate_coin_city_causal_patch_table.py} fails closed on the run's
study name, sealed-test size and stability gate. The power-matched runs are
\path{mechanistic_probe/qwen3_14b_probe_v2/},
\path{mechanistic_probe/qwen3_14b_symbol_probe_v2/},
\path{mechanistic_probe/qwen3_8b_symbol_probe_v2/},
\path{mechanistic_probe/qwen3_4b_symbol_probe_v2/},
\path{mechanistic_probe/llama3_1_70b_symbol_probe_v2/},
\path{mechanistic_probe/qwen2_5_72b_symbol_probe_v2/},
\path{mechanistic_probe/qwen2_5_32b_symbol_probe_v2/} (checkpoint
\texttt{5ede1c97bbab6ce5cda5812749b4c0bdf79b18dd}) and
\path{mechanistic_probe/llama3_1_8b_symbol_probe_v2/} (checkpoint
\texttt{0e9e39f249a16976918f6564b8830bc894c89659}), built by
\path{exp2_v2/biased_news/mechanistic_probe/build_probe_tasks.py} and
\path{build_arbitrary_symbol_probe_tasks.py} with 62 and 31 episodes per factorial
cell; \path{paper/generate_coin_city_probe_power_table.py} regenerates their table
and validates the sealed-test size and task hash. Each run's
activation tensor is regenerable from its frozen task file and pinned checkpoint;
the extraction manifests record the tensor hashes rather than the tensors.
Cross-vendor probe evidence here is the Llama~3.1-8B null rather than a
positive. The diagnostic covers six checkpoints spanning Qwen and Llama at 16 sealed episodes per arm, plus
the matched 88-episode semantic and arbitrary-symbol Qwen3-14B pair. For the semantic
cue it is correlational and does not establish that the decoded direction
mediates the forecast; under arbitrary labels
Section~\ref{app:coin-city-causal-patch} provides a causal intervention for one
checkpoint. Both complement the paired behavioral interventions above.
%============================================================================
% END experiment2_mechanistic_probe_appendix.tex
%============================================================================

\subsection{Robustness across generator populations and deployments}
\label{app:coin-city-robustness}

\subsubsection{Generator-population robustness}
Fresh seeds 940000--940249 target strong/weak means of \(.75/.40\)
with slope SD \(.10\), jointly narrowing the regime gap and widening
within-regime dispersion. The realized target slopes are
\(.752\pm.100\) and \(.404\pm.105\). Draws were neither truncated nor reordered:
two reference pairs reverse their nominal latent ordering and remain in the
analysis. Case noise, factorial balance, prompts, prefix set, local checkpoint
settings, and episode-randomized symbol mapping are unchanged.

All ten open checkpoints were run on the three matched arms, producing 37,500
response records; 37,355 contain a numeric forecast, and unresolved formats remain
missing. Table~\ref{tab:coin-city-generator-robustness} compares the
\(k=0\) symbol-minus-no-context contrast across generator populations. Bold
intervals exclude zero; negative \(\Delta\)MAE would favor symbol context.
Because this population changes separation and dispersion together, it tests the
joint perturbation rather than identifying either parameter's effect separately.

\begin{center}
\begin{minipage}{\linewidth}
\centering
\captionof{table}{Arbitrary-symbol recovery on the original and harder generator
populations. Intervals are 95\% episode-bootstrap intervals with 5,000 draws;
\(n\) is the harder-population complete-case count.}
\label{tab:coin-city-generator-robustness}
\scriptsize
\setlength{\tabcolsep}{2.5pt}
\resizebox{\linewidth}{!}{\input{tables/exp2_generator_population_robustness}}
\end{minipage}
\end{center}

Which Qwen3 checkpoints clear significance changes, while the larger pair retains
higher symbol-arm regime accuracy than the smaller pair.
Qwen3-8B moves from an original null to \(\Delta\rho=.231\) \([.074,.370]\) while
Qwen3-14B falls to \(.155\) \([-.004,.303]\); Qwen3-4B remains null and Qwen3-32B
remains positive. Regime accuracy in the symbol arm, which subtracts no baseline,
is \(.512\), \(.516\), \(.660\), \(.656\) originally against \(.528\),
\(.554\), \(.589\), \(.606\) here. The 14B and 32B checkpoints exchange order,
and the smaller models improve slightly, but the gap between 8B and 14B shrinks
from \(.144\) to \(.035\).
The \(\Delta\rho\) contrast then subtracts a no-context baseline carrying a
standard error of about \(.064\), and baselines of \(-.036\) for Qwen3-8B against
\(+.025\) for Qwen3-14B suffice to place one interval above zero and the other
\(.004\) below it, with the two overlapping by \(78\%\). We read this as a
detection threshold that is no longer locatable to a checkpoint on the compressed
population; it does not support a fixed scale threshold.

All four Qwen2.5 intervals remain positive, from \(.149\) \([.009,.279]\) at 7B to
\(.336\) \([.204,.463]\) at 72B, and both Llama checkpoints keep their
classification, 8B null in each population and 70B positive in each. Seven of ten
intervals therefore exclude zero in each population and eight of ten checkpoints
keep their classification, the Qwen3 8B/14B pair being the two that do not. Mapping
recovery survives the substantially less separable population; the checkpoint at
which it becomes detectable does not, and neither does a scale reading across
families, since Qwen2.5-14B sits below Qwen2.5-7B in both. Part of the shrinkage
is mechanical: the regime explains less of the slope here, capping \(\rho\) at
\(.87\) against \(.996\). The observed \(\Delta\rho\) falls further than that alone
predicts. Numerical use stays distinct:
Qwen2.5-72B's discrimination improves while symbol context raises forecast MAE by
\(1.012\) \([.717,1.316]\).


\subsubsection{One independent hosted repeat}
GPT-5.6 was replayed once on the 250 \(k=0\) prompts in each no-context,
semantic-context, and arbitrary-symbol arm. The deployment, low reasoning effort,
provider-native sampling, prompt hashes, parser, and one-attempt policy were
unchanged; all 750 repeat responses parsed.

\begin{center}
\begin{minipage}{\linewidth}
\centering
\captionof{table}{Independent GPT-5.6 \(k=0\) repeat. Paired intervals resample
episodes; mean \(|\Delta\hat y|\) is the per-task absolute change in forecast.}
\label{tab:coin-city-gpt-repeat}
\scriptsize
\setlength{\tabcolsep}{2.5pt}
\resizebox{\linewidth}{!}{\input{tables/exp2_gpt56_repeat}}
\end{minipage}
\end{center}

Per-task mean absolute forecast changes are \(.166\), \(.193\), and \(.167\)
points across the three arms. The arbitrary-symbol-minus-no-context correlation
effect is \(.696\) originally and \(.661\) on repeat; their difference is
\(-.035\) \([-.075,.003]\). The corresponding forecast-MAE benefit is
\(-.774\) originally and \(-.743\) on repeat, with difference \(+.031\)
\([-.046,.104]\). This one repeat supports stability of the paired conclusions
for this named snapshot, not a general estimate of provider or model variance.

\subsection{Missingness, scope, and artifact map}
\label{app:coin-city-limitations}

\subsubsection{Missingness and record recovery}
Claude Opus~5 has two unparseable primary task cells: one correct-context response
at \(k=1\) and one no-context response at \(k=4\). Both exhausted the output
budget during extended thinking and returned no answer. GPT-5.6 has one
no-context failure at \(k=1\), where the provider returned a model-side HTTP 500
before any response. None was rerun. Complete-case pairing therefore gives
\(n=249\) for the corresponding deployment-prefix rows and \(n=250\) elsewhere.
These three failures constitute 0.013\% of the 22,500 primary task cells.

Claude Opus~4.8, DeepSeek-V4-Pro, Kimi K3, and Gemini 3.6 Flash have 3,750
parseable primary responses each after the one offline Claude key-whitespace
normalization described above. That normalization changed neither the forecast
value nor the raw model answer and made no additional API call.

All six misleading-context analytic files contain exactly 1,250 unique parseable
tasks. Concurrent writes in the Kimi collection produced 11 duplicate records and
14 malformed intermediate lines. Construction of the analytic file retained the
earliest response for each task and identified 11 absent task IDs; those 11 tasks
were collected under the unchanged Kimi configuration. Thus 7,500 is the number
of final unique analytic records for the misleading arm rather than the number of
intermediate records or provider-level attempts.

The expanded arbitrary-symbol evaluation contains fifteen response-record
triplets, or 56,250 records. Within every model, each arm contains 1,250
unique task IDs matched to the frozen task-file prompt hash. Original outputs are
preserved. When a model placed an explicit arithmetic expression rather than a
number in \texttt{predicted\_poll}, a no-call reparse accepts only numeric constants,
parentheses, unary signs, and the four arithmetic operators, and requires
the result to lie in the primary parser's 0--100 range. For Llama~3.1-8B, this
reparse leaves 1,234/1,241/1,232 parseable no-context, semantic-context, and
symbol records, respectively. Its matched complete-case $n$ is 242 at $k=0$ and
ranges from 238 to 246 across the five prefixes. Qwen3-32B retains one
unresolved symbol response, leaving 1,249 parseable symbol records; every other
model parses all 1,250 in every arm, and every model except Llama~3.1-8B has
$n=250$ at $k=0$. The final table reports these complete-case denominators, and unresolved
outputs remain missing rather than being re-asked.

\subsubsection{Interpretive scope}
The aggregate behavior is consistent with context-conditioned response-regime
assignment and integration of direct City C evidence in this controlled synthetic
setting. A label-conditioned analogical-regression benchmark---map the label to a
reference, estimate that reference's slope, transfer it to City C, and revise
with City C's own cases---is sufficient to reproduce the pattern but does not
identify the model's internal algorithm. In the misleading arm, reversing only
the City C label reverses forecast assignment while all numeric content remains
fixed.

The primary binary semantic cue is aligned by construction with the higher- or
lower-response family and has a fixed meaning across episodes. The arbitrary-symbol
control removes that meaning by reversing the \texttt{KIV}/\texttt{ZOR} mapping
across episodes; it therefore tests whether the semantic-arm result can arise
solely from a pre-existing national-greater-than-local prior. The expanded
15-model comparison shows detectable symbol-based discrimination for Qwen3-14B
and Qwen3-32B, all four tested Qwen2.5 checkpoints, Llama~3.1-70B, and all five
hosted systems.
In this original population, intervals span zero for Qwen3-4B, Qwen3-8B, and
Llama~3.1-8B. Within Qwen3 the mapping is recovered only from 14B upward, and
within Llama only at 70B,
providing within-family contrasts on the original population; Qwen2.5-7B's positive
result alongside Qwen3-8B's interval spanning zero, and Qwen3-32B's smaller
point estimate than Qwen3-14B's, show that the transition is neither monotone in
nor determined by parameter count alone.
Four quantitative analyses bear on whether a model induces the arbitrary
label-to-regime mapping, and they are computed from different quantities: the
symbol-minus-no-context correlation contrast on forecasts; the partial correlation
between the implied slope and the cue-named reference's own fitted slope
(Section~\ref{app:coin-city-reference-selection}); held-out linear decoding of
hidden states; and causal patching of those states
(Section~\ref{app:coin-city-causal-patch}). Where more than one was run on a
checkpoint they agree, including Llama~3.1-8B, where every analysis run is null.
Agreement among measures that could have disagreed is the reason to prefer the
induction account over a lexical or pooling one. It does not extend the claim to
the hosted deployments, whose internals are unavailable, and only one checkpoint
carries the causal test. A fifth, descriptive audit asks what models themselves
say in their one-sentence rationales.

Deterministic lexical rules were developed on a hash-selected 50-episode subset and
then applied without revision to the remaining
200 episodes. Across all three comparison arms, 56,057 of 56,250 response records
contain a rationale. At $k=0$ in the arbitrary-symbol arm, 2,992 of 3,000 held-out
records contain one; the coder identifies a selected City A/B reference in 1,832,
and every codable statement names the reference that actually shares City C's
episode-local symbol. In the inverted semantic arm, all 1,200 held-out records
contain a rationale: all 1,089 codable reference statements name the reference
selected by the false cue, and all 1,186 codable national/local statements name
the false regime. These outputs are behavioral self-reports, not privileged
evidence of computation; in particular, the coder abstains whenever a statement
does not make an unambiguous reference or regime claim. The conditional pattern
also appears in the three checkpoints with null behavioral symbol contrasts, so
it establishes declared cue interpretation rather than successful numerical use.

Because discrimination and forecast MAE can move in opposite directions, mapping
induction is not equivalent to calibrated numerical use. Ambiguous, continuous,
and adversarial cues remain untested. All prompts also tell models that within-city
response is stable and explain the matched-pair structure.

The query target is a latent conditional mean with no newly sampled observation
noise. This isolates response recovery but is not equivalent to calibration on a
future noisy poll. The six primary deployments differ in provider, reasoning mode,
output budget, and service tier. The symbol module mixes local
vLLM checkpoints with hosted calls and includes the documented cross-resource
Opus~4.8 arm. Cross-model comparisons are therefore descriptive rather than
configuration-matched capability rankings. The cue-inversion module is a matched
stress test over the same task identities and numeric inputs.

The primary and harder designs are two reproducible seeded populations in the
same synthetic world. The second jointly narrows the regime gap and widens slope
dispersion, but two settings do not characterize generator-parameter uncertainty
or isolate those changes. Bootstrap intervals resample episodes within a
population and are pointwise; they do not incorporate provider drift or correct
for the many model, arm, and prefix comparisons. The evaluated populations also do
not characterize factorial variation in generator parameters, ambiguous cues, or
calibration against noisy held-out outcomes.

\subsubsection{Reproducibility and artifact map}
The frozen data root is
\path{exp2_v2/biased_news/data/coin_city_stable_relationship_claude_n250_v4/},
and the experiment-code root is \path{exp2_v2/biased_news/}. Paths below are
relative to those roots unless a repository-level path is shown.
\begin{description}
  \item[Design and preflight.]
  Under the data root, \path{design/manifest.json} and
  \path{design/validation.json} freeze and validate the primary design;
  \path{design/tasks_abc_wrong_context.manifest.json} records the misleading arm,
  and \path{design/tasks_abc_symbol_context.manifest.json} records the balanced,
  episode-randomized symbol control.

  \item[Generator and task builders.]
  Under the code root, the generator is
  \path{engine/coin_city_stable_relationship_claude_n250.py}; the builders are
  \path{eval/build_coin_city_stable_relationship_claude_n250.py} and
  \path{eval/build_coin_city_wrong_context_tasks.py}. The symbol implementation is
  \path{engine/coin_city_symbol_context_arm.py}, with task builder
  \path{eval/build_coin_city_symbol_context_tasks.py}.

  \item[Execution and parsing.]
  The manifest-locked Opus~4.8 wrapper is
  \path{eval/run_coin_city_stable_relationship_claude_n250.py}; shared
  append-only execution and parsing are in
  \path{eval/run_frozen_task_shard.py}. Model-specific runners and
  \path{eval/run_coin_city_wrong_context.py} preserve configurations
  and response records under the data root. The symbol runner is
  \path{eval/run_coin_city_symbol_context.py}; its contemporaneous comparison arms
  are under \path{responses/symbol_control_matched_20260813/}. The open-model runner
  is \path{eval/run_coin_city_symbol_context_open_model.py}; the module-specific
  Qwen and Llama responses, hosted symbol arms, launchers, and run manifests are under
  \path{local_results/symbol_context_model_comparison_20260824/} relative to the
  code root.

  \item[Primary analysis.]
  \path{analysis/analyze_coin_city_stable_relationship_multimodel.py} under the
  code root produces \path{analysis/multimodel_results.json} under the data root,
  the authority for the primary-arm metrics and paired contrasts.

  \item[Arbitrary-symbol analysis.]
  \path{analysis/analyze_coin_city_symbol_context.py} supplies the common estimator;
  \path{local_results/symbol_context_model_comparison_20260824/build_final_comparison.py}
  writes the paper authority
  \path{analysis/symbol_context_model_comparison_20260824.json}. Focused design
  checks are in \path{tests/test_coin_city_symbol_context.py};
  \path{paper/generate_coin_city_symbol_context_tables.py} generates both appendix
  tables from the final comparison.

  \item[Reference-selection analysis.]
  \path{analysis/analyze_coin_city_reference_selection.py} under the code root
  recomputes the partial correlations of
  Section~\ref{app:coin-city-reference-selection} from the frozen responses and
  writes \path{analysis/reference_selection_results.json} under the data root; it
  makes no model call. Checks are in
  \path{tests/test_coin_city_reference_selection.py}, including one that requires
  the partial correlation and the symbol-arm $\Delta\rho$ contrast to name the same
  systems. \path{paper/generate_coin_city_reference_selection_tables.py} writes
  both appendix tables.

  \item[Probe-target ceilings.]
  \path{analysis/compute_bayes_ceilings.py} under the code root derives the
  Bayes-optimal ceiling for each probe target in
  Table~\ref{tab:coin-city-probe-ceilings} and the forecast-MAE floors drawn in
  Figure~\ref{fig:coin-city-results}, importing the generator constants from
  the engine module so they cannot drift, and writes
  \path{analysis/bayes_ceilings.json} under the data root. It reads no
  response file and makes no model call, and fails closed if its closed form and
  its simulation of the generator's own row logic disagree. Checks are in
  \path{tests/test_bayes_ceilings.py}.

  \item[Misleading arm and paper outputs.]
  The repository-level script
  \path{paper/generate_coin_city_split_figures.py} recomputes misleading-arm
  metrics from the final response files and scoring key.
  \path{paper/generate_coin_city_appendix_tables.py} generates the primary tables
  and pairwise contrasts.

  \item[Generator and deployment robustness.]
  \path{engine/coin_city_robustness_population.py} and
  \path{eval/build_coin_city_robustness_population.py} build the fresh-seed
  population; its design, responses, scheduler ledgers, and analysis are
  under \path{data/coin_city_population_075_040_sd010_v1/}.
  \path{analysis/analyze_coin_city_rationales.py} writes
  \path{analysis/rationale_audit_v1.json} under the primary data root.
  \path{eval/run_coin_city_gpt56_k0_repeat.py} and
  \path{analysis/analyze_coin_city_gpt56_k0_repeat.py} isolate and score the
  hosted repeat. \path{paper/generate_coin_city_robustness_tables.py} renders
  both appendix tables from the saved analysis artifacts.
\end{description}
%============================================================================
% END experiment2_appendix.tex
%============================================================================
%============================================================================
% BEGIN experiment3_appendix.tex
%============================================================================
% ============================================================
%  APPENDIX — EXPERIMENT 3 FULL DETAILS  (training & transfer)
%  Place after the Experiment 2 appendix.
% ============================================================

\input{tables/exp3_rebuilt_data}
\section{Experiment 3: controlled relationship transfer}
\label{app:exp3}

As the controlled transfer step, Experiment~3 asks whether outcome-based training
learns to select and apply an episode-matched relationship beyond its training
setting. The primary Qwen3-8B endpoint jointly changes the domain and response
mechanism; Qwen3-4B and Llama-3.1-8B provide scale and architecture comparisons.
The original three-seed design and results are retained below as an execution
record. Section~\ref{app:exp3-eight-seed} reports the later eight-seed Qwen
extension used for the main-text conclusions; the Llama and structureless
comparators retain their original seed counts.

\subsection{Design and evaluation}
\label{app:exp3-coin-structural}

\subsubsection{Data-generating processes}

Training uses only Structure~A, the direct relationship from Experiment~2. On
the displayed scale of domain $d$,
\begin{equation}
  y_{t+1}=b_d+g_jx_t+\epsilon_t,\qquad
  \epsilon_t\sim\mathcal N(0,\sigma_d^2).
  \label{eq:exp3-coin-a}
\end{equation}
A pulse acts in the next period and has no residual effect under zero subsequent
input. The unseen evaluation mechanism, Structure~B, is
\begin{align}
  m_{t+1}&=\rho m_t+g_jx_t,\\
  y_{t+1}-b_d&=\phi(y_t-b_d)+\lambda m_{t+1}+\epsilon_t.
  \label{eq:exp3-coin-b}
\end{align}
This is a linear-Gaussian latent-state process in the sense of
\citet{kalman1960new}: it introduces a persistent mediator and outcome carry-over,
so a single pulse has both immediate and delayed effects. Model prompts disclose neither equation,
parameter values, internal variable, nor structure label.

Coin City uses $b=50$, range $[0,100]$, driver pool
$\{-12,-8,-4,0,4,8,12\}$, and $\sigma=1.6$. Held-out Coin Harbor uses
$b=180$, range $[40,320]$, driver pool
$\{-24,-16,-8,0,8,16,24\}$, and $\sigma=4.5$, together with new terminology,
units, and qualitative descriptions. Evaluation gains follow
$g\sim\mathrm{Uniform}(.22,.92)$. Structure~B independently samples
$\rho\sim\mathrm{Uniform}(.48,.74)$,
$\phi\sim\mathrm{Uniform}(.18,.42)$, and
$\lambda\sim\mathrm{Uniform}(.68,.94)$. A reference and its matched target share
these parameters but have independent inputs and noise.

\begin{center}
\begin{minipage}{\linewidth}
\centering
\captionof{table}{\textbf{Exact Coin City--Coin Harbor domain contrast.} Both
surfaces instantiate the same Structure~A and B simulators and task template; Coin
Harbor changes the semantic setting, time unit, numeric scale, interventions, noise,
and pathway descriptions.}
\label{tab:exp3-domain-contrast}
\footnotesize
\setlength{\tabcolsep}{3pt}
\renewcommand{\arraystretch}{1.05}
\begin{tabular}{p{.20\linewidth}p{.36\linewidth}p{.36\linewidth}}
\toprule
\textbf{Field} & \textbf{Coin City (train and test)} & \textbf{Coin Harbor (test only)} \\
\midrule
Driver $\to$ outcome & Net campaign news $\to$ candidate support & Net shipping
bulletin $\to$ cargo-flow index \\
Time unit & Week & Tide cycle \\
Resting level / support & 50 / $[0,100]$ & 180 / $[40,320]$ \\
Observed-input pool & $\{-12,-8,-4,0,4,8,12\}$ &
$\{-24,-16,-8,0,8,16,24\}$ \\
Evaluation shocks & $\{-12,-6,0,6,12\}$ & $\{-24,-12,0,12,24\}$ \\
Observation-noise SD & 1.6 support points & 4.5 cargo-index points \\
Structure A description & Synchronized citywide alerts; reactions occur mostly in
the same week & Synchronized fleetwide radio; responses occur mostly in the same
tide cycle \\
Structure B description & Neighborhood conversations and local organizations;
reactions build gradually and persist & Dockside crews and harbor groups; responses
build gradually and persist \\
\bottomrule
\end{tabular}
\end{minipage}
\end{center}

\subsubsection{Training episodes and matched control}

Each training prompt contains two 14-case Coin City reference systems, both
governed by Equation~\ref{eq:exp3-coin-a}. One draws
$g\in[.18,.46]$ and the other $g\in[.68,.98]$; their order is randomized.
The target shares exactly one reference gain, selected in balanced fashion, and
shows $k\in\{0,2,4,8\}$ independently generated cases. Its national/local
background description identifies the matching weak or strong response as in
Experiment~2. There are 1,200 rows at each $k$, or 4,800 per training run.

The episode-matched and population-prior datasets have byte-identical prompts.
Episode-matched reward uses the target's noise-free forecasts. For the control, targets are
deranged within each $k$ stratum with a fixed permutation: no row keeps its own
target, while the complete multiset of target vectors is exactly preserved.
Consequently the control can learn the output format and population distribution
but cannot improve reward by using that row's references or target history.
The Qwen3-4B structureless diagnostic additionally scrambles displayed drivers
and rewards a flat intervention response. It isolates task-format learning and is
excluded from the episode-matched-versus-prior estimand.

\input{tables/exp3_training_example}

\subsubsection{Evaluation interface and qualitative hint}

Every evaluation prompt contains a complete 14-period Structure~A reference, a
complete 14-period Structure~B reference, and a target governed by one of them.
For Coin City the two reference descriptions are:
\begin{quote}\small
Residents receive campaign developments through synchronized citywide alerts.
Reactions occur mostly during the same week as the news.

Campaign developments usually spread through neighborhood conversations and
local organizations. Reactions may build gradually and persist after the
original news has faded.
\end{quote}
Coin Harbor uses the matched fleetwide-radio versus dockside-network descriptions.
The target receives the correct description, no description, or the other
pathway's description. The three prompt variants reuse the exact reference and
target trajectories, scenario grid, and gold forecasts; the manifest records a
common numeric hash for every cue triplet.

Each task requests ten counterfactuals: shocks
$\{-2u_d,-u_d,0,u_d,2u_d\}$ at horizons one and three, with $u_d=6$ in Coin City
and $u_d=12$ in Coin Harbor. The shock occurs in the next period and all later
inputs are zero. The response vector consists of the eight nonzero-shock minus
zero-shock contrasts. The full endpoint contains
\[
2\ {\rm domains}\times2\ {\rm structures}\times3\ {\rm hints}\times
4\ k{\rm\ values}\times30\ {\rm episodes}=1,440
\]
paired prompts. Training and evaluation seed bands begin at 81,000,000 and
91,000,000, respectively.

\noindent\textbf{Exact task template and prompt variants.}
\label{app:exp3-coin-prompts}
The following is the task string emitted by the prompt generator before the
model-specific chat wrapper is applied. Bracketed uppercase fields below denote
row-specific substitutions; they are not shown to the model. Line wrapping is
typographic, while the blank-line-separated sections and all fixed wording are
part of the prompt.

\begin{quote}
\begingroup\ttfamily\scriptsize\raggedright
Forecasting task in [DOMAIN].

\smallskip
The resting level of [OUTCOME] is [BASELINE]. Values are bounded to
[LOWER, UPPER].

\smallskip
The reference systems are complete examples. The target system may share a
stable response pattern with one reference. Infer from the information shown;
no response equation or coefficient is provided.

\smallskip
REFERENCE SYSTEM 1

Background: [REFERENCE 1 CONTEXT]

[REFERENCE 1 TABLE: period, driver, and 14 observed outcomes]

\smallskip
REFERENCE SYSTEM 2

Background: [REFERENCE 2 CONTEXT]

[REFERENCE 2 TABLE: period, driver, and 14 observed outcomes]

\smallskip
TARGET SYSTEM

[TARGET BACKGROUND, WHEN PRESENT]

[TARGET TABLE WITH $k$ OBSERVATIONS, OR THE EXACT NO-OBSERVATION SENTENCE]

\smallskip
COUNTERFACTUAL FORECASTS

For each scenario, apply the stated [DRIVER] in the next [PERIOD] and then set it
to 0 in later periods. Forecast [OUTCOME] at the stated horizon.

\smallskip
A: shock=[VALUE]; horizon=1\quad $\ldots$\quad
E: shock=[VALUE]; horizon=1

F: shock=[VALUE]; horizon=3\quad $\ldots$\quad
J: shock=[VALUE]; horizon=3

\smallskip
Return only strict JSON with exactly ten finite numbers in scenario order:

\{"forecasts":[A,B,C,D,E,F,G,H,I,J]\}
\endgroup
\end{quote}

The renderer realizes the variants as follows. Training always includes the
correct strong/weak background and uses the same input string for episode-matched and
population-prior training; only the hidden reward target differs. At evaluation,
\texttt{correct} inserts the target pathway's description, \texttt{none} omits
the background line entirely, and \texttt{misleading} inserts the other
pathway's description. Every cue triplet otherwise has identical tables,
interventions, and targets. If $k=0$, the target block contains the exact
sentence \emph{No target-system observations are available yet.}; otherwise it
contains the same table format as the references with the first $k$ rows. The
Qwen3-4B \texttt{structureless} diagnostic retains this shell but swaps every
displayed $-12.0\leftrightarrow12.0$ and $-8.0\leftrightarrow8.0$ token through
a collision-safe substitution. Untrained bases receive the same held-out
prompts as trained models. Thus cue, reward-arm, structureless, and base
comparisons do not introduce an unreported instruction change.

\subsubsection{Estimands, roster, and validation criteria}

The original factorial crosses Qwen3-4B, Qwen3-8B, and Llama-3.1-8B, two reward
arms, and three training seeds, for 18 runs. Three Qwen3-4B structureless runs
and three untrained endpoints provide diagnostic comparisons.
The primary estimand is population-prior minus episode-matched response MAE on
Structure~B in Coin Harbor with the correct hint, averaging the four $k$ values
equally. Positive values favor episode-matched training. The same estimator defines
factorial contrasts for Structure~A/Coin Harbor (domain transfer),
Structure~B/Coin City (structural transfer), and the in-distribution cell.
Paired absent-minus-correct and misleading-minus-correct contrasts measure hint
use. Their paired $k=8$ minus $k=0$ difference tests whether observations override
a misleading prior; negative values indicate a shrinking misleading-hint penalty.

Every reported round-300 endpoint uses five draws per prompt at temperature $.7$;
temperature-zero checkpoints serve only as online training monitors.

Each episode's score is the mean of its five per-draw response MAEs, so ``averaging
draws'' averages errors and not forecasts. The hierarchical bootstrap
resamples the three training seeds at the top level and paired episodes within
each selected seed. Parse failures receive the full observable-range error
penalty rather than being discarded. That penalty is the clip width---280 points
in Coin Harbor, roughly fifty times a typical error---so it carries large
leverage: a single unparsed draw of 150 raises the Qwen3-8B base joint-cell MAE
from $5.26$ to $5.72$, which is most of that arm's reported gap to matched
training. Both trained arms parse completely in that cell, so the primary
estimand is unaffected; the base comparison is not.
The launch ledger contains all 24 endpoints and 172,800 stochastic rows;
172,732 parse (99.96\%), and every endpoint has at least 99.82\% coverage.

\noindent\textbf{Original three-seed uncertainty analysis.}
\label{app:exp3-coin-cluster-sensitivity}
The original hierarchical interval resamples three top-level clusters, and a nonparametric
cluster bootstrap at $G=3$ is known to understate between-cluster uncertainty
\citep{cameron2008bootstrap}. We therefore report a deliberately conservative
alternative alongside it: a Student-$t$ interval on the three seed-level means with
two degrees of freedom, which uses the same paired episode differences but charges
the full small-sample penalty for estimating a variance from three numbers. Under
that alternative none of the twelve transfer-cell intervals excludes
zero, the primary Qwen3-8B contrast becoming $.251$ $[-.259,.762]$ and the Qwen3-4B
comparator $.319$ $[-.120,.758]$. Nine of the twelve intervals that exclude zero
under the hierarchical bootstrap therefore do not survive the $t$ alternative, and we
do not treat any original three-seed interval as robust to the choice of
small-cluster procedure.

In that original roster, the sign pattern was consistent without estimating a
variance. Across
the four transfer cells and three training seeds there are twelve seed-level
contrasts per model. All twelve favor episode-matched training for Qwen3-8B and all
twelve for Qwen3-4B, so all 24 Qwen seed-by-cell contrasts point the same way;
these reuse training seeds across cells and are not independent replications. The corresponding count for
Llama-3.1-8B is three of twelve. This sign pattern motivated the later precision
extension; its interpretation is superseded by the eight-seed results in
Section~\ref{app:exp3-eight-seed}.

Before training, the dataset passed 17 fail-closed checks: exact row counts and factorial
balance; byte-identical episode-matched/population-prior prompts and held-out datasets; exact
within-$k$ target marginals with no fixed reward links; absence of Structure~B
from training prompts and metadata; cue-triplet numeric identity; train/evaluation
identifier disjointness; hidden-equation leak scans; strict gold round trips; and
oracle separation from the blind pathway mixture. Mean A--B response separation
is 2.65 points in Coin City and 4.80 in Coin Harbor. Across both tokenizers,
the longest prompt is 1,071 tokens against the configured 2,304-token limit.

We estimate population-prior-minus-episode-matched intervals in all four transfer
cells, paired intervals for both cue penalties at every $k$, and the paired
$k=8$ minus $k=0$ change in misleading-minus-correct error to evidence override.
The protocol archive records versioned input, analysis, and output hashes. Full
optimization and compute parameters are in
Section~\ref{app:exp3-complete-training-setup}.

\subsubsection{Original three-seed scale and architecture comparators}
\label{app:exp3-coin-qwen-results}

This subsection preserves the original estimates. Its Qwen contrasts and
cross-size synthesis precede the eight-seed extension and do not establish the
current transfer claim; see Section~\ref{app:exp3-eight-seed}.

\noindent\textbf{Qwen3-4B scale replication.}
\label{app:exp3-coin-qwen4b}
The Qwen3-4B comparison includes three seeds in each of the episode-matched,
population-prior, and structureless arms. Each is evaluated on the same 1,440
tasks at round 300 with five stochastic draws per task; the untrained base uses
the same endpoint. Parse rates are 100\% for every trained arm
and 99.97\% for the stochastic base.

Joint-cell MAE is $6.15$, $6.76$, $4.81$, and $4.49$ for base,
structureless, population-prior, and episode-matched training, respectively.
The prior-minus-episode-matched difference is $.319$ $[.134,.489]$;
episode-matched training has the lowest point estimate in all four cells, and
all four original intervals exclude zero. The original Qwen3-8B estimate is likewise positive,
$.251$ $[.058,.470]$, giving the same interval-excluding-zero result at both
tested Qwen sizes.

\noindent\textbf{Qwen3-family summary.}
In the hardest joint-transfer cell, Qwen3-8B response MAE is $5.72$ for the
untrained base, $4.71$ after population-prior training, and $4.46$ after
episode-matched training. The prior-minus-matched contrasts are $.319$
$[.134,.489]$, $.251$ $[.058,.470]$, and $.120$ $[-.027,.261]$ at 4B, 8B,
and 14B. An equal-weight synthesis across the Qwen3 family is $.230$
$[.128,.333]$, a 4.9\% MAE reduction. All three size-level estimates, eight of
nine size--seed estimates, and 11 of 12 size--cell estimates agree in direction.
The original estimates suggested an advantage that attenuated with scale;
the eight-seed extension does not sustain that cross-size transfer conclusion.
At three seeds and 8B, about one fifth of the improvement
over the untrained base comes from preserving which evidence belongs to which
outcome rather than learning the format or common answers.

Averaged over evidence depths, omitting the hint raises episode-matched response
MAE by $.151$ $[.079,.235]$, and a misleading hint raises it by $.313$
$[.228,.415]$. The misleading-hint penalty is smaller at $k=8$ than at $k=0$,
but the evidence-override contrast $-.116$ $[-.309,.100]$ includes zero.

\noindent\textbf{Llama-3.1-8B architecture comparator.}
In the joint-transfer cell the prior-minus-episode-matched training contrast is
positive, $.420$ $[-.115,.958]$, but its interval spans zero. The three seed-level
effects are $+.762$, $+.622$, and $-.125$, so this result does not establish a
robust cross-architecture training advantage.

The joint cell is also the most favorable of the four Llama cells, and
reporting it alone would misrepresent the factorial, so we give all four. The
contrasts are $-.351$ $[-.680,-.130]$ in the in-distribution cell, $-.167$
$[-.862,.306]$ under structure-only shift, $-.581$ $[-.868,-.243]$ under
domain-only shift, and $+.420$ $[-.115,.958]$ under joint shift. In the two
Structure~A cells the population-prior arm is therefore better than the
episode-matched arm by an interval that excludes zero, and only three of the twelve
Llama seed $\times$ cell contrasts favor episode-matched training. On this
architecture the factorial runs against the hypothesis in the cells
whose response mechanism was seen during training, and recovers a positive point
estimate only where both shifts apply. We report this as a failure of the
episode-matching effect to replicate on Llama-3.1-8B rather than as a
directionally-consistent-but-underpowered result. The later Qwen extension
also weakens the original cross-size transfer interpretation.

Llama nevertheless uses the qualitative
pathway cue at the stochastic endpoint:
absent-minus-correct MAE is $.219$ $[.123,.313]$ and
misleading-minus-correct is $.210$ $[.092,.320]$. Its misleading-hint penalty
shrinks from $k=0$ to $k=8$ by $-.621$ $[-.957,-.323]$. All three cue-related
intervals exclude zero. These contrasts indicate sensitivity to the supplied
mechanism cue and reduced misleading-cue effects as observations accumulate.

\noindent\textbf{Additional Llama-3.1-70B scale check.}
To distinguish the Llama-8B result from a general failure of the architecture,
we launched a Llama-3.1-70B extension using the same training
and held-out data. We report a fixed exploratory snapshot: seed 42 at round 100
of 300, temperature zero, one draw on each of the same 1,440 held-out tasks.
The intervals below bootstrap paired tasks conditional on this single training
seed; they do not measure between-seed uncertainty and are excluded from the
primary three-seed factorial estimates.

The capacity result is already pronounced on the response mechanism seen in
training. In the correct-hint Coin City/Structure~A cell, response MAE falls
from $4.114$ for the untrained 70B checkpoint to $.410$ under population-prior
training and $.450$ under episode-matched training at round 100. For comparison,
the completed round-300 Llama-8B temperature-zero endpoints average $4.159$ and
$4.328$ in the corresponding arms. Thus, after only one third of the planned
schedule, 70B has roughly one tenth the direct-response error of 8B; both 70B
arms parse every output and classify Structure~A correctly on every task in
this cell.

This snapshot does not yet establish the treatment-specific transfer effect.
Across the in-distribution, domain-only, structure-only, and joint-transfer
cells, the round-100 prior-minus-episode-matched contrasts are respectively
$-.040$ $[-.098,.020]$, $+.132$ $[-6.166,6.820]$, $-.068$
$[-.156,.018]$, and $+1.778$ $[-.143,5.406]$. Positive values favor
episode-matched training. The large joint-cell point estimate is driven by one
additional unparsed population-prior response: among tasks parsed by both arms,
the contrast is only $+.039$ $[-.209,.276]$, and Structure~B classification is
near zero in both arms. The in-distribution contrast also changes from $+.291$
$[.209,.371]$ at round 50 to $-.040$ at round 100. We therefore draw the
narrow conclusion from this snapshot that scale lets Llama execute the seen numerical
relationship far more accurately, while episode-conditioned structural
transfer remains unresolved. This additional snapshot is distinct from
the primary three-seed, five-draw round-300 comparison.

\begin{center}
\begin{minipage}{\linewidth}
\centering
\captionof{table}{\textbf{Stochastic Coin City appendix comparators in the
joint-transfer cell, original three-seed roster.} The updated eight-seed Qwen
results appear in Table~\ref{tab:exp3-eight-seed}.
Entries are response MAE under the correct hint after averaging five draws per
task; lower is better. Positive prior-minus-episode-matched contrasts favor
episode-matched training. Intervals use the paired seed-first
bootstrap with a fixed cell-specific bootstrap seed.}
\label{tab:exp3-coin-complete-roster}
\scriptsize
\setlength{\tabcolsep}{3pt}
\resizebox{\linewidth}{!}{\input{tables/exp3_coin_structural_primary}}
\end{minipage}
\end{center}


\subsubsection{Later eight-seed precision extension}
\label{app:exp3-eight-seed}

Qwen3-8B and Qwen3-4B were extended to seeds 42--49 in both reward arms, holding
fixed the 300-round endpoint, five-draw decoder, data, and estimands. The added
seed rosters and four estimators were recorded before their added endpoints
existed. The seed count was chosen after seeing the original three-seed results;
this is a later precision extension, not an eight-seed preregistration of the
original study. The Qwen3-4B extension followed the disclosed Qwen3-8B interim
look at seven seeds, whose joint-cell estimate was $.108$ $[-.035,.265]$ while
a crashed seed-48 run was repeated. Eight seeds remained the fixed endpoint.
The amendment and interim records are in
\path{exp3_training_transfer/five_seed_extension/}.

\begin{center}
\begin{minipage}{\linewidth}
\centering
\captionof{table}{\textbf{Eight-seed prior-minus-episode-matched response MAE.}
Positive values favor episode matching; brackets give hierarchical 95\% intervals.
All cells use the correct cue and average evidence depths equally.}
\label{tab:exp3-eight-seed}
\small
\setlength{\tabcolsep}{5pt}
\begin{tabular}{lcc}
\toprule
Cell & Qwen3-8B & Qwen3-4B \\
\midrule
In distribution & $\eEightIdEst{}$ $\eEightIdCi{}$ & $\eFourbIdEst{}$ $\eFourbIdCi{}$ \\
Domain shift & $\eEightDomainEst{}$ $\eEightDomainCi{}$ & $\eFourbDomainEst{}$ $\eFourbDomainCi{}$ \\
Structure shift & $\eEightStructureEst{}$ $\eEightStructureCi{}$ & $\eFourbStructureEst{}$ $\eFourbStructureCi{}$ \\
Joint shift & $\eEightJointEst{}$ $\eEightJointCi{}$ & $\eFourbJointEst{}$ $\eFourbJointCi{}$ \\
\bottomrule
\end{tabular}
\end{minipage}
\end{center}

The Qwen3-8B joint estimate falls from $\ePubEightJointEst{}$
$\ePubEightJointCi{}$ to $\eEightJointEst{}$ $\eEightJointCi{}$; all four
Qwen3-4B intervals now span zero. The sole cell supported by all four procedures
is Qwen3-8B in distribution: $\eEightIdEst{}$ $\eEightIdCi{}$, Student-$t$
interval $\eEightIdTci{}$, wild-cluster $p=\eEightIdWild{}$, and
$\eEightIdSign{}$ positive seeds (one-sided sign-test $p=\eEightIdSignP{}$).
The hierarchical bootstrap uses 5,000 repetitions; the Student-$t$ interval has
seven degrees of freedom; the wild cluster bootstrap uses 9,999 Webb-weight
repetitions; the sign test is exact. Seed-level contrasts and all four
procedures for every cell are retained in
\path{qwen3_8b_step300_stochastic_eight_seed.json} and
\path{qwen3_4b_step300_stochastic_eight_seed.json} under the Coin City reports
root. These artifacts supply \path{tables/exp3_rebuilt_data.tex}.

The wrapper deviation record separately documents changed checkpoint-saving
flags in the Qwen3-8B added-seed block and the crossover reruns of seeds 42--44.
Their mean joint-cell contrast is $.255$, close to the original $.251$, although
individual seed contrasts change. This supports robustness of the mean at three
seeds, not exact equivalence of the wrappers. Those diagnostic runs are excluded
from the endpoint roster. Qwen3-4B's added
runs restored the parent flags before producing endpoints. The original
Llama-8B and Qwen3-4B structureless comparisons remain three-seed analyses.

\subsubsection{Structure-selection experiment}
\label{app:exp3-structure-selection}

Protocol \texttt{coin\_city\_structure\_selection\_v1}. The parent Coin City
design cannot test whether cue-driven selection transfers on functional form,
because its training set contains only the direct structure. This dataset makes
structure selectable. Each episode displays one direct and one mediated
reference whose horizon-1 responses are matched exactly---the direct gain is
scaled by the mediated system's $\lambda$, since a mediated horizon-1 response
is $\lambda\,g\,s$ against a direct $g\,s$---so persistence at horizon~3 is the
only discriminating feature and no model can select on response magnitude
instead. Mediated parameters are resampled every episode, so a memorised
template does not solve the task.

Solvability was verified before any data was generated, the check the parent
design lacked: a one-feature classifier on the coefficient of the lagged
outcome, applied to the fourteen reference rows, separates the structures with
$d=4.47$ in Coin City and $3.10$ in Coin Harbor and classifies at $99.0\%$ and
$97.2\%$. Choosing correctly is worth $2.18$ response-MAE points.

Training is Coin City only, semantic labels only, correct cue only, balanced
2{,}400/2{,}400 across target structures. Held out entirely: Coin Harbor, and an
arbitrary-label variant in which each structure is denoted by a symbol whose
mapping is redrawn every episode, so it can be recovered only from the reference
trajectories. Twenty-eight fail-closed checks gate the dataset, including that
the dimension under test varies in training and that the arbitrary symbols never
appear in a training prompt. Evaluation is fully crossed over two domains, two
label kinds, two target structures, three cue conditions and four evidence
depths, 120 tasks per cell.

The primary measure is the implied persistence ratio, mean $|$predicted
horizon-3 response$|$ over mean $|$predicted horizon-1 response$|$, computed as a
ratio of means so it stays finite when individual horizon-1 predictions are near
zero. Its true value is $0$ for direct targets and approximately $.65$ for
mediated. The cue-following index is the difference in that ratio between
episodes whose cue names a mediated structure and those whose cue names a direct
one.

The untrained reference was evaluated on the same 2{,}880-task set at $99.91\%$
parse. Its index is \eSelCitySemBase{} and \eSelCityArbBase{} on Coin City under
semantic and arbitrary labels and \eSelHarborSemBase{} and \eSelHarborArbBase{}
on Coin Harbor, with implied persistence between $1.05$ and $1.21$ in every
condition: the untrained model predicts full persistence irrespective of cue or
truth.

Results are a three-seed pilot of the episode-matched arm, reported as a pilot.
No claim about the episode-matched minus population-prior contrast is made at
three seeds; that contrast requires the eight-seed treatment applied to the
parent roster. A mid-training interim look at the greedy monitors was taken at
the author's direction and is recorded in
\texttt{INTERIM\_LOOK\_DISCLOSURE.md}; it agreed with the endpoint.



%============================================================================
% BEGIN experiment3_diagnostics_appendix.tex
%============================================================================
\subsubsection{Component diagnosis with paired cue flips}
\label{app:exp3-components}

The preceding RL pilot motivates a fixed, exploratory decomposition, frozen
before these inference outputs. We sample 96 fresh latent episodes, 24 in each
of the two domain $\times$ two label cells. Each episode is presented with both
cue assignments under four interfaces, giving 768 prompts per checkpoint:
(i) identify the named reference without forecasting; (ii) forecast with exact
unit-shock response coefficients at horizons 1 and 3 supplied for both
references; (iii) forecast with the selected reference and its coefficients
supplied; and (iv) the original forecasting task. The coefficients are latent
simulator summaries, not the requested ten forecasts; this assistance tests
whether a supplied relationship is usable, not whether exact coefficients are
recoverable from fourteen noisy observations. We use $k=0$ target observations
so numerical evidence is identical across the original-task cue flips. This
diagnostic does not test integration with target observations at $k>0$.

The fixed roster comprises Qwen3-8B base, its matched and fixed-prior RL
checkpoints at seeds 42--44, and an untrained Qwen3-32B feasibility comparison.
All 6,144 greedy generations parse. Selection-only accuracy is $100\%$ for
every checkpoint in every domain/label cell. Table~\ref{tab:exp3-components}
reports the frozen scores, including every assistance condition. At 8B,
supplying patterns does not reduce average response error in any domain/label
cell for the base or either trained arm. This locates a failure beyond merely
identifying the reference named by the cue; it does not uniquely separate
numerical application from sensitivity to the assisted prompt format.

\begin{table}[htbp]
\centering\small
\input{tables/exp3_components}
\caption{\textbf{Reference naming succeeds while forecasting remains difficult.}
Name is selection-only accuracy; Original, Both, and Selected are response MAE
for the three forecast interfaces. Trained rows average all three seeds, not
selected endpoints. Each checkpoint contributes 24 paired episodes per cell.
Domain scales differ, so compare interfaces within a domain.}
\label{tab:exp3-components}
\end{table}

Forecast structure accuracy assigns each predicted eight-coordinate response
vector to the closer of the episode's direct and mediated simulator templates
in MAE; ties count as incorrect. The paired cue-change calibration is
\[
\alpha=\frac{1}{N}\sum_{i=1}^N
\frac{(\widehat r_{i,M}-\widehat r_{i,D})^\top(r_{i,M}-r_{i,D})}
     {\lVert r_{i,M}-r_{i,D}\rVert_2^2},
\]
where $r$ subtracts the zero-shock forecast separately within each horizon,
and $M,D$ index the mediated and direct cue assignments. Zero denotes
invariance to the cue; one denotes the simulator's change along this direction,
but need not imply a correct full response vector. This differs from the
preceding persistence-ratio index, whose selecting reference is approximately
$.65$. Matched RL's mean $\alpha$ is $.000$, $.017$, $.136$, and $.008$ for
City semantic, City arbitrary, Harbor semantic, and Harbor arbitrary,
respectively. The corresponding fixed-prior means are $-.008$, $.027$,
$.085$, and $.035$; these three-seed descriptions are not a new confirmatory
comparison between training arms.

\paragraph{Stronger-model feasibility and output semantics.}
The frozen scorer clips forecasts to the valid domain range before computing
responses. Inspection revealed that Qwen3-32B often returns deviations from rest
instead of absolute levels; clipping those values distorts the response vector,
especially in Harbor, whose lower bound is 40. We therefore report a separately
labelled \emph{post hoc} raw-output audit (Table~\ref{tab:exp3-baseline-audit}),
without replacing the frozen scores or counting offset-corrected answers as
correct forecasts. With the selected pattern supplied in City/semantic, raw
response MAE is $.384$ versus $4.794$ on the original interface. Complete
relative vectors are accurate to $.01$ in $52.1\%$ of cases, but $95.8\%$ omit
the resting level. This shows partial conditional numerical ability, not a
solution of the original task. The 8B assistance failures remain predominantly
response-pattern errors after a common-offset correction.

\begin{table}[htbp]
\centering\small
\input{tables/exp3_baseline_audit}
\caption{\textbf{Post hoc Qwen3-32B audit of incorrectly offset answers.}
The first three numeric columns are unclipped response MAE. Exact means all
eight relative response coordinates are within $.01$ of truth; Omit rest means
both zero-shock predictions are near zero. These last two columns concern the
selected-pattern interface only. They do not measure correct absolute forecasts.}
\label{tab:exp3-baseline-audit}
\end{table}

\subsubsection{Supervised learnability and its transfer boundary}
\label{app:exp3-learnability}

We train two fixed forecast-SFT recipes from the same Qwen3-8B base on the
original 4,800 balanced Coin City/semantic prompts, with simulator-correct
absolute forecasts as targets. Both use seed 42, one epoch, effective batch 16,
300 optimizer updates, completion-only token cross entropy, LoRA rank 32 and
alpha 64 on all seven projection modules, AdamW $(.9,.95)$, no weight decay,
constant learning rate, and gradient norm cap 1.0. Rates are $10^{-6}$ (the RL
rate) and $2\times10^{-5}$ (an optimization-feasibility condition). Neither
Harbor nor arbitrary labels appear in training. Final endpoints alone are
evaluated, on the unchanged 768-prompt component diagnostic.

This matches distinct prompts and prompt presentations, not computation or
optimizer updates: RL samples eight completions per prompt and uses about
2,400 updates. The higher-rate SFT recipe also changes learning rate. Both
supervised outcomes are reported in Table~\ref{tab:exp3-learnability}; the
seed-42 RL checkpoint is the primary diagnostic comparison. The other two
matched RL seeds yield City/semantic structure accuracies $52.1\%$ and
$50.0\%$, with response MAE $1.029$ and $1.053$.

\begin{table}[htbp]
\centering\small
\setlength{\tabcolsep}{4pt}
\input{tables/exp3_learnability}
\caption{\textbf{A supervised recipe learns in-domain structure selection but
fails to transfer.} All trained recipes here use seed 42. Structure accuracy
is derived from forecasts, not reference naming. All outputs parse. Intervals
resample the 24 paired episodes within each cell (2,000 draws); they do not
capture training-seed uncertainty. Both fixed SFT learning rates are retained.}
\label{tab:exp3-learnability}
\end{table}

At the higher rate, City/semantic structure accuracy is $47/48$ ($97.9\%$),
response MAE is $.658$, and forecast-level MAE is $.527$. Calibration
$\alpha=1.413$ indicates overshoot despite accurate choice of response form.
The lower-rate recipe reaches only $58.3\%$ structure accuracy. Transferred
from City/semantic training, higher-rate SFT does not recover structure
selection under arbitrary labels or in Harbor, and its Harbor response errors
increase substantially. Original-task SFT response scores are unchanged when
clipping is removed, and an additive-offset audit does not remove the transfer
failure. Reference naming remains perfect. Thus the successful recipe
establishes in-domain learnability, while a semantic cue-to-pattern
association remains a possible explanation. It does not establish transferable
induction from reference trajectories or an isolated effect of SFT versus RL.

\paragraph{The boundary is arbitrary labels, not surface domain.}
Transfer failure alone cannot say whether the untrained cells are unreachable
by this recipe or merely unreached from City/semantic training. We therefore
trained the identical recipe natively in each of the three cells held out
above---City/arbitrary, Harbor/semantic and Harbor/arbitrary---at three seeds
each, changing only the training cell. Each model is evaluated on its own
cell's original task. Results are in
Table~\ref{tab:exp3-native-cells}.

The two conditions separate cleanly. Harbor/semantic reaches $100.0\%$
structure accuracy at every seed, matching the $97.9\%$ obtained in
City/semantic, so the second surface domain is fully learnable by this recipe
and the earlier Harbor failure was a transfer failure rather than a ceiling.
Both arbitrary-label cells remain at chance under native supervision---$41.7\%$
in City and $53.5\%$ in Harbor, with seed ranges of $[37.5,43.8]$ and
$[52.1,54.2]$ that exclude the semantic result---at the same recipe, budget
and seed count. Reference naming stays at $100.0\%$ in all three cells, so the
failure is not one of identifying the named reference. The boundary this
recipe encounters is therefore the arbitrary label mapping specifically, and
it survives native training; surface domain shift is not a barrier to it.

\begin{table}[htbp]
\centering\small
\setlength{\tabcolsep}{4pt}
\input{tables/exp3_native_cells}
\caption{\textbf{Native training separates the label condition from the
surface domain.} The higher-rate SFT recipe trained in each cell and evaluated
on that cell's original task, averaged over seeds 42--44. Seed range is the
minimum and maximum structure accuracy across those seeds. Name is
selection-only reference identification. Harbor is learnable under semantic
labels; neither arbitrary-label cell exceeds chance.}
\label{tab:exp3-native-cells}
\end{table}

\paragraph{Adaptive decision disclosure.}
The initially specified numerical launch gate was not met. Before any SFT
training, a documented outcome-informed amendment used the stronger model's
conditional numerical ability and absolute-level errors to justify running the
two already prepared, unchanged recipes. This is an exploratory learnability
pilot, not a passed preregistered gate or a confirmatory method comparison.
There was no SFT endpoint selection. Separate reference-identification
supervision was unnecessary given the perfect selection-only results.

\subsubsection{Mechanism-family evidence use under coherent gain changes}
\label{app:exp3-evidence-use}

The mechanism-family accuracy comparison uses a fixed population-prior training
control. It is distinct from the parent Coin City shuffled-pair control and
does not preserve the matched arm's exact target-vector distribution. To test
whether the accuracy gain depends on displayed episode-specific relationships,
we freeze a fresh gain-intervention set before inference: 16 pairs in each of
12 mechanisms, comprising 128 pairs from eight seen mechanisms and 64 from
four held-out compositions or parameter-extrapolation mechanisms. Both
mechanism-disclosure conditions use the same episodes, with $k=9$ target
observations.

Within each pair the target gain changes from the 20th to the 80th percentile
of that world's native uniform support. Inputs, exact noise draws, calibration
trajectories, baseline, and mechanism are fixed; affected target observations,
terminal states, and future outcomes are regenerated coherently. The disclosed
description provides the same gain range on both sides, never the actual gain.
This tests a gain change, not a change of mechanism identity. Native supports
are retained, including the narrow, noisy extrapolation world.

The primary metric is response-change MAE: subtract the zero-shock forecast
within each horizon, take the high-minus-low change of the resulting eight
responses, and compare with the simulator's change. This removes purely
additive level tracking. Greedy inference evaluates Qwen3-4B-Instruct-2507
base and the final matched and shuffled-target adapters at all five training
seeds, 42--46. All pairs parse. The mean contrasts and bootstrap intervals in
Table~\ref{tab:exp3-gain-evidence} resample paired training seeds and
world-stratified episodes (2,000 draws). The comparator is the
distribution-matched shuffled-target arm of
Appendix~\ref{app:exp3-shuffled-control}, not the fixed population prior, so
the contrast isolates episode-to-target pairing at a preserved answer
distribution. Individual seed results are retained in the released analysis.

\begin{table}[htbp]
\centering\footnotesize
\setlength{\tabcolsep}{3pt}
\resizebox{\linewidth}{!}{\input{tables/exp3_gain_evidence}}
\caption{\textbf{Better gain tracking on seen mechanisms does not establish
transferred evidence use.} Entries are response-change MAE (lower is better);
trained entries average the five training seeds 42--46. No change is an
invariant predictor. Slope M and Slope S are matched and shuffled tracking
slopes (ideal 1, zero for no response to changed evidence). Positive
shuffled-minus-matched contrasts favor matched training. Both held-out
intervals include zero, and matched held-out errors exceed the no-change baseline.}
\label{tab:exp3-gain-evidence}
\end{table}

Matched training improves on the shuffled-target control for seen mechanisms
by $.262$ $[.212,.319]$ (disclosed) and $.373$ $[.318,.426]$ (undisclosed),
but held-out contrasts are $-.024$ $[-.079,.027]$ and $-.016$
$[-.061,.027]$, both spanning zero. The tracking slopes separate the same way:
on seen mechanisms matched training reaches $.323$ and $.461$ against $.084$
and $.107$ for the shuffle, whereas on held-out mechanisms matched tracking
falls to $.122$ and $.238$ and its errors exceed the $.826$ no-change baseline
in both conditions. At five seeds, therefore, the evidence-use benefit is
established on seen mechanisms and absent on held-out ones. The previously
reported ordinary forecast-accuracy advantage remains a narrower result than
transferred evidence use: these interventions neither show a complete absence
of sensitivity nor establish robust tracking of held-out episode-specific
gains.

\subsubsection{Distribution-matched shuffled-target control}
\label{app:exp3-shuffled-control}

The accuracy comparison above uses a fixed population-prior control, which does
not preserve the matched arm's exact target-vector distribution. We therefore
ran a separately frozen shuffled-target control. It holds the prompts, the
training budget and the exact joint target-vector multiset fixed within
compatible world/evidence-depth groups, changing only which compatible episode
receives which target vector, so the two arms differ in the episode-to-target
pairing and in nothing else. Qwen3-4B-Instruct-2507 trains for 300 rounds in
each of \eShufG{} seeds per disclosure condition. The matched arm is rerun on
the same A5000 hardware and actor-memory allocation, so both arms share
hardware; endpoints use a frozen fresh held-out set of 720 tasks per seed under
greedy and five-draw stochastic decoding. The pre-registered primary estimand
is the stochastic shuffled-minus-matched response MAE. All 40 endpoint/decode
cells completed and every response parsed.

\begin{table}[htbp]
\centering\footnotesize
\setlength{\tabcolsep}{3pt}
\input{tables/exp3_shuffled_control}
\caption{\textbf{Matched training retains a smaller advantage over a
distribution-matched shuffled-target control.} Entries are shuffled-minus-matched
response MAE at \eShufG{} training seeds; positive favors matched training.
The hierarchical interval resamples seeds and trajectories; Webb $p$ is a
two-sided wild cluster bootstrap. The stochastic inferred-mechanism cell is the
pre-registered primary endpoint and the only cell clearing every estimator.}
\label{tab:exp3-shuffled-control}
\end{table}

Matched training retains an advantage on the primary endpoint in both
disclosure conditions. With the mechanism inferred it is
\eShufUndiscStochEst{} (\eShufUndiscStochSign{} seeds positive, Student-$t$
\eShufUndiscStochTci{}, hierarchical \eShufUndiscStochCi{}, wild cluster
$p=\eShufUndiscStochWild{}$, exact sign test $p=\eShufUndiscStochSignP{}$),
the only cell that clears every estimator. With the mechanism described the
point estimate is larger, \eShufDiscStochEst{} \eShufDiscStochCi{}, but rests
on \eShufDiscStochSign{} seeds, one of which reverses sign, and does not clear
the cluster-robust test ($p=\eShufDiscStochWild{}$). Greedy decoding is
directionally consistent in both conditions---\eShufDiscGreedyEst{} described
and \eShufUndiscGreedyEst{} inferred---but neither cell clears the
cluster-robust test.

\begin{table}[htbp]
\centering\footnotesize
\setlength{\tabcolsep}{3pt}
\input{tables/exp3_shuffled_control_level}
\caption{\textbf{The matched advantage does not appear in level error.}
Shuffled-minus-matched level MAE, same runs and estimators as
Table~\ref{tab:exp3-shuffled-control}. The inferred-mechanism cells show no
matched benefit. This contrast was not pre-registered and is descriptive.}
\label{tab:exp3-shuffled-control-level}
\end{table}

The benefit is concentrated in response error, which is what the manipulation
predicts. Because donors are drawn from the same compatibility group, the
shuffle changes an episode's target \emph{level} almost not at all while
changing its response \emph{shape} almost entirely: averaged over all 4,800
episodes, the recipient-donor level difference is $.0086$--$.0093$ against a
shape difference of $.892$--$.908$, a ratio of $96$--$104$ across both
disclosure conditions and all five seeds, with a within-group level standard
deviation of $.0076$. Both arms therefore receive essentially correct level
information and only the matched arm receives the correct episode-specific
relationship, so a level-error advantage is not expected. Its absence when the
mechanism is inferred (\eShufUndiscStochLevelEst{} \eShufUndiscStochLevelCi{},
$p=\eShufUndiscStochLevelWild{}$ stochastic; \eShufUndiscGreedyLevelEst{}
\eShufUndiscGreedyLevelCi{} greedy) is a manipulation check rather than a
negative result, and it argues against the benefit being generic improvement,
which would raise both metrics.

We report this as supporting rather than decisive, for three reasons. The
level contrast was not pre-registered. Level MAE is computed on the raw
forecast vector, so it partly contains the response component and is roughly
twice its magnitude, diluting any response-driven gain. And the pattern is not
uniform: the described condition does show positive level contrasts
(\eShufDiscGreedyLevelEst{} greedy, \eShufDiscStochLevelEst{} stochastic). The
asymmetry above is a property of the frozen builder, not of any prediction;
\path{measure_level_shape_asymmetry.py} recomputes it from the training targets
alone and records it in \path{level_shape_asymmetry.json}. None of this speaks
to held-out mechanisms, where gain tracking remains unestablished
(App.~\ref{app:exp3-evidence-use}). These estimates use the same-hardware
comparator; a supplementary historical-hardware comparison reusing the earlier
A6000 matched seeds 45--46 is retained in the released results and supports no
claim here.

\paragraph{Control status and reproducibility.}
Only completed diagnostics enter the tables and claims here. Their artifacts,
manifests, raw
outputs, and the outcome-informed SFT amendment reside under
\path{exp3_training_transfer/overnight_diagnostics_20260921/}, in
\path{structure/}, \path{learnability/}, \path{evidence_use/}, and
\path{shuffled_control/}.
\path{paper/generate_exp3_diagnostic_tables.py} renders the four diagnostic
tables from completed summaries and records source hashes in
\path{paper/tables/exp3_diagnostics_sources.json};
\path{paper/generate_exp3_shuffled_control_tables.py} renders the two control
tables and records hashes, the design freeze and the analysis-code hash in
\path{paper/tables/exp3_shuffled_control_sources.json}. Both refuse to emit a
partial roster.
%============================================================================
% END experiment3_diagnostics_appendix.tex
%============================================================================

\section{Experiment 4: historical-market training and held-out evaluation}
\label{app:exp4}
\label{app:exp3b-protocol}

Experiment~4 asks whether base-relative gains extend from the controlled setting
to events and normalized question families unseen during training. It uses
proper-score training, a point-in-time cutoff, and family-disjoint held-out data.

\subsection{Sample construction and leakage controls}

\noindent\textbf{Source archive and sample.}
The source export\footnote{%
  \micon{huggingface.png}\;%
  \href{https://huggingface.co/datasets/anonymous-review/polymarket-full-market-data}%
  {Full Polymarket archive}. Link anonymized for review.%
} has 1,788,703 rows and SHA-256
\texttt{5697f767\ldots f638e7d} (the full digest is stored in the manifest).
The builder scans the entire file and fails on checksum drift. The calendar
splits contain 1,736 train, 512 development, and 1,024 test markets, with YES rates
.311, .275, and .301. Train contains 991 distinct events and 1,456 normalized
question templates; development contains 292 and 453; test contains 493 and 866.
The selected archive rows have no official series identifiers, so event and
template components supply the operative family groups.

\noindent\textbf{Eligibility and point-in-time censoring.}
We admit literal Yes/No markets resolved by an extreme final token price. A
boundary-aware domain rule selects political, governmental, macroeconomic, and
geopolitical questions while explicitly rejecting culture, awards, sports,
esports, and crypto categories. This rule is fixed before model training. One task
is built 30 days before the earliest archived terminal timestamp. At least four
distinct UTC dates must have genuine pre-cutoff YES candles, the latest candle must
be no more than seven days stale, and its price must remain in $[.02,.98]$. We keep
at most the latest 14 dates. Empty, carry-forward, and post-cutoff candles are
rejected. The prompt omits the outcome, final prices and volume, market status,
realized terminal time, and the tags used for domain selection.

The archive records timestamped price history but not revisions to static question
or resolution-rules text. We therefore treat those strings as time-invariant. This
is a residual point-in-time limitation: the audit verifies that no timestamped or
outcome field crosses the cutoff, but cannot determine whether an archived rule
changed between the forecasting cutoff and resolution.

\noindent\textbf{Family-disjoint temporal split.}
Before applying sample caps, every development row sharing an event, available
series, or normalized template with train is removed; the same rule removes test
rows matching train or development. Numeric values and month names are normalized
in template keys. The final intersection is zero for each pair of splits. Sampling
within a split is a stable hash of task identifier and protocol version, so row
order in the export cannot change membership.

\subsection{Training and held-out evaluation}

\noindent\textbf{Matched update budget.}
Training comprises 1,736 unique markets arranged in a deterministic, balanced
schedule of 4,800 presentations, matching Experiment~3's 300-round update budget.
Every market appears twice and 1,328 appear a third time. Schedule order and task
files are checksum-verified. With batch size 16, this gives exactly 300
rollout-and-update rounds (2,400 AdamW steps). The manifests and tables distinguish
unique markets from presentations. The longest Qwen3-4B wrapped prompt contains
1,012 tokens against the 1,792-token training limit. Data, schedule, reward,
optimizer, update count, seeds, and scoring rule are held constant across models;
each uses its native chat template.
Shared optimizer, batching, and runtime details appear in
Section~\ref{app:exp3-complete-training-setup}.

\begin{center}
\begin{minipage}{\linewidth}
\centering
\captionof{table}{Experiment~4 training configuration. Held-out data are absent
from training and evaluated once per model after all three training seeds.}
\label{tab:exp3b-hparams}
\scriptsize
\begin{tabular}{p{.34\linewidth}p{.59\linewidth}}
\toprule
\textbf{Quantity} & \textbf{Value} \\
\midrule
Model roster / adapter & Qwen3-4B-Instruct-2507, Qwen3-8B, and
Llama-3.1-8B-Instruct; LoRA rank 32, $\alpha=64$ \\
Algorithm & Dr.~GRPO; $\beta=0$; AdamW, lr $10^{-6}$ \\
Unique train / presentations & 1,736 / 4,800 balanced repetitions \\
Rounds / seeds & 300; $42,43,44$ \\
Rollout & 16 prompts, 8 samples each; temperature 1.3 \\
Output & one JSON probability; 128 new-token limit \\
Reward & $1-(\hat p-y)^2$ only \\
Chat formatting & Qwen3-4B uses its frozen ChatML wrapper; Qwen3-8B uses
its non-thinking template; Llama uses its native instruct template. \\
Development & 512 tasks; temperature-zero evaluation every 25 rounds \\
Held-out evaluation & 1,024 tasks; temperature zero; one evaluation per model \\
\bottomrule
\end{tabular}
\end{minipage}
\end{center}

\subsection{Held-out results}
\label{app:exp3b-results}
\noindent\textbf{Decision rule.}
The primary learned contrast is the mean loss over three trained seeds minus base
loss. Invalid JSON receives Brier loss one. Intervals use 5,000 bootstrap
resamples of connected event/template components.

All Qwen3-8B test forecasts parse. Its three-seed trained mean obtains Brier
.11642 versus .12867 for the base, a difference of $-.01225$ with 95\% interval
$[-.01948,-.00602]$. Every training seed improves on the base. The interval
excludes zero, supporting transfer to held-out events for Qwen3-8B.

\begin{center}
\begin{minipage}{.94\linewidth}
\centering
\captionof{table}{\textbf{Qwen3-8B on the 1,024-market held-out set.}
The mean-row interval resamples connected event/template families; lower is
better.}
\label{tab:exp4-endpoint-results}
\scriptsize
\setlength{\tabcolsep}{3.5pt}
\resizebox{\linewidth}{!}{\input{tables/exp3b_qwen3_8b_endpoint_results}}
\end{minipage}
\end{center}

\begin{center}
\begin{minipage}{.9\linewidth}
\centering
\captionof{table}{\textbf{Qwen3-4B and Llama held-out comparisons.}
The Qwen3-8B result is immediately above. $\Delta$ columns are trained
mean minus comparator, so negative values favor training. Parse is the minimum
coverage across base and trained endpoints; Llama's 12 unparsed base outputs
receive Brier loss one.}
\label{tab:exp3b-model-roster-results}
\scriptsize
\setlength{\tabcolsep}{3.5pt}
\resizebox{\linewidth}{!}{%
\input{tables/exp3b_model_roster_results}
}
\end{minipage}
\end{center}

Qwen3-4B replicates the base-relative improvement: its trained mean is $.11589$
versus $.13160$ for base, a difference of $-.01572$
$[-.02528,-.00667]$. Llama-8B improves from $.15203$ to $.11502$, a difference
of $-.03701$ $[-.05376,-.02258]$; all three trained Llama endpoints parse.

The latest pre-cutoff market price obtains Brier $.11369$, and logistic
recalibration of its logit fit only on the training split obtains $.11232$.
Neither comparator informs language-model training. The trained-minus-market
interval spans zero for every model. Qwen3-4B is also statistically
indistinguishable from the calibrated price, but Qwen3-8B and Llama-8B have
higher error by $+.00409$ $[+.00098,+.00785]$ and $+.00269$
$[+.00067,+.00539]$, respectively. Thus the complete roster shows
base-relative improvement rather than superiority to market-based forecasts.

\subsection{Scale and architecture comparison}
\label{app:exp4-scale}

The scale comparison uses all 318 eligible 2026-window tasks whose identifiers and
event, series, and normalized question-template families are absent from the
training, development, and 1,024-market evaluation sets. The holdout contains 227
event families and 288 normalized templates; selection does not use outcomes or
market probabilities.

Qwen3 checkpoints at 4B, 8B, 14B, and 32B use three seeds, a common reward, a
4,800-presentation schedule, and the same development rule. The prospective 32B
extension was registered before the shared holdout was opened; it retained the
14B scientific settings and changed only memory distribution across four A6000
GPUs. Each untrained base and trained adapter receives five non-thinking draws
per task at temperature $.7$, top-$p$ $.8$, and top-$k$ 20. All 25,440 Qwen
draws satisfy the strict JSON contract.

\begin{center}
\begin{minipage}{.94\linewidth}
\centering
\captionof{table}{\textbf{Exact values behind the main-text scale figure.}
Lower Brier is better. Trained means average seeds 42--44. Missing five-draw task
groups receive Brier loss one.}
\label{tab:exp4-qwen-scale-results}
\scriptsize
\setlength{\tabcolsep}{3.5pt}
\resizebox{\linewidth}{!}{\input{tables/exp4_qwen_scale_results}}
\end{minipage}
\end{center}

The trained means are $.13664$, $.12859$, $.12695$, and $.12652$ at 4B, 8B,
14B, and 32B; the corresponding bases are $.15224$, $.13686$, $.12812$, and
$.12659$. The 32B seed Briers are $.12653$, $.12651$, and $.12651$.
Trained-minus-base contracts from $-.01559$ at 4B to $-.00827$ at 8B,
$-.00117$ at 14B, and $-.00007$ $[-.00045,+.00029]$ at 32B. Thus the larger
untrained checkpoints approach the same absolute performance reached by the
trained endpoints under the shared recipe. The registered 14B-minus-8B trained
contrast remains $-.00164$ $[-.00332,-.00024]$.

On this holdout, contemporaneous Polymarket prices have Brier $.12655$.
Trained 14B minus market is $+.00040$ $[-.00028,+.00139]$; trained 32B minus
market is $-.00003$ $[-.00009,+.00002]$. Both intervals cross zero. The scale
profile therefore shows convergence to the contemporaneous market frontier,
not a resolved advantage over it. Because data volume, update count, and
optimizer settings are held fixed, the comparison does not identify separately
tuned optima for each model size.

Llama-3.1-8B uses the same holdout and five-draw endpoint. All 4,770 trained draws
parse; nine of the 1,590 base draws do not, causing seven base task groups to
receive Brier loss one. Trained seed Briers are $.12638$, $.12629$, and
$.12662$, for a mean of $.12643$ versus $.16417$ for base. Trained minus base is
$-.03774$ $[-.05860,-.01946]$; trained minus Polymarket is
$-.00012$ $[-.00046,+.00012]$. The latter interval crosses zero.

\begin{center}
\begin{minipage}{.62\linewidth}
\centering
\includegraphics[width=\linewidth]{figures/exp4_llama_scale_results.pdf}
\captionof{figure}{\textbf{Llama scale replication.}
At 3B and 8B, open/filled squares show base/trained means, small squares show
seeds, and the dashed line shows Polymarket. The parameter axis is logarithmic.
The 3B base is descriptive because its strict-JSON recovery was specified after
observing syntax failures; all other endpoints follow the common protocol.}
\label{fig:exp4-llama-scale-results}
\end{minipage}
\end{center}

Qwen3-1.7B and Llama-3.2-3B preserve the same training/development data, three
seeds, 4,800-presentation schedule, five stochastic draws, and holdout.
Qwen3-1.7B seed Briers are $.12413$, $.12610$, and
$.12415$ (mean $.12479$) versus $.22531$ for base; trained minus Polymarket is
$-.00175$ $[-.00644,+.00265]$. Llama-3.2-3B seed Briers are $.13197$, $.13666$,
and $.14149$ (mean $.13671$); trained minus Polymarket is $+.01016$
$[+.00121,+.02060]$. The pinned public BF16 Llama mirror is recorded because the gated upstream
repository was unavailable to the workspace account. Its unconstrained base
initially produced only 4/318 complete five-draw groups at the 128-token cap;
an output-cap-only rerun reached 21/318. A disclosed, holdout-informed syntax
recovery therefore repeated the full base endpoint with identical prompts and
sampling settings while constraining output to a JSON object containing one
numeric probability. All 1,590 recovered draws parse. Base Brier is $.21887$,
and trained minus base is $-.08216$ $[-.11062,-.05367]$. Because the syntax
constraint was introduced after observing the failures and differs from the
trained endpoints, this base-relative contrast is descriptive. Both trained
models also parse every draw.

\noindent\textbf{Hosted frontier references.}
The API panel evaluates the same 318-task holdout. It is categorical rather than a parameter-scale
curve: provider parameter counts are unavailable or not comparable, and the
APIs do not expose identical decoding controls. DeepSeek-V4-Pro uses the common
temperature $.7$/top-$p$ $.8$ stochastic request. GPT-5.6-Sol uses provider-native
low reasoning without user-set temperature or top-$p$; Claude Opus~4.8 likewise
uses provider-native sampling. Each endpoint receives five independent calls
per task and uses the same strict parser, five-draw mean, Brier scoring, and
family-clustered 5,000-repetition interval procedure as the local panel.

DeepSeek parses all 1,590 draws and has Brier $.12660$, differing from the crowd by
$+.00005$ $[-.00105,+.00102]$. Claude parses all draws and has Brier $.12927$,
a difference of $+.00272$ $[-.00049,+.00614]$. GPT receives all 1,590
responses but six are malformed; two task groups therefore receive Brier loss
one, and task-level coverage is 99.37\%.
Its Brier is $.13335$, a difference of $+.00680$
$[+.00016,+.01635]$.

Kimi uses a 1,024-token output budget. Thirteen transport failures are retried
without replacing the 1,577 successful responses; all 1,590 final draws parse.
Its Brier is $.12682$, differing from the crowd by $+.00027$
$[-.00046,+.00109]$. Hosted results provide performance context, not evidence
for a common scale law or an identical-decoder comparison.

\subsection{Qwen3-4B evidence-updating evaluation}
\label{app:exp3b-update}
This evaluation applies the Experiment~1 update prompt and all nine fixed packets
for each of its 100 markets (three pro-$H_1$, three
anti-$H_1$, and three orthogonal). The packet and initial-context SHA-256 digests
are \texttt{6741be0a\ldots8796d36} and
\texttt{c550cf05\ldots0c098}. Every condition receives the same canonical first
structured Qwen~2.5-7B forecast for each market, so this paired analysis isolates
conditional revision rather than web retrieval or initial event-model
construction. Search and thinking are disabled; decoding is greedy with a
1,500-token output cap. The longest prompt is 1,480 tokens against a 6,144-token
context limit. These contrasts characterize revision behavior and are separate
from the held-out forecasting evaluation.

The paired evaluation produces all 900 packet-conditioned revisions for the base
and each final adapter, with 100\% parse coverage and ICS equal to one. Base
EHC/HFC is $.909$ on 591 eligible directional updates. Seeds 42, 43, and 44 obtain
$.917$, $.922$, and $.918$ (592, 588, and 591 eligible updates), for a
three-seed mean of $.919$. The paired market-bootstrap trained-minus-base EHC
difference is $+.0101$ with 95\% interval $[+.0004,+.0207]$; HFC has the same
estimate and interval $[+.0009,+.0210]$. Training reduces mean absolute
directional revision by $.52$ percentage points $[-.70,-.36]$ and orthogonal
revision by $.22$ points $[-.38,-.06]$; the directional-minus-orthogonal gap
narrows by $.30$ points $[-.52,-.09]$. The sensitivity ratio changes from
$4.64\times$ to $4.83\times$ ($+.19$ $[-.07,+.48]$). These components show
improved directional consistency alongside smaller revisions overall; they do
not indicate an increase in absolute packet selectivity.
\begin{center}
\begin{minipage}{.82\linewidth}
\centering
\captionof{table}{\textbf{Paired evidence-updating contrasts.}
Movement rows are mean absolute probability revisions in percentage points;
intervals resample the 100 markets.}
\label{tab:exp3b-update-seed-results}
\scriptsize
\setlength{\tabcolsep}{4pt}
\resizebox{\linewidth}{!}{%
\input{tables/exp3b_update_seed_results}
}
\end{minipage}
\end{center}

%============================================================================
% BEGIN experiment3_training_setup_appendix.tex
%============================================================================
\section{Experiments 3--4: shared training configuration and runtime}
\label{app:exp3-complete-training-setup}

Tables~\ref{tab:exp3-oat-common}--\ref{tab:exp3-runtime} record the effective
OAT/Dr.~GRPO~\citep{liu2024oat} arguments for the reported Coin City and Polymarket studies, including framework defaults omitted from the command line.
Execution manifests hash the data, trainer, launcher, evaluator, prompt
renderer, and each cell's environment.

\noindent\textbf{Step accounting.}
A ``step'' in the figures and earlier tables is one rollout-and-update round rather than one
AdamW call.  Each full run consumes 16 prompts per round and samples eight responses
per prompt (128 trajectories).  With one learner GPU, per-device microbatch one, one
PPO epoch, and gradient accumulation 16, OAT performs eight AdamW parameter updates per
round. Thus 4,800 prompt presentations give 300 reported rounds and 2,400 AdamW
updates. The command line supplied a .03 warmup ratio, but OAT selected
Transformers' \texttt{constant} scheduler, which does not use the warmup argument.
The effective learning rate was therefore $10^{-6}$ from the first optimizer
update. We use ``round'' below to avoid counting the eight within-round parameter
steps as independent data exposures.

\begin{center}
\begin{minipage}{\linewidth}
\centering
\captionof{table}{\textbf{Shared Dr.~GRPO implementation parameters.} These
values apply to the Coin City and Polymarket studies; the study-specific
tables below record exceptions.}
\label{tab:exp3-oat-common}
\scriptsize
\setlength{\tabcolsep}{3pt}
\begin{tabular}{p{.29\linewidth}p{.66\linewidth}}
\toprule
\textbf{Quantity} & \textbf{Effective value} \\
\midrule
Trainable parameters & Frozen base model; LoRA rank 32, $\alpha=64$, dropout 0,
no bias; targets \texttt{q\_proj}, \texttt{k\_proj}, \texttt{v\_proj},
\texttt{o\_proj}, \texttt{down\_proj}, \texttt{up\_proj}, and
\texttt{gate\_proj}; no quantization. \\
Algorithm & OAT Dr.~GRPO (\texttt{critic\_type=drgrpo}); eight-response group
advantage $r_i-\bar r$ with no group-standard-deviation division; policy loss is
token-summed with the generation limit as a fixed normalizer (no response-length
normalization). \\
Policy objective & PPO ratio clipping 0.2; truncated importance-sampling cap 2.0;
one PPO epoch; KL coefficient $\beta=0$; no critic/value model; reward whitening,
advantage whitening, and REINFORCE mode all off. \\
Optimizer & PyTorch AdamW \citep{loshchilov2019decoupled} (substituted for
DeepSpeed fused/CPU Adam); learning rate
$10^{-6}$; $\beta_1=.9$, $\beta_2=.95$; weight decay 0; maximum gradient norm 1.0. \\
Schedule & Constant $10^{-6}$ learning rate from the first optimizer update; the
supplied .03 warmup ratio is ignored by Transformers' \texttt{constant} scheduler;
one prompt epoch; 4,800 prompt presentations $=$ 300 rollout-and-update rounds
$=$ 2,400 AdamW steps. \\
Batching & 16 prompts per round; 8 sampled responses per prompt; 128 trajectories;
training microbatch 1; gradient accumulation 16; one learner GPU. \\
Training sampling & Temperature 1.3; top-$p=1$; top-$k=-1$; one independently
sampled completion at a time, repeated eight times per prompt. \\
Truncation and stopping & Model EOS only (no external stop string); a completion
ending at the length limit receives reward zero.  Missing or invalid required JSON
also receives reward zero. \\
Precision and memory & bfloat16; FlashAttention~2; gradient checkpointing;
DeepSpeed ZeRO-2 \citep{rajbhandari2020zero}; reference-model offload; vLLM prefix
caching \citep{kwon2023vllm}; dropout disabled. \\
Seeds and shuffling & Training seeds 42, 43, and 44.  Each dataset is deterministically
shuffled once, then OAT shuffles prompt batches independently under the training seed. \\
Checkpointing and logging & Evaluation includes step zero and the final round;
intermediate evaluation at the setting-specific frequency; one final LoRA adapter retained; metrics
logged each round; Weights \& Biases disabled. \\
\bottomrule
\end{tabular}
\end{minipage}
\end{center}


\subsection{Coin City A-to-B factorial}
\label{app:exp3-coin-structural-hparams}

Table~\ref{tab:exp3-coin-structural-hparams} records the original three-seed
configuration for Experiment~3. The current Qwen3-8B and Qwen3-4B analyses extend
both reward arms to eight seeds (Section~\ref{app:exp3-eight-seed}), retaining
the five-draw stochastic round-300 endpoint. The structureless diagnostic and
Llama-8B architecture comparator retain their original three seeds. The
table supplements the shared optimizer details in Table~\ref{tab:exp3-oat-common}; where the two tables differ, this
setting-specific table governs. The execution manifest stores the environment;
the protocol archive hashes the preflight manifest, data files,
generator, prompt renderer, scorer, and trainer entry point.

\begin{center}
\begin{minipage}{\linewidth}
\centering
\captionof{table}{\textbf{Original three-seed configuration for Coin City
hidden-structure transfer.} The later Qwen extensions are documented in
Section~\ref{app:exp3-eight-seed}.}
\label{tab:exp3-coin-structural-hparams}
\scriptsize
\setlength{\tabcolsep}{3pt}
\begin{tabular}{p{.29\linewidth}p{.66\linewidth}}
\toprule
\textbf{Quantity} & \textbf{Effective value} \\
\midrule
Model roster & \texttt{Qwen/Qwen3-4B-Instruct-2507},
\texttt{Qwen/Qwen3-8B}, and \texttt{meta-llama/Llama-3.1-8B-Instruct}. \\
Chat formatting & Qwen3-4B uses the explicit ChatML wrapper; Qwen3-8B and
Llama-3.1-8B use their tokenizer templates, with Qwen thinking disabled. All
include an assistant-generation prefix. \\
Original grid & 3 models $\times$ 2 reward arms
(episode-matched reward, stored under the legacy artifact key \texttt{causal}, and \texttt{population\_prior}) $\times$ 3 seeds
($42,43,44$) $=18$ full training runs. \\
Diagnostics & Qwen3-4B \texttt{structureless} at the same three seeds, plus one
untrained endpoint for each model. These six diagnostics are excluded from the primary
episode-matched-versus-prior estimand. \\
Training data & Coin City and Structure~A only; 4,800 prompts, balanced across
$k\in\{0,2,4,8\}$ and weak/strong target match. Two reference trajectories of
length 14 and one target trajectory of maximum length 8. \\
Comparator target construction & Episode-matched reward uses the episode truth.
Population-prior reward targets are a fixed no-self-match derangement within
each $k$; its target-vector multiset is exactly identical to episode-matched training.
Prompt strings are byte-identical across these two arms. \\
Held-out data & $2$ domains $\times2$ target structures $\times3$ hint arms
$\times4$ evidence depths $\times30$ episode seeds $=1,440$ rows. The complete
held-out dataset is byte-identical across all training arms. \\
Mechanism ranges & Evaluation $g\in[.22,.92]$; Structure~B
$\rho\in[.48,.74]$, $\phi\in[.18,.42]$, and $\lambda\in[.68,.94]$.
Training reference gains are $[.18,.46]$ and $[.68,.98]$. \\
Domain ranges & Coin City: baseline 50, support $[0,100]$, noise 1.6, driver
pool $\{-12,-8,-4,0,4,8,12\}$. Coin Harbor: baseline 180, support $[40,320]$,
noise 4.5, driver pool $\{-24,-16,-8,0,8,16,24\}$. \\
Interventions & Five shocks at horizons one and three; Coin City
$\{-12,-6,0,6,12\}$, Coin Harbor $\{-24,-12,0,12,24\}$. Ten joint forecasts
produce eight nonzero-minus-zero response contrasts. \\
Reward & $.4\,r_{\rm level}+.6\,r_{\rm response}$. Each component averages
$\max(0,1-|e|/S)$; level tolerance is 10\% of the domain's observable range
(10 Coin City points; 28 Coin Harbor units), and response tolerance is half one
shock unit (3 and 6, respectively). Invalid or length-truncated output receives
zero online reward. \\
Output contract & Exactly one JSON object with sole key
\texttt{forecasts} and exactly ten finite JSON numbers. The same strict parser
is used online and offline; no prose, extra/duplicate key, Boolean, non-finite
value, or wrong-length list is accepted. \\
Structured decoding & Syntax-only XGrammar GBNF in training, online evaluation,
endpoint evaluation, and base evaluation. It fixes object shape and arity only;
numeric values and their relationship to targets remain unconstrained. \\
Optimization & Frozen base plus LoRA rank 32, $\alpha=64$; OAT Dr.~GRPO;
PyTorch AdamW learning rate $10^{-6}$ with an effective constant schedule and no
warmup; $\beta=0$; one PPO epoch; gradient norm 1; ZeRO-2 and reference offload. \\
Batch and update accounting & 16 prompts per round; eight rollouts per prompt;
128 trajectories; microbatch 1; gradient accumulation 16; 4,800 presentations
$=$ 300 rollout-and-update rounds $=$ 2,400 AdamW steps. \\
Training generation & Temperature 1.3, top-$p=1$, top-$k=-1$, model EOS only;
192 new-token limit; prompt limit 2,304; total model context 3,072; vLLM prefix
caching enabled. \\
Precision / memory & BF16, FlashAttention~2, gradient checkpointing, dropout
off, no quantization; vLLM GPU fraction .38; expandable-segments CUDA allocator. \\
Online evaluation & All 1,440 rows at temperature 0, one draw, top-$p=1$,
192-token limit, batch 120, at step zero and every 50 rounds including round 300. \\
Reported endpoint evaluation & Five draws per prompt at temperature .7 and
top-$p=.95$; trained decode seed is training seed $+20{,}260{,}818$ and the
base seed is 20,260,819. Temperature-zero checkpoints are retained only as
online monitoring artifacts and are not reported as endpoints. \\
Offline scoring & Level MAE, response MAE, response error normalized by the
blind 50/50 pathway mixture, response correlation, strict parse coverage, and
nearest behavioral pathway. A parse failure receives the full domain range as
both level and response error. \\
Variance and estimands & Five-draw outputs first average within episode.
Original hierarchical intervals resample three training seeds, then paired
episodes; the precision extension resamples eight seeds and adds the
small-cluster checks in Section~\ref{app:exp3-eight-seed}.
Primary: prior-minus-episode-matched response MAE for correct-hint Structure~B in Coin
Harbor, equal over $k$. Domain, structure, joint, hint, and hint-by-$k$ contrasts
form the factorial analysis. \\
Preflight validation & A 10-round run for every model and evaluated arm must
save an adapter, produce the full temperature-zero monitoring and five-draw
stochastic evaluation sets, achieve
at least .95 strict parse rate under each decoder, and exhibit nonconstant finite
reward. The 17-check data preflight must also pass before estimating the factorial. \\
Compute & Training and trained-checkpoint evaluation use two NVIDIA A6000 48-GB
GPUs (one actor, one learner), 8 CPU cores, and 56--100 GB allocated host memory.
Untrained base evaluation uses one A6000, 8 CPU cores, and 24 GB host memory.
All runs load models from the same pinned local cache. \\
\bottomrule
\end{tabular}
\end{minipage}
\end{center}

\begin{center}
\begin{minipage}{\linewidth}
\centering
\captionof{table}{\textbf{Runtime environment for the reported training runs.}
Package versions come from the job execution environment.}
\label{tab:exp3-runtime}
\scriptsize
\setlength{\tabcolsep}{3pt}
\begin{tabular}{p{.29\linewidth}p{.66\linewidth}}
\toprule
\textbf{Component} & \textbf{Version / configuration} \\
\midrule
Python / CUDA build & Python 3.10.20; PyTorch 2.6.0 built for CUDA 12.4. \\
Training stack & OAT 0.2.4 at Git revision
\texttt{869706620a35ec5304fa80a0\allowbreak{}c74f0c566cc2bc57}; DeepSpeed 0.16.8; PEFT 0.15.2. \\
Generation stack & Transformers 4.51.3; vLLM 0.8.4; FlashAttention 2.7.4.post1. \\
Data stack & Datasets 2.16.1; Accelerate 1.14.0. \\
Allocator / logging & \texttt{PYTORCH\_CUDA\_ALLOC\_CONF=expandable\_segments:True};
Weights \& Biases disabled; stdout/stderr and per-round metrics retained with each
run. \\
\bottomrule
\end{tabular}
\end{minipage}
\end{center}
%============================================================================
% END experiment3_training_setup_appendix.tex
%============================================================================

\section{Experiments 3--4: reproducibility and artifact map}
\label{app:exp3-repro}

The artifact roots \path{exp3_training_transfer/coin_city_structural/}
and \path{exp3_training_transfer/polymarket/scripts/} contain the generators, prompt builders,
trainers, evaluators, and analysis code for the Coin City and historical-market
studies. Execution ledgers record effective hyperparameters, Git revisions,
adapter identities, and SHA-256 hashes of data and executable protocol files;
result renderers verify these fields before aggregation. For Coin City,
\path{make_qwen3_8b_endpoint_outputs.py} validates the seven original five-draw
files (50,400 scored draws). The later eight-seed reports in
Section~\ref{app:exp3-eight-seed} supply Figure~\ref{fig:exp3-transfer-results};
\path{render_exp3_paper_macros.py} renders the current prose values from those
reports and the recorded structure-selection and mechanism-family endpoints.
\path{make_paper_outputs.py} validates the original 24-endpoint roster and all
48 recorded endpoint files before writing the machine-readable result and the
stochastic appendix comparator table. Temperature-zero files are checked as
execution provenance but are excluded from manuscript estimates. For the
historical study,
\path{analyze_exp1_reapplication.py} validates four matched 900-record files and
computes the paired-market direction, sensitivity, and movement-component
contrasts. \path{generate_exp4_scale_figure.py} verifies the frozen 32B protocol,
summary, and raw-output hashes before adding the
point to Figure~\ref{fig:exp4-scale-results} and
Table~\ref{tab:exp4-qwen-scale-results}; it also writes the separate Llama
replication in Figure~\ref{fig:exp4-llama-scale-results}. The Polymarket
manifests preserve sample
membership, prompt and adapter hashes, parse coverage, and bootstrap settings.
%============================================================================
% END experiment3_appendix.tex
%============================================================================


\end{document}
